\documentclass{amsart}
\usepackage{amsmath,amssymb,amsthm,mathtools}
\usepackage{enumitem}
\usepackage{booktabs}
\usepackage{longtable}
\usepackage{float}
\usepackage[hidelinks]{hyperref}
\renewcommand{\UrlBreaks}{\do\_\do/\do.\do-}
\newcommand{\file}[1]{\nolinkurl{#1}}
\hypersetup{
  pdftitle={General line cones and the leading-form map},
  pdfauthor={Knowledge Sense},
  pdfsubject={Mathematics},
  pdfkeywords={line cone, negative cotangent cohomology, nonsmoothable curve singularity, Hilbert scheme of points, leading-form map, Horace method}
}

\newtheorem{theorem}{Theorem}[section]
\newtheorem{proposition}[theorem]{Proposition}
\newtheorem{lemma}[theorem]{Lemma}
\newtheorem{corollary}[theorem]{Corollary}
\theoremstyle{definition}
\newtheorem{definition}[theorem]{Definition}
\newtheorem{construction}[theorem]{Construction}
\theoremstyle{remark}
\newtheorem{remark}[theorem]{Remark}
\numberwithin{equation}{section}
\setlength{\emergencystretch}{2em}
\raggedbottom

\newcommand{\CC}{\mathbb{C}}
\newcommand{\PP}{\mathbb{P}}
\newcommand{\A}{\mathbb{A}}
\newcommand{\ZZ}{\mathbb{Z}}
\newcommand{\QQ}{\mathbb{Q}}
\DeclareMathOperator{\ev}{ev}
\DeclareMathOperator{\Gr}{Gr}
\DeclareMathOperator{\rank}{rank}
\DeclareMathOperator{\Hom}{Hom}
\DeclareMathOperator{\Spec}{Spec}
\DeclareMathOperator{\Fitt}{Fitt}

\begin{document}

\title{General line cones and the leading-form map}
\author{Knowledge Sense}
\address{Knowledge Sense}
\email{math@knowledgesense.jp}
\date{12 October 2026}
\subjclass[2020]{Primary 14B07; Secondary 14C05, 14H20, 13D10}
\keywords{line cone, negative cotangent cohomology, nonsmoothable curve singularity,
Hilbert scheme of points, leading-form map, Horace method}

\begin{abstract}
Let $Z$ be a general set of $r$ points in projective three-space and $C(Z)\subset\A^4$ the
affine cone over its standard homogeneous coordinate ring, a reduced curve consisting of
$r$ lines through the origin. Write $r=\binom{d+2}3+s$ with $0\le s<\binom{d+2}2$. We prove
that all cotangent groups $T^1_j(C(Z))$ of negative degree vanish for every fixed pair
with $d\ge84$ and $d+3\le s\le\binom{d+2}2-7d-427$. Equivalently in degree $-1$, the
differential of the map sending a configuration of affine points to the leading forms of
its degree-$d$ equations has the maximal rank $3r-4$ allowed by translations and dilation,
as predicted by Douvropoulos, Jelisiejew, N\o dland and Teitler and by Stevens.
Consequently these cones admit neither an algebraic nor an analytic smoothing, and the set
of branch numbers $r$ covered has natural density one. The lower end is sharp: vanishing
holds at $s=d+3$ for every $d\ge27$, the first position permitted by the dimension count,
whereas at $s=d+2$ the group $T^1_{-1}$ of a general cone has dimension $4(d+1)$ for every
$d\ge4$. The proof combines Fourier analysis on degenerate cyclic plane arrangements,
finite characteristic-zero certificates, Horace-type degenerations to plane point sets, and
Pinkham's equivariant versal deformation.
\end{abstract}

\maketitle

% =====================================================================
\section{Introduction}

A reduced curve singularity consisting of $r$ lines through the origin of $\A^n$ is the
affine cone over a set of $r$ points of $\PP^{n-1}$. Such cones are among the oldest
sources of non-smoothable curve singularities; the examples obtained by dimension counts,
going back to Pinkham \cite{Pinkham}, are surveyed and extended by Stevens
\cite[Sections~3.1--3.2]{Stevens}, and the subject has recently been revisited also by
Bozlee \cite{Bozlee}. A
second, more rigid mechanism is negative grading. The cotangent space $T^1$ of a cone is
graded, and if all its pieces of negative degree vanish, Pinkham's equivariant versal
deformation has only parameters of non-positive weight; the vertex then persists in every
deformation, and no smoothing exists. For finite schemes the analogous statement is part of
Jelisiejew's theory of trivial negative tangents \cite[Theorem~1.2]{JelElem}.

For lines in $\A^4$ this paper concerns the second mechanism. Work over $\CC$. For an ordered
set $Z$ of $r$ distinct points of $\PP^3$ let $C(Z)\subset\A^4$ be the affine cone defined by
the homogeneous ideal of the points, that is, the spectrum of the standard graded coordinate
ring of the given embedding. Put
\begin{equation}\label{eq:LM}
  L_d=\binom{d+2}3,\qquad M_d=\binom{d+2}2,\qquad r=L_d+s,\qquad m=M_d-s .
\end{equation}
Every $r\ge1$ is uniquely of this form with $0\le s<M_d$; for general $Z$ the ideal then
has no equations of degree less than $d$ and an $m$-dimensional space of equations of degree
$d$. By ``general'' we always mean: in a nonempty Zariski open subset of the space of ordered
configurations, for the fixed pair $(d,s)$ under consideration; we do not claim a common
open set for infinitely many pairs.

\begin{theorem}\label{thm:main}
Let $d\ge84$ and $d+3\le s\le M_d-7d-427$. For general $Z$,
\[
  T^1_j(C(Z))=0\qquad\text{for every }j<0,
\]
and the reduced curve $C(Z)$ admits neither an algebraic nor an analytic smoothing. The set
of branch numbers $r=L_d+s$ covered by these pairs has natural density one.
\end{theorem}

The lower end of the range is sharp, and it is reached in every degree from $d=27$ on.

\begin{theorem}\label{thm:lower}
For general $Z$, $T^1_j(C(Z))=0$ for every $j<0$, and $C(Z)$ admits neither an algebraic
nor an analytic smoothing, in each of the following cases:
\begin{enumerate}[label=\textup{(\alph*)}]
\item $d\ge27$ and $s=d+3$;
\item $d\ge27$ and $d+3\le s\le2d-24$;
\item $d\ge28$ and $d+3\le s\le3d-52$.
\end{enumerate}
On the other hand, for $d\ge4$ and $s=d+2$ the group $T^1_{-1}(C(Z))$ of a general cone has
dimension $4(d+1)$.
\end{theorem}

The last assertion only locates the onset of vanishing; it makes no statement about
smoothability at $s=d+2$.

The degree $-1$ part has a concrete meaning. Choose the chart $X_0\ne0$ and write the points
as $(1:p_i)$ with $P=(p_1,\dots,p_r)\in(\A^3)^r$. Let $F_j$ be the polynomials of degree at
most $j$ in three variables and $S_d$ the forms of degree $d$. For general $P$, evaluation is
injective on $F_{d-1}$ and surjective on $F_d$; the equations $K(P)=\ker(\ev_P|_{F_d})$ form
an $m$-dimensional space whose leading forms span
\[
  \Phi(P)=\{f_d:f\in K(P)\}\in\Gr(m,S_d).
\]
The ideal $\Phi(P)\oplus S_{\ge d+1}$ is the limit of the homothetic configurations $tP$ as
$t\to0$, so $\Phi$ is the homothety-limit map to the Grassmannian of very compressed algebras
studied in \cite[Section~5]{DJNT} and \cite[Section~2.3]{JelOpen}. Translations and the
dilation of $\A^3$ do not change $\Phi$, so its differential has rank at most $3r-4$. We
show (Proposition~\ref{prop:T1model}) that $T^1_{-1}(C(Z))$ is the kernel of this differential
modulo these four motions, and we prove the corresponding maximal-rank statement.

\begin{theorem}\label{thm:rank}
In all cases of Theorems~\ref{thm:main} and~\ref{thm:lower}\textup{(a)--(c)}, the
differential of $\Phi$ at a general configuration has rank $3r-4$; equivalently, the image of
$\Phi$ has dimension $3r-4$, and a general fibre has dimension four. For $d\ge4$ and
$0\le s\le d+2$ the differential at a general configuration is surjective.
\end{theorem}

The statement about the fibres is local: it does not assert that a general fibre is a single
orbit, nor any global injectivity. A dimension count (Section~\ref{sec:cone}) shows that the
rank $3r-4$ is possible only if
\[
  F_d(s)=s(M_d-s)-3(L_d+s)+4\ge0 ,
\]
and for $d\ge11$ this holds exactly for $d+3\le s\le M_d-d-6$; moreover
$F_d(d+2)=-4(d+1)$ and $F_d(d+3)=(d-11)(d+2)/2$. Thus Theorem~\ref{thm:lower}(a) is the
first dimensionally possible position in each degree, and Theorems~\ref{thm:main}
and~\ref{thm:lower} leave $6d+421$ positions open at the upper end of each layer $d\ge84$. In
the language of \cite{Stevens}, whose Proposition~3.5 gives the lower bound
$\dim T^1_{-1}\ge\max(0,-F_d(s))$ for general lines in $\A^4$, our results establish the
equality predicted in \cite[Conjecture~3.6]{Stevens} for all $s\le d+2$ ($d\ge4$) and
throughout the ranges above, where both sides vanish. They also establish, in these ranges,
the maximal rank of the homothety-limit map conjectured in \cite[Remark~5.4]{DJNT}, where it
was verified by computer for configurations of at most $95$ points. Existence results for
non-smoothable reduced curve singularities with prescribed genus and number of branches are
proved in \cite[Theorem~1.5]{Bozlee}; the present results concern the specific family of general line cones
in $\A^4$ and the vanishing of all their negative cotangent groups. To our knowledge, this
vanishing had previously been established for line cones in $\A^4$ only for finitely many
values of $r$, by direct computation \cite[Section~3.2]{Stevens}.

\subsection*{Outline of the proof}
The identification of $T^1_{-1}$ with the kernel of the moving differential
$J_P$ uses the description of $T^1$ of a curve through its normalization
(Section~\ref{sec:cone}). Vanishing in degree $-1$ propagates to all negative degrees by an
explicit obstruction map whose injectivity is witnessed by a lattice configuration
(Section~\ref{sec:lower}), and non-smoothability follows from Pinkham's construction and a
Fitting-ideal argument (Section~\ref{sec:nonsmooth}).

The rank statements are proved by exhibiting, for each pair $(d,s)$, one configuration of
maximal rank. For the lower boundary $s=d+3$ we use the $L_d=\binom n3$ triple points of
$n=d+2$ planes, $n$ further points, one on each plane, and the origin. Products of the plane
equations interpolate the triple points, and the velocities of the triple points can be
eliminated exactly (Section~\ref{sec:star}). We specialize to a cyclic arrangement, which is
not a genuine arrangement but on which the cyclic group of order $n$ acts; Fourier
decomposition reduces the remaining equations, for each character, to a system with six or
seven unknowns, whose kernel (the \emph{special kernel}) has total dimension $2n-2$
(Section~\ref{sec:cyclic}). Moving the auxiliary points off the planes by a homothety, we
show that the first-order perturbation has rank $2n-6$ on the special kernel; the four
remaining directions are the motions (Section~\ref{sec:opening}). For $n\ge89$ this is proved
uniformly by a three-column detector and exact estimates for three Fourier moments; for
$29\le n\le88$ it is verified by exact modular computations lifted to characteristic zero.
Further positions are reached by coupled Horace transitions from degree $e$ to $e+1$, which
add $b=M_e+c$ points on a plane (Section~\ref{sec:horace}); their two auxiliary rank
conditions are proved with canonical modules and syzygies of general plane points. The bands
of Theorem~\ref{thm:lower}(b),(c) are obtained from the lower boundary by sparse transitions
with $c=2,3$, and the broad range of Theorem~\ref{thm:main} by a finite propagation from an
explicit rational configuration in degree $11$, followed by an analytic induction
(Sections~\ref{sec:seed} and~\ref{sec:proof}).

\subsection*{Dependencies}
Besides standard facts, the proof uses the following results that we do not reprove:
Pinkham's construction of an equivariant formal versal deformation \cite{Pinkham}; the
cotangent sequence and the residue description of the canonical module of a curve of lines
as presented in \cite{Stevens}; Sauer's count of minimal generators for general plane points
\cite[Proposition~1.7]{Sauer}; and the Buchsbaum--Rim exactness criterion
\cite{Kustin}. It also uses finite exact computations: modular certificates for
$29\le n\le88$, a rational configuration in degree $11$, a finite list of transitions, and
a fixed finite set of explicit symbolic identities and nonvanishing certificates; their
scope is described in Appendices~\ref{app:finite}--\ref{app:data}. The statements for
infinitely many degrees are proved by uniform arguments and inductions which use these
fixed certificates as input (for instance as the base of an induction, or as identities of
rational functions in a parameter); they are never inferred from tests of finitely many
degrees.

\subsection*{Organization}
Section~\ref{sec:cone} sets up the cone ring and the moving differential.
Sections~\ref{sec:lower} and~\ref{sec:nonsmooth} treat lower negative degrees and
non-smoothability. Section~\ref{sec:star} introduces plane arrangements and proves the
statement at $s\le d+2$. Sections~\ref{sec:cyclic} and~\ref{sec:opening} prove the lower
boundary. Section~\ref{sec:horace} proves the Horace transitions, Section~\ref{sec:seed} the
seed and the propagation, and Section~\ref{sec:proof} completes the proofs.
Appendix~\ref{app:finite} describes the modular certificates, Appendix~\ref{app:boundary}
lists the leading coefficients and interval bounds for the boundary characters,
Appendices~\ref{app:seed} and~\ref{app:propagation} give the seed and the finite
propagation, and Appendix~\ref{app:data} describes the supplementary programs and data.

% =====================================================================
\section{The cone ring and the moving differential}\label{sec:cone}

\subsection{Conventions}
Throughout, $d\ge 2$ and $0\le s<M_d$ are integers, and $L_d$, $M_d$, $r$ and
$m$ are as in \eqref{eq:LM}; thus $m\ge1$, and the pair $(d,s)$ is the unique
representative of $r$. We write $F_j\subset\CC[x,y,z]$ for the space of
polynomials of degree at most $j$ (with $F_j=0$ for $j<0$) and $S_j$ for the space of
forms of degree $j$; thus $\dim F_{d-1}=L_d$, $\dim S_d=M_d$ and
$\dim F_d=L_d+M_d$. For $f\in F_j$ we write $f=\sum_{k\le j}f_k$ with $f_k\in S_k$,
and $D_bf$ denotes the derivative of $f$ in the direction $b\in\CC^3$.

Let $\operatorname{Conf}_r\subset(\PP^3)^r$ be the space of ordered $r$-tuples of
distinct points; it is irreducible. For $Z=(a_1,\dots,a_r)\in\operatorname{Conf}_r$
let $I_Z\subset S=\CC[X_0,\dots,X_3]$ be the homogeneous ideal of the reduced
set $\{a_1,\dots,a_r\}$, let $R=R(Z)=S/I_Z$ be its standard graded coordinate ring,
and let $C(Z)=\Spec R\subset\A^4$ be the affine cone; it is the union of the $r$ lines
through the origin corresponding to the points. We always use this ring, which is
the coordinate ring of the specified embedding; it differs in general from the full
section ring $\bigoplus_j H^0(\mathcal O_Z(j))$. The cotangent groups
$T^1_j=T^1(C(Z))_j$ are taken with respect to the grading of $R$, so that a
derivation of degree $j$ raises degrees by $j$.

Most computations take place in the affine chart $X_0\ne 0$. For
$P=(p_1,\dots,p_r)\in(\A^3)^r$ with distinct entries put $a_i=(1:p_i)$, and let
\[
  \ev_P\colon F_j\to\CC^r,\qquad f\mapsto (f(p_1),\dots,f(p_r)),\qquad
  V_j(P)=\ev_P(F_j).
\]
The \emph{evaluation open} $U^{\mathrm{ev}}_{d,s}\subset(\A^3)^r$ is the set of
configurations of distinct points for which $\ev_P|_{F_{d-1}}$ is injective and
$\ev_P|_{F_d}$ is surjective. On $U^{\mathrm{ev}}_{d,s}$ we put
\begin{equation}\label{eq:KQ}
  K(P)=\ker(\ev_P|_{F_d}),\qquad Q(P)=\CC^r/V_{d-1}(P),
\end{equation}
so that $\dim K(P)=m$ and $\dim Q(P)=s$. Every nonzero $f\in K(P)$ has degree
exactly $d$, since $K(P)\cap F_{d-1}=0$; hence the leading-form map
$f\mapsto f_d$ is injective on $K(P)$, and
\[
  \Phi(P)=\{f_d: f\in K(P)\}\in\Gr(m,S_d)
\]
is well defined. The \emph{moving differential} is the linear map
\begin{equation}\label{eq:JP}
  J_P\colon(\CC^3)^r\to\Hom(K(P),Q(P)),\qquad
  J_P(v)(f)=\bigl[(D_{v_i}f(p_i))_{i=1}^r\bigr].
\end{equation}
Its source has dimension $3r$ and its target dimension $sm$. The subspace
\begin{equation}\label{eq:WP}
  W_P=\{(b+a\,p_i)_{i=1}^r: b\in\CC^3,\ a\in\CC\}
\end{equation}
consists of the velocities of translations and of the dilation of $\A^3$.

\begin{lemma}\label{lem:newton}
Let $\Lambda\subset\ZZ_{\ge0}^3$ be a finite set that is closed under passing to
coordinatewise smaller multi-indices, and let $P_\Lambda=(\gamma)_{\gamma\in\Lambda}$,
regarded as a configuration of $|\Lambda|$ distinct points of $\A^3$. For
$\alpha\in\ZZ_{\ge0}^3$ put
\[
  N_\alpha(x_1,x_2,x_3)=\prod_{j=1}^3\prod_{i=0}^{\alpha_j-1}(x_j-i).
\]
Then $(N_\alpha)_{\alpha\in\Lambda}$ restricts to a basis of functions on
$P_\Lambda$, and $N_\alpha$ vanishes on $P_\Lambda$ whenever $\alpha\notin\Lambda$.
\end{lemma}

\begin{proof}
$N_\alpha(\gamma)=0$ unless $\alpha\le\gamma$ coordinatewise, and
$N_\gamma(\gamma)=\gamma_1!\gamma_2!\gamma_3!\ne 0$. Ordering $\Lambda$ by a linear
extension of the coordinatewise order makes the matrix
$(N_\alpha(\gamma))_{\alpha,\gamma\in\Lambda}$ triangular with nonzero diagonal. If
$\alpha\notin\Lambda$ and $\gamma\in\Lambda$, then $\alpha\le\gamma$ would give
$\alpha\in\Lambda$; hence $N_\alpha(\gamma)=0$.
\end{proof}

\begin{corollary}\label{cor:evopen}
For $d\ge2$ and $0\le s<M_d$ the evaluation open $U^{\mathrm{ev}}_{d,s}$ contains
configurations whose points affinely span $\A^3$; in particular it is nonempty.
\end{corollary}

\begin{proof}
Let $\Lambda$ consist of all multi-indices of degree at most $d-1$ and of $s$
multi-indices of degree $d$. By Lemma~\ref{lem:newton} the polynomials $N_\alpha$,
$\alpha\in\Lambda$, which lie in $F_d$, give a basis of $\CC^r$ after evaluation, and
those with $|\alpha|\le d-1$ form a basis of $F_{d-1}$ that evaluates to independent
vectors. Since $d\ge2$, $\Lambda$ contains $0,e_1,e_2,e_3$. Both conditions defining
$U^{\mathrm{ev}}_{d,s}$ are given by the nonvanishing of minors and are therefore open.
\end{proof}

\subsection{The graded pieces of the cone ring}
Let $P\in U^{\mathrm{ev}}_{d,s}$ with $r\ge4$, and suppose that
the points affinely span $\A^3$. Parametrize the $i$-th line of $C(Z)$ by
$t_i\mapsto t_ia_i$. Restriction to the lines embeds $R$ into the graded ring
\[
  K=\prod_{i=1}^r\CC[t_i,t_i^{-1}],\qquad K_j=\{(c_it_i^j)_i\}\cong\CC^r,
\]
which is the ring of homogeneous fractions of $R$. In degree $j\ge0$ the image of
$R_j$ is $V_j(P)$, since a form $G$ of degree $j$ restricts to
$(G(1,p_i)t_i^j)_i$. Hence $R_j=K_j$ for $j\ge d$, $\dim R_j=\dim F_j$ for
$j\le d-1$, and $I_Z$ has no nonzero elements of degree less than $d$. The
dehomogenizations $f(x)=G(1,x)$ of the forms $G\in(I_Z)_d$ are exactly the elements
of $K(P)$. Each line is smooth away from the vertex and two lines meet only at the
vertex, so $C(Z)$ is a reduced curve with an isolated singular point; since $d\ge2$
it has embedding dimension four. In particular all $T^1_j$ are finite dimensional
and vanish for all but finitely many $j$.

\begin{proposition}\label{prop:T1model}
Let $P$ be as above and $Z=((1:p_i))_i$.
\begin{enumerate}[label=\textup{(\alph*)}]
\item There is a natural isomorphism $T^1_{-1}\cong\ker J_P/W_P$, and $\dim W_P=4$.
\item Let $j\le-2$. Every class in $T^1_j$ has a unique representative given by a
vector $v=(v_i)\in(\CC^3)^r$, and the class is zero if and only if $v=0$. If
$j=-2$, the representative satisfies
\begin{equation}\label{eq:minus2}
  (D_{v_i}f(p_i))_{i=1}^r\in V_{d-2}(P)\qquad\text{for every } f\in K(P).
\end{equation}
\item For $j\le-2$, multiplication by $X_0$ maps $T^1_j\to T^1_{j+1}$ and
preserves the representing vectors $v$.
\end{enumerate}
\end{proposition}

\begin{proof}
We use the exact sequence obtained from $0\to R\to K\to K/R\to0$; see
\cite[Section~1.4, (1)]{Stevens}. Since all points lie in the chart $X_0\ne0$, the
vertex is the only point of $C(Z)$ with $X_0=0$, and $K=R[X_0^{-1}]$ is the coordinate
ring of the disjoint union of the $r$ punctured lines; it is smooth over $\CC$. Cotangent
cohomology is compatible with localization, so $T^1(R,K)=T^1(R[X_0^{-1}]/\CC,K)=0$, and
$T^1(R)$ is the cokernel of the map
$\operatorname{Der}(R,K)\to\operatorname{Der}(R,K/R)$. A derivation into $K/R$
is represented by a vector $\xi\in K^4$, determined modulo $R^4$, such that
$\xi\cdot\nabla G\in R$ for every $G\in I_Z$; it is enough to impose this for
a set of generators of $I_Z$. On the $i$-th line the Jacobian matrix of $I_Z$
has rank three over $\CC(t_i)$, with kernel spanned by the radial vector $a_i$,
so the image of $\operatorname{Der}(R,K)$ consists of the vectors whose $i$-th
component is a multiple of $a_i$. In degree $j$ the $i$-th component of $\xi$ is
$\xi_it_i^{j+1}$ with $\xi_i\in\CC^4$; subtracting the radial multiple
$(\xi_i)_0\,a_it_i^{j+1}$ we may assume that $\xi_i=(0,v_i)$, and this normalization
is unique.

Let $G\in(I_Z)_q$ and $f=G(1,x)$. For $k=1,2,3$ we have
$\partial_kG(t,tp)=t^{q-1}\partial_kf(p)$, so the $i$-th component of
$\xi\cdot\nabla G$ is $t_i^{q+j}D_{v_i}f(p_i)$. The condition is therefore
$(D_{v_i}f(p_i))_i\in V_{q+j}(P)$. For $j=-1$ and $q\ge d+1$ it is empty, because
$V_{q-1}(P)=\CC^r$; for $q=d$ it is the condition $J_P(v)(f)=0$. Thus the space of
normalized representatives in degree $-1$ is $\ker J_P$. The indeterminacy $R^4_{0}$
consists of the constant vectors $e=(e_0,e')\in\CC^4$; after normalization
they become $v_i=e'-e_0p_i$, which is $W_P$. Since the points affinely span $\A^3$, the
map $(b,a)\mapsto(b+ap_i)_i$ is injective. This proves~(a).

If $j\le-2$, then $R^4_{j+1}=0$, so there is no indeterminacy, and the
condition for $q=d$ is~\eqref{eq:minus2} when $j=-2$. Multiplication by $X_0$
multiplies the $i$-th component by $t_i$, which in the normalization keeps $v_i$
and raises the degree by one. This proves (b) and (c).
\end{proof}

\begin{lemma}\label{lem:Wker}
For every $P\in U^{\mathrm{ev}}_{d,s}$ we have $W_P\subset\ker J_P$. Consequently
$\rank J_P\le 3r-4$ when the points affinely span $\A^3$, and $T^1_{-1}=0$ if and
only if $\rank J_P=3r-4$.
\end{lemma}

\begin{proof}
Let $f\in K(P)$. For $v_i=b$ we get $D_bf(p_i)$, and $D_bf\in F_{d-1}$. For
$v_i=p_i$, Euler's formula gives $D_{p}f(p)=\sum_kkf_k(p)$, and since
$f(p_i)=0$ this equals $\bigl(\sum_k(k-d)f_k\bigr)(p_i)$, the value of a polynomial
in $F_{d-1}$. In both cases $J_P(v)(f)=0$. The rest follows from
Proposition~\ref{prop:T1model}(a).
\end{proof}

Comparing dimensions, $\dim T^1_{-1}\ge 3r-4-sm=-F_d(s)$, where
\begin{equation}\label{eq:F}
  F_d(s)=s(M_d-s)-3(L_d+s)+4 .
\end{equation}
Thus $T^1_{-1}=0$ is possible only if $F_d(s)\ge0$. A direct computation gives
\begin{equation}\label{eq:Fvalues}
  F_d(d+2)=-4(d+1),\qquad F_d(d+3)=\tfrac12(d-11)(d+2).
\end{equation}
The function $F_d$ is a concave quadratic in $s$ which is invariant under
$s\mapsto M_d-3-s$. Hence, for $d\ge11$, the integers $s$ with $F_d(s)\ge0$ form
exactly the interval
\begin{equation}\label{eq:window}
  d+3\le s\le M_d-d-6,
\end{equation}
whose endpoints are exchanged by this symmetry, as are $d+2$ and $M_d-d-5$.

\subsection{Openness}
\begin{lemma}\label{lem:open}
The set $U^{(-1)}_{d,s}$ of $P\in U^{\mathrm{ev}}_{d,s}$ whose points affinely span
$\A^3$ and satisfy $\rank J_P=3r-4$ is Zariski open in $(\A^3)^r$. Its image in
$\operatorname{Conf}_r$ under $P\mapsto((1:p_i))_i$ is open, and it consists of
configurations $Z$ with $T^1_{-1}(C(Z))=0$.
\end{lemma}

\begin{proof}
Cover $U^{\mathrm{ev}}_{d,s}$ by the open sets on which a fixed $r\times r$ minor of
$\ev_P|_{F_d}$ and a fixed $L_d\times L_d$ minor of $\ev_P|_{F_{d-1}}$ do not vanish.
On such a set Cramer's rule gives a basis of $K(P)$ and a complement-projection for
$V_{d-1}(P)$ whose entries are regular functions, so $J_P$ is represented by a
$sm\times 3r$ matrix of regular functions. The condition $\rank J_P\ge3r-4$ is the
nonvanishing of one of its minors of that size, and $\rank J_P\le3r-4$ holds
everywhere by Lemma~\ref{lem:Wker}. The map to $\operatorname{Conf}_r$ is an open
immersion onto the configurations missing the plane $X_0=0$, and the last assertion
is Lemma~\ref{lem:Wker}.
\end{proof}

Since $(\A^3)^r$ is irreducible, $U^{(-1)}_{d,s}$ is either empty or dense, and it is
nonempty as soon as one configuration in it is exhibited. All of our witnesses are
of this kind; some of them are only points of a larger irreducible parameter space
mapping to $(\A^3)^r$, and we explain the transfer in each case.

\subsection{The leading-form map}
The map $\Phi\colon U^{\mathrm{ev}}_{d,s}\to\Gr(m,S_d)$ is a morphism: on the charts in
the proof of Lemma~\ref{lem:open} it is given by the leading forms of a regular
basis of $K(P)$.

\begin{proposition}\label{prop:leading}
For $P\in U^{\mathrm{ev}}_{d,s}$ we have $\ker d\Phi_P=\ker J_P$. In particular,
in characteristic zero,
\[
  \rank J_P=3r-4\ \text{for general }P\iff
  \dim\overline{\Phi(U^{\mathrm{ev}}_{d,s})}=3r-4 .
\]
\end{proposition}

\begin{proof}
Let $v\in(\CC^3)^r$ and let $P_\varepsilon=(p_i+\varepsilon v_i)$ over
$\CC[\varepsilon]/(\varepsilon^2)$. A basis of $K(P_\varepsilon)$ has the form
$f_a+\varepsilon g_a$, where $(f_a)$ is a basis of $K(P)$ and $g_a\in F_d$ satisfies
$g_a(p_i)+D_{v_i}f_a(p_i)=0$ for all $i$. The tangent vector $d\Phi_P(v)$ is the
homomorphism $\Phi(P)\to S_d/\Phi(P)$ sending $(f_a)_d$ to the class of $(g_a)_d$.
If $d\Phi_P(v)=0$, then after subtracting a combination of the $f_b$ from $g_a$ we
may assume $(g_a)_d=0$, so $g_a\in F_{d-1}$ and $J_P(v)(f_a)=0$. Conversely, if
$J_P(v)=0$, then each $g_a$ may be chosen in $F_{d-1}$, and
$d\Phi_P(v)=0$. The second assertion follows from generic smoothness of $\Phi$
onto its image and from Lemma~\ref{lem:Wker}.
\end{proof}

\begin{remark}\label{rem:homothety}
The ideal $\Phi(P)\oplus S_{\ge d+1}\subset\CC[x,y,z]$ is the ideal of the flat
limit of the homothetic configurations $tP$ as $t\to0$: indeed
$t^df(t^{-1}x)\to f_d$ for $f\in K(P)$, and the Hilbert functions agree. Its
quotient has Hilbert function $(1,3,\dots,M_{d-1},s)$, of the very compressed type
discussed in \cite[Section~2.3]{JelOpen}. Proposition~\ref{prop:leading}
therefore identifies $T^1_{-1}=0$ with the statement that the map
$P\mapsto\lim_{t\to0}tP$ to the Grassmannian $\Gr(m,S_d)$ has image of the
maximal dimension $3r-4$ allowed by translations and dilation; this is the map and
the maximal-rank question studied in \cite[Section~5]{DJNT}. It means that a general
fibre has dimension four. It does not assert that a general fibre is a single orbit. In
the language of de Boor and Ron, $\Phi(P)\oplus S_{\ge d+1}$ is the space of leading forms
of the ideal of $P$, and \cite[Theorem~(4.8)]{deBoorRon} identifies it with the orthogonal
complement, for the apolarity pairing, of the least interpolation space of $P$, which is a
correct interpolation space by \cite[Theorem~(5.8)]{deBoorRon}; this duality says nothing
about the rank of the differential of $\Phi$.
\end{remark}

% =====================================================================
\section{Cotangent groups of lower negative degree}\label{sec:lower}

In this section we show that vanishing in degree $-1$ propagates to all negative
degrees. For $P\in U^{\mathrm{ev}}_{d,s}$ define
\begin{equation}\label{eq:Psi}
  \Psi_P\colon\CC^3\oplus\CC\to\Hom(K(P),S_{d-1}),\qquad
  \Psi_P(b,a)(f)=D_bf_d-a\,f_{d-1}.
\end{equation}
On the charts of Lemma~\ref{lem:open}, $\Psi_P$ is a matrix of regular functions, so
its injectivity is an open condition on $U^{\mathrm{ev}}_{d,s}$; it does not depend on
the choice of basis of $K(P)$.

\begin{proposition}\label{prop:Psi}
Let $d\ge2$ and $0\le s\le M_d-4$. Then $\Psi_P$ is injective for a nonempty open set
of $P\in U^{\mathrm{ev}}_{d,s}$.
\end{proposition}

\begin{proof}
Since $m=M_d-s\ge4$ and $d\ge2$, the four multi-indices
\[
  (d,0,0),\quad(0,d,0),\quad(0,0,d),\quad(d-1,1,0)
\]
are distinct, and we may choose a set $A$ of $m$ multi-indices of degree $d$ containing
them. Let $\Lambda$ consist of all multi-indices of degree at most $d-1$ together with
the multi-indices of degree $d$ not in $A$, and let $P=P_\Lambda$ be as in
Lemma~\ref{lem:newton}. Then $|\Lambda|=L_d+M_d-m=r$, the configuration lies in
$U^{\mathrm{ev}}_{d,s}$ by the argument of Corollary~\ref{cor:evopen}, and its points
affinely span $\A^3$. By Lemma~\ref{lem:newton} the polynomials $N_\alpha$,
$\alpha\in A$, vanish on $P$; their leading monomials $x^\alpha$ are distinct, so they
form a basis of $K(P)$.

The polynomials attached to the four chosen indices are
\begin{align*}
  N_{(d,0,0)}&=x^d-\tbinom d2x^{d-1}+\cdots, &
  N_{(0,d,0)}&=y^d-\tbinom d2y^{d-1}+\cdots,\\
  N_{(0,0,d)}&=z^d-\tbinom d2z^{d-1}+\cdots, &
  N_{(d-1,1,0)}&=x^{d-1}y-\tbinom{d-1}2x^{d-2}y+\cdots,
\end{align*}
where the dots denote terms of degree at most $d-2$; the factor $y$ of the last
polynomial contributes no term of lower degree. In $\Psi_P(b,a)$ we select the
coefficient of $x^{d-1}$, $y^{d-1}$, $z^{d-1}$ and $x^{d-2}y$ in the images of these
four polynomials, respectively. In the coordinates $(b_x,b_y,b_z,a)$ the resulting
$4\times4$ minor is
\[
  \det\begin{pmatrix}
    d&0&0&\binom d2\\
    0&d&0&\binom d2\\
    0&0&d&\binom d2\\
    d-1&0&0&\binom{d-1}2
  \end{pmatrix}
  =d^2\Bigl(d\tbinom{d-1}2-(d-1)\tbinom d2\Bigr)
  =-\frac{d^3(d-1)}2\ne0 .
\]
For $d=2$ the last row comes from the coefficient of $y$ in $N_{(1,1,0)}=xy$, and the
formula remains valid. Hence $\Psi_P$ is injective at $P_\Lambda$.
\end{proof}

\begin{proposition}\label{prop:lowerneg}
Let $d\ge2$, $0\le s\le M_d-4$ and $r\ge4$. Let $P\in U^{\mathrm{ev}}_{d,s}$ be a
configuration whose points affinely span $\A^3$ and for which $\Psi_P$ is injective,
and let $Z=((1:p_i))_i$. Then multiplication by $X_0$ induces injections
\[
  T^1_j(C(Z))\hookrightarrow T^1_{-2}(C(Z))\hookrightarrow T^1_{-1}(C(Z))
  \qquad(j\le-3).
\]
In particular, if $T^1_{-1}(C(Z))=0$, then $T^1_j(C(Z))=0$ for every $j<0$.
\end{proposition}

\begin{proof}
For $j\le-3$, multiplication by $X_0^{-2-j}$ preserves the representing vectors of
Proposition~\ref{prop:T1model}(b),(c), and a class in degree at most $-2$ is zero
exactly when its representing vector is zero; hence the first map is injective.

Let $v$ represent a class in $T^1_{-2}$ whose image in $T^1_{-1}$ vanishes. By
Proposition~\ref{prop:T1model}(a),(c) we have $v_i=b+ap_i$ for some $b\in\CC^3$ and
$a\in\CC$. Let $f\in K(P)$. By Euler's formula and $f(p_i)=0$,
\[
  D_{v_i}f(p_i)=H_f(p_i),\qquad
  H_f=D_bf-a\sum_{k=0}^{d-1}(d-k)f_k\in F_{d-1}.
\]
By \eqref{eq:minus2} there is $g\in F_{d-2}$ with $\ev_P(H_f)=\ev_P(g)$. Since
$\ev_P$ is injective on $F_{d-1}$, $H_f=g$, so the degree $d-1$ part of $H_f$
vanishes; this part is $D_bf_d-af_{d-1}=\Psi_P(b,a)(f)$. As $f$ was arbitrary,
$\Psi_P(b,a)=0$, hence $a=0$, $b=0$ and $v=0$.
\end{proof}

\begin{remark}\label{rem:lowerneg}
Proposition~\ref{prop:lowerneg} uses only necessary conditions coming from the
equations of degree $d$; equations of higher degree could only impose further
conditions. It gives no information on $T^1_{-1}$ itself. The Newton configuration
$P_\Lambda$ in the proof of Proposition~\ref{prop:Psi} is used only as a point of the
open set where $\Psi_P$ is injective; we make no claim about its group $T^1_{-1}$.
\end{remark}

% =====================================================================
\section{Nonsmoothability}\label{sec:nonsmooth}

The following statement is the form in which negative vanishing is used. For
the case of finite schemes compare \cite[Theorem~1.2]{JelElem}; the case of curves
is analogous, and we give the argument in the generality needed.

\begin{proposition}\label{prop:nonsmooth}
Let $C=\Spec S/I\subset\A^N$ be an affine cone with an isolated singular point at the
vertex, where $I$ is a homogeneous ideal generated in degrees at least $d\ge2$ and
$\dim C=1$. If $T^1_j(C)=0$ for all $j<0$, then the germ of $C$ at the vertex admits
neither an algebraic nor an analytic smoothing.
\end{proposition}

\begin{proof}
By Pinkham's construction \cite[Proposition~2.3 and (2.6)--(2.10)]{Pinkham}, the
germ has a formal versal deformation on which $\mathbb G_m$ acts compatibly with its
action on $C$. Concretely, choose a homogeneous basis of $T^1$; a basis element of
degree $e$ corresponds to a parameter $u$ of weight $-e$, and the defining equations
$G^0_1,\dots,G^0_k$ of $C$, of degrees $q_1,\dots,q_k$, lift to equations
$G_1,\dots,G_k\in\CC[[u]][X_1,\dots,X_N]$ that are homogeneous of degree $q_l$ for the
total weight. By hypothesis $e\ge0$ for every basis element, so every parameter has
weight at most $0$. A monomial $u^\alpha X^\beta$ in $G_l$ satisfies
$|\beta|+\operatorname{wt}(u^\alpha)=q_l$, hence $|\beta|\ge q_l\ge d$. Thus every
$G_l$ lies in $(X_1,\dots,X_N)^d$. In particular $X=0$ defines a section of the
versal family, and along this section all $G_l$ and all their first partial
derivatives in $X$ vanish, because $d\ge2$. The relative continuous cotangent module
restricted to the section is therefore free of rank $N$.

Suppose that a smoothing exists, algebraic or analytic. Restricting it to a curve in
its base through the special point which meets the locus of smooth fibres, and passing
to the normalization of that curve, we obtain a flat family $\mathcal X\to T$ over a
smooth curve germ $(T,0)$ with uniformizer $\tau$, with special fibre the germ of $C$
and with smooth fibres over $T\smallsetminus\{0\}$ near the singular point $x$. Let
$\Omega=\Omega_{\mathcal X/T}$. By the Jacobian criterion the relative singular locus of
$\mathcal X\to T$ is the zero set of $\Fitt_1(\Omega)$; it is contained in the fibre
$\tau=0$, so $\tau$ lies in the radical of $\Fitt_1(\Omega)\mathcal O_{\mathcal X,x}$
(by the Nullstellensatz in the algebraic case and by R\"uckert's Nullstellensatz in the
analytic case), and since the local ring is Noetherian, $\tau^n\in\Fitt_1(\Omega)
\mathcal O_{\mathcal X,x}$ for some $n$. Fitting ideals commute with base change and with
completion, so $\tau^n\in\Fitt_1(\widehat\Omega)$, where $\widehat\Omega$ is the
continuous relative cotangent module of the completed family over $\CC[[\tau]]$. By
formal versality the completed family is isomorphic, over $\CC[[\tau]]$, to the pullback
$\mathcal Y_\tau$ of the equivariant versal family along a morphism
$\CC[[u]]\to\CC[[\tau]]$; an isomorphism of formal families identifies the continuous
cotangent modules and their Fitting ideals, so $\tau^n\in\Fitt_1(\Omega_{\mathcal
Y_\tau})$. The section $X=0$ of the versal family pulls back to a section
$\iota\colon\CC[[\tau]][[X]]/(G)\to\CC[[\tau]]$ of $\mathcal Y_\tau$, and
$\iota^*\Omega_{\mathcal Y_\tau}$ is free of rank $N\ge2$ by the first paragraph.
Since Fitting ideals commute with base change, $\iota(\Fitt_1(\Omega_{\mathcal Y_\tau}))
=\Fitt_1(\iota^*\Omega_{\mathcal Y_\tau})=0$, whereas $\iota(\tau^n)=\tau^n\ne0$. This
contradiction proves the proposition.
\end{proof}

\begin{remark}\label{rem:formal}
The argument uses only the formal versal family and the formal isomorphism with the
pulled-back smoothing; no convergence of an equivariant coordinate change is
needed. It applies to every smoothing, equivariant or not. For the cones $C(Z)$ of
Section~\ref{sec:cone} the hypotheses on the singular point and on the degrees of
generators were verified before Proposition~\ref{prop:T1model}, with $N=4$.
\end{remark}

% =====================================================================
\section{Plane arrangements}\label{sec:star}

Throughout this section $n=d+2\ge 4$, so that $L_d=\binom n3$ and $M_d=\binom n2$.

\subsection{Star interpolation}
Let $\ell_1,\dots,\ell_n$ be affine linear forms $\ell_i(x)=1+\nu_i\cdot x$ with
$\nu_i\in\CC^3$.

\begin{definition}\label{def:genuine}
The arrangement of planes $\{\ell_i=0\}$ is \emph{genuine} if every three of the normals
$\nu_i$ are linearly independent and no four of the planes have a common point.
\end{definition}

For a $3$-subset $T\subset[n]$ let $a_T$ be the common point of the three planes
indexed by $T$; for a $2$-subset $\{a,b\}$ (an \emph{edge}) and a $3$-subset $T$ put
\[
  E_{ab}=\prod_{j\notin\{a,b\}}\ell_j\in F_d,\qquad
  V_T=\prod_{j\notin T}\ell_j\in F_{d-1},
\]
and write $V_{abk}=V_{\{a,b,k\}}$.

\begin{lemma}\label{lem:star}
Let the arrangement be genuine and let $A=(a_T)_T$. Then the $L_d$ points $a_T$ are
distinct, $(V_T)_T$ is a basis of $F_{d-1}$ with $V_T(a_U)=0$ for $U\ne T$ and
$V_T(a_T)\ne0$, and $(E_{ab})_{ab}$ is a basis of the space $I_A(\le d)$ of
polynomials in $F_d$ vanishing on $A$. In particular $\dim I_A(\le d)=M_d$, and the
leading forms of the $E_{ab}$ are linearly independent.
\end{lemma}

\begin{proof}
The point $a_T$ lies on exactly the planes indexed by $T$, so the points are
distinct, $V_T(a_T)\ne0$ and $V_T(a_U)=0$ for $U\ne T$. Since $\dim F_{d-1}=L_d$, the
$V_T$ form a basis and evaluation on $A$ is an isomorphism on $F_{d-1}$. Every
$E_{ab}$ vanishes on $A$. If $\sum c_{ab}E_{ab}=0$, restriction to a general point
of the line $\ell_a=\ell_b=0$ kills every term except $c_{ab}E_{ab}$, so $c_{ab}=0$.
Since $A$ imposes independent conditions on $F_{d-1}\subset F_d$, we have
$\dim I_A(\le d)=\dim F_d-L_d=M_d$. Finally $I_A(\le d)\cap F_{d-1}=0$.
\end{proof}

\subsection{The Jacobian in star coordinates}
Let $B=(b_j)_{j\in S}$ be $s$ further points, with $A\cup B$ distinct, and put
\[
  \mathcal E=(E_{ab}(b_j))_{j\in S,\ ab}\qquad(s\times M_d).
\]
If $\rank\mathcal E=s$, then $P=A\cup B$ lies in $U^{\mathrm{ev}}_{d,s}$ and
$K(P)=\{\sum_{ab}c_{ab}E_{ab}:\ c\in\ker\mathcal E\}$. At a triple point we use the
coordinates $t_{T,k}=\nu_k\cdot w_T$ ($k\in T$) of a velocity $w_T\in\CC^3$; since the
three normals in $T$ are independent, this is an invertible change of coordinates.

\begin{proposition}\label{prop:starjac}
Let $P=A\cup B$ be as above with $\rank\mathcal E=s$, and let $(w,v)$ be velocities at
$A$ and $B$. For an edge $ab$ and $j\in S$ put
\begin{equation}\label{eq:Xrow}
  X_{j,ab}(w,v)=D_{v_j}E_{ab}(b_j)-\sum_{k\notin\{a,b\}}t_{abk,k}\,V_{abk}(b_j).
\end{equation}
Then $(w,v)\in\ker J_P$ if and only if $\sum_{ab}c_{ab}X_{j,ab}(w,v)=0$ for all
$c\in\ker\mathcal E$ and all $j\in S$.
\end{proposition}

\begin{proof}
Let $f=\sum c_{ab}E_{ab}$ with $c\in\ker\mathcal E$. At $a_T$ only the edges contained in
$T$ contribute to the gradient, and for $T=\{a,b,k\}$ we get
$D_{w_T}E_{ab}(a_T)=V_T(a_T)\,t_{T,k}$. By Lemma~\ref{lem:star} the unique
$g\in F_{d-1}$ with $g(a_T)=D_{w_T}f(a_T)$ for all $T$ is
$g=\sum_{ab}c_{ab}\sum_{k\notin\{a,b\}}t_{abk,k}V_{abk}$. The vector
$(D_{v_i}f(p_i))_i$ lies in $V_{d-1}(P)$ if and only if it equals $\ev_P(g)$, that is,
if and only if $D_{v_j}f(b_j)=g(b_j)$ for all $j\in S$. More precisely, the map
$y\mapsto y_B-\ev_B(g_y)$, where $g_y\in F_{d-1}$ interpolates $y_A$, identifies
$Q(P)$ with $\CC^S$, and under this identification $J_P(w,v)(f)$ is the vector
$\bigl(\sum_{ab}c_{ab}X_{j,ab}(w,v)\bigr)_{j\in S}$.
\end{proof}

The condition of Proposition~\ref{prop:starjac} says that, for each $j$, the vector
$X_j=(X_{j,ab})_{ab}$ is a linear combination of the rows of $\mathcal E$. Introducing
the coefficients of these combinations as unknowns, we obtain the \emph{extended
system}
\begin{equation}\label{eq:extended}
  \Gamma_{j,ab}(w,v,\beta)=X_{j,ab}(w,v)-\sum_{l\in S}\beta_{jl}E_{ab}(b_l)=0
  \qquad(j\in S,\ ab),
\end{equation}
in the unknowns $t=(t_{T,k})$, $v\in(\CC^3)^S$ and $\beta\in\CC^{S\times S}$.

\begin{corollary}\label{cor:extended}
If $\rank\mathcal E=s$, the projection $(t,v,\beta)\mapsto(t,v)$ maps $\ker\Gamma$
isomorphically onto $\ker J_P$ (in star coordinates). Consequently
$\rank J_P=\rank\Gamma-s^2$.
\end{corollary}

\begin{proof}
For $(t,v)\in\ker J_P$ the coefficients $\beta_{jl}$ exist and are unique because the
rows of $\mathcal E$ are independent. The rank identity follows by comparing
$\dim\ker\Gamma$ with the number $3L_d+3s+s^2$ of unknowns.
\end{proof}

The four motions of Lemma~\ref{lem:Wker} have the following description: a
translation by $b$ has $t_{T,k}=\nu_k\cdot b$, $v_j=b$, $\beta=0$; the dilation has
$t_{T,k}=-1$, $v_j=b_j$ and $\beta=(n-2)I$. This follows from
$D_bE_{ab}=\sum_k(\nu_k\cdot b)V_{abk}$ and
$D_xE_{ab}(x)=(n-2)E_{ab}(x)-\sum_kV_{abk}(x)$, which are identities of polynomials.

\subsection{Formal normal coordinates}\label{sec:formal}
The entries of \eqref{eq:extended} are polynomials in $\nu=(\nu_1,\dots,\nu_n)$ and
in $B$, once the $t_{T,k}$ are regarded as independent unknowns. For an arbitrary
point $y=(\nu,B)$ of the affine space $Y_s=(\CC^3)^n\times(\CC^3)^S$, this defines the
\emph{formal} extended system $\Gamma^{\mathrm{f}}(y)$, a matrix with $sM_d$ rows and
$N_s=3L_d+3s+s^2$ columns. We stress that at a point $y$ whose arrangement is not
genuine, the triple points $a_T$ need not exist and $\Gamma^{\mathrm f}(y)$ is not
the differential of anything; it is merely a matrix of polynomials.

Let $Y_s^\circ\subset Y_s$ be the set of $y$ such that the arrangement is genuine,
$A\cup B$ consists of distinct points and $\rank\mathcal E=s$. It is open, and for
$s<M_d$ it is nonempty: the evaluations at points of $\A^3$ span the dual of
$I_A(\le d)$, and any prescribed point at which some $E_{ab}$ does not vanish can be
completed to $s$ points with independent evaluations.

\begin{lemma}\label{lem:transfer}
Let $Z\subset Y_s$ be an irreducible locally closed subset with $Z\cap Y_s^\circ\ne
\emptyset$.
\begin{enumerate}[label=\textup{(\alph*)}]
\item If $\rank\Gamma^{\mathrm f}(z_0)\ge N_s-4$ for some $z_0\in Z$, then for a
general point $z\in Z$ the configuration $P=A\cup B$ satisfies $\rank J_P=3r-4$, and
$U^{(-1)}_{d,s}$ is nonempty.
\item More generally, let $G$ be a matrix of rational functions on $Z$, regular on a
nonempty open subset $Z'$, and let $\rho,\rho'$ be integers such that, for every
$z\in Z'\cap Y_s^\circ$, $\rank G(z)\ge\rho$ implies $\rank J_P\ge\rho'$ for the
configuration $P=A\cup B$ of $z$. If $\rank G(z_0)\ge\rho$ for one $z_0\in Z'$, then
$\rank J_P\ge\rho'$ for a general configuration $P$ of $r$ points.
\end{enumerate}
\end{lemma}

\begin{proof}
In (b) the condition $\rank G\ge\rho$ is open on $Z'$ and nonempty, hence dense in the
irreducible $Z$; so it meets the nonempty open set $Z'\cap Y_s^\circ$, and at a common
point $\rank J_P\ge\rho'$. On the charts of Lemma~\ref{lem:open} the condition
$\rank J_P\ge\rho'$ is open on $(\A^3)^r$, which is irreducible. For (a) take
$G=\Gamma^{\mathrm f}$ and $\rho=N_s-4$: at a point of $Y_s^\circ$ the formal system
is the extended system, so by Corollary~\ref{cor:extended} and
Lemma~\ref{lem:Wker}, $4\le\dim\ker J_P=\dim\ker\Gamma^{\mathrm f}\le4$.
Lemma~\ref{lem:open} gives the conclusion for general configurations.
\end{proof}

Below, $Z$ will be a linear family through a non-genuine (cyclic) point and $Y_s$
itself, or the incidence variety of assigned points described next.

\subsection{Assigned points and the positions \texorpdfstring{$s\le d+2$}{s <= d+2}}\label{sec:outer}
For $S\subset[n]$ let $X_S\subset Y_s$ be the variety of $(\nu,B)$ with
$\ell_i(b_i)=0$ for $i\in S$; we then call $b_i$ the point \emph{assigned} to the plane
$i$. On the affine chart where every assigned point is nonzero, the fibre over $B$ is
a product of affine planes and free factors, so $X_S$ is irreducible. Put
$C_{ij}=\ell_j(b_i)$; on the open set where $C_{ij}\ne0$ for $i\ne j$, the point $b_i$
lies on no other plane, and $\alpha_i=\prod_{j\ne i}C_{ij}\ne0$. Then
$E_{ab}(b_j)=0$ unless $j\in\{a,b\}$, and $E_{ab}(b_a)=\alpha_a/C_{ab}$.

\begin{theorem}\label{thm:outer}
Let $d\ge4$ and $0\le s\le d+2$. For a general configuration $P$ of $r=L_d+s$ points
the map $J_P$ is surjective, $\rank J_P=s(M_d-s)$, and
\[
  \dim T^1_{-1}(C(Z))=3r-4-s(M_d-s)=-F_d(s).
\]
In particular $\dim T^1_{-1}(C(Z))=4(d+1)$ for general $Z$ at $s=d+2$.
\end{theorem}

\begin{proof}
We work on $X_S$ with $|S|=s$ and use only the columns of the triple coordinates
$t_{T,k}$; the auxiliary velocity columns are retained in $J_P$ but are not needed
for surjectivity. By the identification $Q(P)\cong\CC^S$ used in the proof of
Proposition~\ref{prop:starjac}, the target of $J_P$ is $\Hom(\ker\mathcal E,\CC^S)$,
whose dual is $\ker\mathcal E\otimes(\CC^S)^*$. A dual vector is thus a family
$(\lambda_i)_{i\in S}$ with $\lambda_i\in\ker\mathcal E$, and it annihilates the column
$t_{T,k}$ if and only if $\sum_{i\in S}\lambda_{i,ab}V_T(b_i)=0$ for $ab=T\smallsetminus
k$. Since $V_T(b_i)=0$ for $i\notin T$ and $V_T(b_i)=\alpha_i/(C_{ij}C_{ik})$ for
$T=\{i,j,k\}$, the substitution $\mu_{i,ab}=\alpha_i\lambda_{i,ab}$ turns these
conditions and the conditions $\lambda_i\in\ker\mathcal E$ into
\begin{align}
  \sum_{i\in S\cap T}\frac{\mu_{i,uv}}{C_{ij}C_{ik}}&=0
    &&(T\text{ a triple},\ uv\subset T,\ \{j,k\}=T\smallsetminus i),\label{eq:stress1}\\
  \sum_{k\ne j}\frac{\mu_{i,jk}}{C_{jk}}&=0&&(i,j\in S).\label{eq:stress2}
\end{align}
We call a solution a \emph{stress}. Let $G$ be the coefficient matrix of the linear
system \eqref{eq:stress1}--\eqref{eq:stress2} in the $sM_d$ unknowns $\mu$; its entries
are rational functions on $X_S$, regular where $C_{ij}\ne0$ for $i\ne j$, and the absence
of nonzero stresses means $\rank G=sM_d$. At points of $X_S\cap Y_s^\circ$ with
$\rank\mathcal E=s$, the absence of stresses is equivalent to the surjectivity of the
restriction of $J_P$ to the triple-coordinate columns, and therefore implies that $J_P$
itself is surjective, that is, $\rank J_P\ge sm$. By Lemma~\ref{lem:transfer}(b) with
$Z=X_S$, $\rho=sM_d$ and $\rho'=sm$, it suffices to find one point of $X_S$ with
$C_{ij}\ne0$ off the diagonal at which every stress vanishes. The set $X_S\cap Y_s^\circ$ is nonempty
by a general choice of a genuine arrangement and of points on the planes. The formula
for $\dim T^1_{-1}$ follows from Proposition~\ref{prop:T1model}, and the last
assertion from \eqref{eq:Fvalues}. The case $s=0$ is empty. We now produce stress-free
points.

\emph{Case $1\le s\le n-2$.} Let $a\ne b$ be indices outside $S$. Applying
\eqref{eq:stress1} to triples with exactly one index $i$ in $S$ shows that $\mu_i$
vanishes on every edge $\{i,a\}$ with $a\notin S$ and on every edge with both ends
outside $S$. Applying it next to triples $\{i,j,a\}$ with $i,j\in S$, $a\notin S$,
shows that every $\mu_i$ vanishes on each edge meeting the complement of $S$. For an
edge $\{i,j\}\subset S$ the two triples $\{i,j,a\}$ and $\{i,j,b\}$ give a $2\times2$
system for $\mu_{i,ij},\mu_{j,ij}$ that is nonsingular when
\begin{equation}\label{eq:outer4}
  C_{ja}/C_{ia}\ne C_{jb}/C_{ib};
\end{equation}
then both vanish, and the triple $\{i,j,l\}$ gives $\mu_{l,ij}=0$ for every other
$l\in S$. Condition \eqref{eq:outer4} holds simultaneously for all pairs for a general
choice of the two unused planes. Here \eqref{eq:stress2} was not used.

\emph{Case $s=n$.} Put $x_{ab}=\mu_{a,ab}/C_{ab}$ for $a\ne b$. Equation
\eqref{eq:stress1} for the triple $\{a,b,k\}$ and the edge $ab$ determines
$\mu_{k,ab}=-C_{ka}C_{kb}(x_{ab}/C_{ak}+x_{ba}/C_{bk})$. Equations \eqref{eq:stress2}
then become $\sum_{k\ne i}x_{ik}=0$ and
\[
  x_{ij}-\sum_{k\notin\{i,j\}}\Bigl(\frac{C_{ik}}{C_{jk}}x_{jk}
    +\frac{C_{ji}C_{ik}}{C_{jk}C_{ki}}x_{kj}\Bigr)=0\qquad(i\ne j).
\]
Suppose $C$ is symmetric, and put $Z_{ij}=x_{ij}/C_{ij}$, $Z_{ii}=0$ and
$c=Z^{\mathsf t}\mathbf 1$; no symmetry of $Z$ is assumed. Restoring the omitted terms
$k=i$ in the column sums, the equations become
\begin{equation}\label{eq:outer6}
  2(C\circ Z)=CZ^{\mathsf t}+C\operatorname{diag}(c),\qquad (C\circ Z)\mathbf1=0,\qquad
  Cc=0,
\end{equation}
where $\circ$ is the entrywise product; the diagonal of the first identity is the
row-sum condition, and the last follows by applying the first to $\mathbf1$.

Specialize to the circular point: $b_i=(\cos\theta_i,\sin\theta_i,0)$ and
$\ell_j(x)=1-\cos\theta_j\,x_1-\sin\theta_j\,x_2$ with $\theta_i=2\pi i/n$. Then
$\ell_i(b_i)=0$ and $C_{ij}=\tfrac12(2-\omega^{i-j}-\omega^{j-i})$ with
$\omega=e^{2\pi\sqrt{-1}/n}$; this matrix is symmetric with nonzero off-diagonal
entries. The equations \eqref{eq:stress1}, \eqref{eq:stress2} are invariant under a
common rescaling of $C$, so we may replace $C$ by $2-\omega^{i-j}-\omega^{j-i}$. All
normals lie in a plane, so this point of $X_S$ is not a genuine arrangement; it is
used only as a point of $X_S$. With $e_p(i)=\omega^{pi}$, the matrix $C$ has
eigenvalues $\lambda_0=2n$, $\lambda_{\pm1}=-n$ and $\lambda_p=0$ otherwise. Write
$Z=\sum z_{pq}e_pe_q^{\mathsf t}$, fix a total index $k$ modulo $n$, and put
$h_p=z_{p,k-p}$, $H=h_0$. Then \eqref{eq:outer6} becomes
\begin{equation}\label{eq:outer9}
  -2(h_{p-1}-2h_p+h_{p+1})=\lambda_p(h_{k-p}+H),\qquad \sum_ph_p=0,\qquad
  \lambda_kH=0 .
\end{equation}
For $p\notin\{0,\pm1\}$ the right side vanishes, so $h$ is affine on $1\le p\le n-1$:
with $A=h_1$, $B=h_{n-1}$ and $D=B-A$,
$h_p=A+(p-1)D/(n-2)$, and $\sum h_p=0$ gives $A+B=-2H/(n-1)$.

If $2\le k\le n-2$, subtracting the equations at $p=1$ and $p=-1$ gives
$4nD/(n-2)=0$, so $D=0$; the equation at $p=0$ then gives $2n(n-4)H=0$, so $H=0$ for
$n\ge5$, and $h=0$. If $k=\pm1$, then $H=0$ and $A+B=0$, and the equation at $p=0$
gives $h_k=0$, so $A=B=0$. If $k=0$, then $H=0$, $A+B=0$, and the equation at $p=1$
reduces to $n(n-4)D/(2(n-2))=0$, so $D=0$. Thus $Z=0$ for $n\ge5$, hence $x=0$ and
$\mu=0$. (For $n=4$ the stress space is two-dimensional. The case $s=n-1$ below needs
$n-1\ge5$, which is why we assume $d\ge4$.)

\emph{Case $s=n-1$.} Let $0$ be the unused index and $q=n-1\ge5$. On the assigned
indices take the circular point for $q$, and specialize $\ell_0$ to the constant
$1$, so $C_{i0}=1$. For assigned $i\ne j$ put $x_{ij}=\mu_{i,ij}/C_{ij}$; the triple
$\{0,i,j\}$ gives $x_{ij}=-x_{ji}$, so $Z_{ij}=x_{ij}/C_{ij}$ is skew-symmetric. With
$a_j=\mu_{j,0j}$ the same triple gives $\mu_{i,0j}=-a_j$ for $i\ne j$, and
\eqref{eq:stress2} becomes
\begin{equation}\label{eq:outer12}
  2(C\circ Z)+CZ+C\operatorname{diag}(Z\mathbf1)-\mathbf1a^{\mathsf t}=0,\qquad
  a_i=(CZ)_{ii};
\end{equation}
for instance the $(i,j)$ entry is $2C_{ij}Z_{ij}+(CZ)_{ij}+C_{ij}(Z\mathbf1)_j-a_j$,
which includes the term of the unused edge. In Fourier coordinates as before, skew
symmetry gives $h_p=-h_{k-p}$ and $\sum h_p=0$, and \eqref{eq:outer12} gives, away
from $p=0$, $-2(h_{p-1}-2h_p+h_{p+1})+\lambda_p(h_p-H)=0$. Hence $h$ is affine on
$1\le p\le q-1$, and the equations at $p=\pm1$ read
\[
  (q-2)(A-H)+\tfrac{2D}{q-2}=0,\qquad (q-2)(B-H)-\tfrac{2D}{q-2}=0 .
\]
Adding gives $A+B=2H$, and with $H+(q-1)(A+B)/2=0$ this gives $H=0$. Subtracting
gives $q(q-4)D/(q-2)=0$, so $D=0$ and all modes vanish. Then $Z=0$, $x=0$, $a=0$, and
\eqref{eq:stress1} gives $\mu=0$.

Finally, at both circular points $\rank\mathcal E=s$. Indeed, let $\sum_iy_iE_{ab}(b_i)=0$
for all edges $ab$. Since $E_{ab}(b_i)=0$ unless $i\in\{a,b\}$ and
$E_{ij}(b_i)=\alpha_i/C_{ij}$, the column of an assigned edge $\{i,j\}$ gives
$\alpha_iy_i/C_{ij}+\alpha_jy_j/C_{ji}=0$. Put $y'_i=\alpha_iy_i$; since $C$ is
symmetric, $y'_i+y'_j=0$ for all assigned $i\ne j$, and three assigned indices force
$y'=0$, hence $y=0$ because $\alpha_i\ne0$. The required open conditions are therefore
nonempty, and the theorem follows.
\end{proof}

\begin{remark}\label{rem:outer}
Theorem~\ref{thm:outer} does not include the first position $s=d+3$ at which
$F_d(s)\ge0$, and it implies nothing about that position; this is the content of
Sections~\ref{sec:cyclic} and~\ref{sec:opening}.
\end{remark}

% =====================================================================
\section{Cyclic arrangements and the special kernel}\label{sec:cyclic}

From now on until the end of Section~\ref{sec:opening}, $n=d+2$, $s=n+1=d+3$ and
$r=\binom n3+n+1$. The auxiliary points are $n$ points, one assigned to each plane,
together with the origin $o$, and they will be moved off the planes by a homothety.
The relevant configurations are not genuine, and all statements are about the
formal extended system of Section~\ref{sec:formal}.

\subsection{The cyclic family}\label{sec:cycfam}
Let $G=\mu_n\subset\CC^\times$ and fix a generator $\omega$. Let $F=\{0,1,3,7\}$ and
consider coefficient vectors $(a_1,a_3,a_7)$ in the affine plane
$\Sigma=\{a_1+a_3+a_7=-1\}$; put $a_0=1$ and
\[
  f(u)=\sum_{e\in F}a_eu^e,\qquad f(1)=0 .
\]
For $i\in\ZZ/n$ let $t_i=\omega^i$ and
\[
  p_i=(t_i,t_i^3,t_i^7),\qquad \nu_i=(a_1t_i^{-1},a_3t_i^{-3},a_7t_i^{-7}),\qquad
  \ell_i(x)=1+\nu_i\cdot x .
\]
Then $\ell_j(p_i)=f(t_i/t_j)$; in particular $\ell_i(p_i)=0$. For $\varepsilon\in\CC$ let
$y(\varepsilon)\in Y_{n+1}$ consist of the normals $\nu_i$ and the auxiliary points
$b_i(\varepsilon)=(1-\varepsilon)p_i$ ($i\in\ZZ/n$) and $b_o=o$, the origin; we write
$b_j=b_j(0)$, so $b_i=p_i$. At $\varepsilon=0$ the point $p_i$ is assigned to the plane
$i$. The linear map $\sigma(x)=(\omega x_1,\omega^3x_2,\omega^7x_3)$ satisfies
$\sigma(p_i)=p_{i+1}$ and $\ell_{i+1}\circ\sigma=\ell_i$, and it commutes with
homotheties. Hence $\ZZ/n$ acts on the unknowns and the equations of
$\Gamma^{\mathrm f}(y(\varepsilon))$ compatibly, for every $\varepsilon$, and the
system decomposes into its character components $\chi\in\ZZ/n$.

Put
\[
  r(x)=\frac{f(x)}{f(x^{-1})}\ (x\in G\smallsetminus\{1\}),\qquad r(1)=-1,\qquad
  \lambda_j=\sum_{x\in G}r(x)x^{-j}\quad(j\in\ZZ/n).
\]
If $f'(1)\ne0$ the value $r(1)=-1$ is the limit of $f(x)/f(x^{-1})$ at $x=1$. Let
$\Omega_n\subset\Sigma$ be the set of coefficient vectors such that
\begin{equation}\label{eq:denomopen}
\begin{gathered}
  a_1a_3a_7\ne0,\qquad f'(1)\ne0,\qquad f(x)\ne0\ \ (x\in G\smallsetminus\{1\}),\\
  \lambda_j\ne0\ \ (j\in\ZZ/n),\qquad \gcd\bigl(f(u),u^7f(u^{-1})\bigr)=u-1 .
\end{gathered}
\end{equation}

\begin{lemma}\label{lem:denom}
For every $n\ge4$ the set $\Omega_n$ is a nonempty Zariski open subset of $\Sigma$.
\end{lemma}

\begin{proof}
Each condition is the nonvanishing of a rational function or of a resultant, so it
suffices to show that each is not identically violated; since $\Sigma$ is
irreducible, finitely many nonempty open sets meet. The last condition holds for
$f=1+2u+3u^3-6u^7$: here $f'(1)=-31$, for $|u|>1$ we have
$6|u|^7>1+2|u|+3|u|^3$, and on $|u|=1$ equality in $|1+2u+3u^3|\le6$ forces $u=1$;
thus all roots other than $1$ lie in the open unit disc, and those of $u^7f(u^{-1})$
other than $1$ lie outside the closed disc.

Let $\zeta\in G$ be a primitive $n$-th root and let $\Lambda\subset\Sigma$ be the line
$f(\zeta)=0$; it is a line because $\zeta^3-\zeta\ne0$. If $\eta\in G\smallsetminus
\{1,\zeta\}$ and $f(\eta)$ vanished identically on $\Lambda$, the vector
$(1,\eta,\eta^3,\eta^7)$ would lie in the span of $(1,1,1,1)$ and
$(1,\zeta,\zeta^3,\zeta^7)$; comparing the first three coordinates gives
$(\eta-\zeta)(\eta+\zeta+1)=0$, so $\eta+\zeta=-1$, and two complex numbers of modulus
one with sum $-1$ are the primitive cube roots of unity, which is impossible for
$n\ge4$. In the same way $f'(1)\equiv0$ on $\Lambda$ would force
$\zeta^2+\zeta+1=3$. Hence a general point of $\Lambda$ has $f'(1)\ne0$ and no zero
in $G$ other than $1$ and $\zeta$. Near such a point, with transversal parameter
$\tau=f(\zeta)$, exactly one summand of $\lambda_j$ has a pole, namely
$x=\zeta^{-1}$, with residue $f(\zeta^{-1})\zeta^j\ne0$. So each $\lambda_j$ is a
nonzero rational function on $\Sigma$. Finally $a_1,a_3,a_7$ and $f(x)$ ($x\ne1$) are
nonzero functions.
\end{proof}

\subsection{Star pivots and the residual system}\label{sec:pivots}
From now on we assume $n\ge29$; Lemma~\ref{lem:denom} itself holds for all $n\ge4$.
At $\varepsilon=0$ write $C_{kl}=\ell_l(p_k)=f(t_k/t_l)$ and
$\Pi=\prod_{x\in G\smallsetminus\{1\}}f(x)=\prod_{l\ne k}C_{kl}$, which is independent
of $k$. The auxiliary points satisfy $E_{ab}(p_l)=0$ for $l\notin\{a,b\}$,
$E_{ab}(p_a)=\Pi/C_{ab}$ and $E_{ab}(o)=1$, while $V_{abk}(p_k)=\Pi/(C_{ka}C_{kb})$
and $V_{abk'}(p_k)=0$ for $k'\ne k$. The rows of \eqref{eq:extended} are of three
kinds: the \emph{pivot rows} $(k,ab)$ with $k\notin\{a,b\}$, the \emph{assigned rows}
$(a,ab)$ and the \emph{origin rows} $(o,ab)$. The pivot row $(k,ab)$ contains exactly
one triple unknown, $t_{abk,k}$, with the nonzero coefficient $-V_{abk}(p_k)$, and
each triple unknown occurs in exactly one pivot row. Solving,
\begin{equation}\label{eq:pivot}
  t_{abk,k}=\nu_k\cdot v_k-C_{ka}C_{kb}\Bigl[\frac{\beta_{ka}}{C_{ab}}
    +\frac{\beta_{kb}}{C_{ba}}+\frac{\beta_{ko}}{\Pi}\Bigr].
\end{equation}
\begin{definition}\label{def:special}
The linear system in the unknowns $(v,\beta)$ obtained by substituting \eqref{eq:pivot} in
the assigned and origin rows of $\Gamma^{\mathrm f}(y(0))$ is the \emph{residual system}
$\mathcal R$. Its kernel $\ker\mathcal R$ is the \emph{special kernel}.
\end{definition}

By construction
\begin{equation}\label{eq:residual}
  \ker\Gamma^{\mathrm f}(y(0))\cong\ker\mathcal R,\qquad
  \rank\Gamma^{\mathrm f}(y(0))=3L_d+\rank\mathcal R .
\end{equation}
The left kernel of $\Gamma^{\mathrm f}(y(0))$ is identified with that of $\mathcal R$ in the same way.

In the character $\chi$ the unknowns are
\[
  v_i=(V_1t_i^{\chi+1},V_3t_i^{\chi+3},V_7t_i^{\chi+7}),\quad
  \beta_{il}=t_i^\chi\beta(t_l/t_i),\quad \beta_{io}=\Pi c\,t_i^\chi,\quad
  \beta_{oi}=\eta\,t_i^\chi/\Pi,
\]
together with $\rho=\beta_{oo}$ if $\chi=0$ and the component $O_e$ of the origin
velocity if $\chi=-e$, $e\in\{1,3,7\}$. Here $\beta\colon G\to\CC$ is an arbitrary
function; its value $\beta(1)$ is the diagonal coefficient. Thus there are $n+5$
unknowns in each character, plus one in the four \emph{global} characters
$\chi\in\{0,-1,-3,-7\}$; in total $n^2+5n+4=3(n+1)+(n+1)^2$, as it should be. We put
$W_e=a_eV_e$, $\alpha=W_1+W_3+W_7$ and $\kappa_e=a_eO_e$.

Using rotation invariance, the assigned rows of the character $\chi$ are indexed by
the separation $x=t_b\in G\smallsetminus\{1\}$ of the edge $\{0,b\}$ with assigned end
$0$, and the origin rows by the edges $\{0,b\}$ up to the reversal
$b\leftrightarrow-b$ (for even $n$ the diameter $b=n/2$ contributes only in even
characters). Dividing the assigned row by $E_{0b}(p_0)$, the residual equations read,
with $y=t_k$ running over $G\smallsetminus\{1,x\}$,
\begin{align}
  &\sum_{y}\frac{\sum_{e>0}W_ey^{-e}-\alpha y^\chi}{f(y^{-1})}
   +\sum_y\frac{f(y)f(y/x)}{f(y^{-1})}\,y^\chi\Bigl[\frac{\beta(y^{-1})}{f(x^{-1})}
   +\frac{\beta(x/y)}{f(x)}+c\Bigr]\notag\\
  &\qquad-\beta(1)-\beta(x)\frac{f(x^{-1})}{f(x)}-c\,f(x^{-1})=0,\label{eq:assigned}\\
  &\sum_y\bigl(\nu_k\cdot v_o-\alpha y^\chi\bigr)
   +\sum_yf(y)f(y/x)\,y^\chi\Bigl[\frac{\beta(y^{-1})}{f(x^{-1})}
   +\frac{\beta(x/y)}{f(x)}+c\Bigr]\notag\\
  &\qquad-\frac{\eta}{f(x^{-1})}-\frac{\eta x^\chi}{f(x)}-\rho=0 .\label{eq:origin}
\end{align}
Here $\nu_k\cdot v_o=\kappa_ey^{-e}$ in the character $-e$ and is zero otherwise, and
$\rho$ occurs only for $\chi=0$.

\subsection{The assigned equations}
For $x\ne1$ put $h(x)=f(x^{-1})\beta(x)$ and set $h(1)=0$. Let
\[
  k_e=\sum_{x\ne1}x^{-\chi-e}\frac{h(x)}{f(x)}\ (e\in F),\qquad
  \Theta=\sum_{y\ne1}\frac{\sum_{e>0}W_ey^{-e}-\alpha y^\chi}{f(y^{-1})},\qquad
  \vartheta=\Theta-\beta(1),
\]
and let $\mathcal C_\chi h(x)=x^\chi\sum_{y\in G}r(x/y)y^{-\chi}h(y)$. For a function
$g$ on $G$ write $\hat g_k=\sum_xg(x)x^{-k}$; then
$\widehat{\mathcal C_\chi h}_k=\lambda_{k-\chi}\hat h_k$.

\begin{lemma}\label{lem:convolution}
Multiplied by $f(x)$, the assigned equation \eqref{eq:assigned} is equivalent to
\begin{multline}\label{eq:conv}
  \mathcal C_\chi h(x)+r(x)\sum_{e\in F}a_ek_ex^{-e}
  +c\,f(x)\sum_{e\in F}a_e\lambda_{-\chi-e}x^{-e}\\
  +\vartheta f(x)-r(x)\Bigl[\sum_{e>0}W_ex^{-e}-\alpha x^\chi\Bigr]=0 .
\end{multline}
The left side of \eqref{eq:conv} at $x=1$ equals $\sum_{e>0}W_e-\alpha=0$. Hence the
assigned equations are equivalent to \eqref{eq:conv} on all of $G$.
\end{lemma}

\begin{proof}
We treat the terms of \eqref{eq:assigned} after multiplication by $f(x)$. The first
sum is $f(x)\Theta$ minus the omitted term $y=x$, which gives
$-r(x)[\sum W_ex^{-e}-\alpha x^\chi]$; together with $-\beta(1)f(x)$ it gives the
terms with $\vartheta$ and $W$. In the second sum write
$f(y/x)=\sum_ea_ey^ex^{-e}$ and note $f(y)\beta(y^{-1})=h(y^{-1})$. The
$\beta(y^{-1})$ part becomes $r(x)\sum_ea_ex^{-e}\sum_{u\ne1,x^{-1}}u^{-\chi-e}h(u)/f(u)$;
the omitted $u=x^{-1}$ contributes a multiple of $\sum_ea_e=f(1)=0$, so this is
$r(x)\sum a_ek_ex^{-e}$. The $\beta(x/y)$ part, with $u=x/y$, is
$x^\chi\sum_{u\ne1,x}r(x/u)u^{-\chi}h(u)=\mathcal C_\chi h(x)+h(x)$, because
$h(1)=0$ and $r(1)=-1$; the extra $h(x)$ cancels the term $-\beta(x)f(x^{-1})$. The
$c$ part is $cf(x)\sum_ea_ex^{-e}\sum_{y\ne1,x}r(y)y^{\chi+e}$, and
$\sum_{y\ne1,x}r(y)y^{\chi+e}=\lambda_{-\chi-e}+1-r(x)x^{\chi+e}$. The term with
$r(x)$ contributes $x^\chi f(1)=0$, while the term $+1$ gives $cf(x)f(x^{-1})$, which
cancels the last term $-cf(x)f(x^{-1})$ of \eqref{eq:assigned}. This proves
\eqref{eq:conv} for $x\ne1$. At $x=1$ we have $f(1)=0$, $r(1)=-1$ and, since
$r(y^{-1})=\sum_ea_ey^{-e}/f(y)$,
$\mathcal C_\chi h(1)=\sum_{y\ne1}r(y^{-1})y^{-\chi}h(y)=\sum_ea_ek_e$; so the left
side is $\sum_{e>0}W_e-\alpha$.
\end{proof}

Taking Fourier coefficients in \eqref{eq:conv} and putting
$z_0=-k_0$, $z_e=W_e-a_ek_e$ ($e>0$) and $H_k=\hat h_k$, we obtain, with indices
modulo $n$ and $a_k$ the coefficient of $u^k$ in $f$ reduced modulo $n$,
\begin{equation}\label{eq:Hk}
  H_k=\frac{\sum_{e\in F}z_e\lambda_{k+e}-n\vartheta a_k-nc\,\mathcal B_k}{\lambda_{k-\chi}}
     -\alpha,\qquad
  \mathcal B_k=\sum_{u,v\in F,\ u-v=k}a_ua_v\lambda_{-\chi-v}.
\end{equation}

\begin{proposition}\label{prop:eight}
Let the coefficients lie in $\Omega_n$. Given
$(z_0,z_1,z_3,z_7,\alpha,\vartheta,c,\eta)$, define $H$ by \eqref{eq:Hk}, $h$ by inverse
Fourier transformation, $k_e$ as above and $W_e=z_e+a_ek_e$ for $e\in F$. Then
$h(1)=\sum_{e\in F}W_e-\alpha$, and the solutions of all assigned equations of the
character $\chi$ correspond bijectively to the data satisfying
\begin{equation}\label{eq:twoconstr}
  h(1)=0,\qquad W_0=z_0+k_0=0,
\end{equation}
the original unknowns being recovered by $V_e=W_e/a_e$, $\beta(x)=h(x)/f(x^{-1})$ for
$x\ne1$, $\beta(1)=\Theta-\vartheta$, with $c,\eta$ and the global unknowns unchanged.
\end{proposition}

\begin{proof}
By construction the Fourier transform of \eqref{eq:conv}, with $\sum_{e>0}W_ex^{-e}$
replaced by $\sum_{e\in F}W_ex^{-e}$, holds; evaluating its inverse transform at
$x=1$ as in Lemma~\ref{lem:convolution} gives
$h(1)=\sum_FW_e-\alpha$. Under \eqref{eq:twoconstr} we get $W_0=0$ and
$\alpha=W_1+W_3+W_7$, so \eqref{eq:conv} holds on $G$ and the reconstructed unknowns
solve the assigned equations by Lemma~\ref{lem:convolution}. Conversely a solution
determines all eight quantities, and \eqref{eq:Hk} is forced because every
$\lambda_{k-\chi}$ is nonzero.
\end{proof}

The Fourier transform of the identity $r(x)f(x^{-1})=f(x)$, valid on all of $G$, is
\begin{equation}\label{eq:R}
  \sum_{e\in F}a_e\lambda_{k+e}=n\,a_k\qquad(k\in\ZZ/n).
\end{equation}
Writing $z_e=z'_e+\vartheta a_e$, \eqref{eq:R} removes $\vartheta$ from
\eqref{eq:Hk}:
\begin{equation}\label{eq:Htilde}
  H_k=\tilde H_k-\alpha,\qquad
  \tilde H_k=\frac{\sum_{e\in F}z'_e\lambda_{k+e}-nc\,\mathcal B_k}{\lambda_{k-\chi}} .
\end{equation}
The inverse transform of the constant $\alpha$ is $\alpha$ times the indicator of
$x=1$. Hence $h(1)=0$ has the unique solution $\alpha=\tilde h(1)$; the $k_e$ do not
depend on $\alpha$ or $\vartheta$; and $W_0=0$ has the unique solution
$\vartheta=-z'_0-k_0$.

\begin{corollary}\label{cor:six}
For coefficients in $\Omega_n$ the assigned subsystem of every character has full
row rank, and its solution space is parametrized isomorphically by the six unknowns
$(z'_0,z'_1,z'_3,z'_7,c,\eta)$, together with $\rho'=\rho+n\alpha$ if $\chi=0$ and
$\kappa_e$ if $\chi=-e$.
\end{corollary}

\subsection{The origin equations}
For $j\in F$ put
\begin{equation}\label{eq:Aj}
  \mathcal A_j=\sum_{e\in F}z'_e\lambda_{\chi+j+e}-nc\,\mathcal B_{\chi+j}
     -\eta\,[j=0]\,\lambda_0 ,
\end{equation}
and for the character $\chi$ let
$S_\chi(x)=\sum_{u,v\in F,\,u+v+\chi\equiv0}a_ua_vx^{-v}$. The ten sums $u+v$ with
$u,v\in F$ are $0,1,2,3,4,6,7,8,10,14$. We call $\chi$ \emph{quadratic} if
$\chi\in\{-2,-4,-6,-8,-10,-14\}$, global if $\chi\in\{0,-1,-3,-7\}$, and
\emph{ordinary} otherwise.

\begin{proposition}\label{prop:origin}
Let $n\ge29$ and let the coefficients lie in $\Omega_n$. On the solution space of the
assigned equations, the origin equations of the character $\chi$ are equivalent to
the following systems:
\begin{enumerate}[label=\textup{(\roman*)}]
\item ordinary $\chi$: $\mathcal A_j=0$ for $j\in F$;
\item quadratic $\chi$: $\mathcal A_j=0$ for $j\in F$, and $c=0$;
\item $\chi=-e$, $e\in\{1,3,7\}$: $\mathcal A_j-(\kappa_e-nca_e)\lambda_j=0$ for
$j\in F$;
\item $\chi=0$: $\mathcal A_j-\tfrac12(\rho'-nc)\lambda_j=0$ for $j\in F$.
\end{enumerate}
\end{proposition}

\begin{proof}
We evaluate the sums in \eqref{eq:origin}. In the character $-e$,
$\sum_{y\ne1,x}\kappa_ey^{-e}=-\kappa_e(1+x^{-e})$, and
$\sum_{y\ne1,x}\alpha y^\chi=\alpha(n[\chi=0]-1-x^\chi)$. Using
$f(y)\beta(y^{-1})=h(y^{-1})$ and $f(y/x)\beta(x/y)=h(x/y)$, and the vanishing of
$f(1)$ and $h(1)$, the two $\beta$ sums are
$f(x^{-1})^{-1}\sum_ea_eH_{\chi+e}x^{-e}$ and $x^\chi f(x)^{-1}\sum_ea_eH_{\chi+e}x^e$;
the omitted terms contain $f(1)$ or $h(1)$. The $c$ sum is
$c\sum_{y\in G}f(y)f(y/x)y^\chi=nc\,S_\chi(x)$. Put
\[
  q_0=H_\chi+\alpha-\eta,\qquad q_e=H_{\chi+e}+\alpha\ (e>0),\qquad
  q(x)=\sum_{e\in F}a_eq_ex^e .
\]
Multiplying \eqref{eq:origin} by $f(x)f(x^{-1})$, the terms with $\alpha$ combine,
and we obtain on all of $G$ (at $x=1$ both sides vanish)
\begin{multline}\label{eq:originq}
  f(x)q(x^{-1})+x^\chi f(x^{-1})q(x)\\
   =f(x)f(x^{-1})\bigl[n\alpha[\chi=0]+\rho+\kappa_e(1+x^{-e})-nc\,S_\chi(x)\bigr],
\end{multline}
with $\rho$ and $\kappa_e$ present only in the global characters $0$ and $-e$.

If $\chi$ has a representative with $|\chi|\le14$, both sides of \eqref{eq:originq}
are Laurent polynomials whose supports lie in an interval of length at most $28<n$,
so \eqref{eq:originq} is an identity of Laurent polynomials. If $15\le\chi\le n-15$
the right side is zero (no global or quadratic character lies in this range) and the
supports $[-7,7]$ and $\chi+[-7,7]$ of the two terms on the left are disjoint modulo
$n$; hence $f(x)q(x^{-1})=0$ and $q=0$.

In the first case suppose the left side equals $f(x)f(x^{-1})\mathcal T(x)$. Then
$f(x)$ divides $x^\chi f(x^{-1})q(x)$ up to a monomial, and since
$\gcd(f,u^7f(u^{-1}))=u-1$, the polynomial $g=f/(1-u)$ of degree six divides $q$.
Write $q=(\mu'-\mu u)g$. From $f=(1-u)g$ and the vanishing of the coefficients of $f$ in
degrees $4,5,6$, the coefficients of $g$ in degrees $3,4,5,6$ are all equal to $-a_7$; as $q$
has no terms of degrees $4,5,6$, we get $(\mu'-\mu)a_7=0$, so $q=\mu(1-u)g=\mu f$.
Then $\mathcal T=\mu(1+x^\chi)$. For ordinary $\chi$, $S_\chi=0$ and $\mathcal T=0$,
so $q=0$. For quadratic $\chi$, $S_\chi$ is $a_e^2x^{-e}$ or $a_ea_{e'}(x^{-e}+x^{-e'})$
with $e,e'>0$, which is not a multiple of $1+x^\chi$; hence $c=0$ and $q=0$. For
$\chi=-e$ we have $S_\chi=a_e(1+x^\chi)$, so $q_j=\kappa_e-nca_e$ for all $j$. For
$\chi=0$, $S_0=1$ and $q_j=\tfrac12(n\alpha+\rho-nc)$. Finally
$q_j=\tilde H_{\chi+j}-\eta[j=0]$ by \eqref{eq:Htilde}, and
$\lambda_jq_j=\mathcal A_j$ by \eqref{eq:Htilde} and \eqref{eq:Aj}; dividing by
$\lambda_j\ne0$ gives the four systems.
\end{proof}

\subsection{The rank of the origin systems}
Let $\mathcal M_\chi=(\lambda_{\chi+j+e})_{j\in\{1,3,7\},\,e\in F}$, a $3\times4$
matrix. The column of $\eta$ in the systems of Proposition~\ref{prop:origin} is
$-\lambda_0\ne0$ in the row $j=0$ and zero in the other rows. Hence, if
$\rank\mathcal M_\chi=3$, the origin system has rank four in the ordinary and global
characters and rank five in the quadratic ones (the extra equation $c=0$ involves a
column not used by the rank-four minor).

\begin{theorem}\label{thm:special}
Let $n\ge29$. For coefficients in a nonempty open subset of $\Omega_n$ we have
$\rank\mathcal M_\chi=3$ for every $\chi\in\ZZ/n$. Consequently the special kernel
$\ker\mathcal R$ has dimension $2$ in each of the $n-10$ ordinary characters, $1$ in
each of the six quadratic characters and $3$ in each of the four global characters;
in total
\[
  \dim\ker\mathcal R=2(n-10)+6+3\cdot4=2n-2 .
\]
\end{theorem}

The rest of this subsection proves Theorem~\ref{thm:special}. For each $\chi$ it
suffices to show that some $3\times3$ minor of $\mathcal M_\chi$ is a nonzero rational
function on $\Sigma$, since there are finitely many characters. We restrict to the
curve
\[
  a=-\tfrac16,\qquad g_b(u)=1+au+au^2+b(u+u^2+\dots+u^6),
\]
\[
  f_b(u)=(1-u)g_b(u)=1+(a+b-1)u-au^3-bu^7 ,
\]
and study minors near $b=0$. The point $b=0$ is not in $\Omega_n$, since $a_7=0$ there;
we only use that a minor with a nonzero coefficient in its expansion at $b=0$ is not
identically zero. Put $q_b(u)=u^6g_b(u^{-1})=u^4h(u)+b\,q_1(u)$ with $h(u)=u^2+au+a$
and $q_1=1+u+\dots+u^5$, so that $r_b(u)=-u^7g_b(u)/q_b(u)$. At $b=0$ the poles of
$r_0$ are $v_1=1/2$ and $v_2=-1/3$, with residues $R_1=-21/160$ and
$R_2=-56/1215$. For small $b\ne0$ there are two poles near $v_1,v_2$ and four small
poles near $0$, all simple for $b$ in a punctured neighbourhood and all inside the
unit disc; sums over the four small poles are analytic in $b$.

\begin{lemma}\label{lem:poleformula}
Let $P_b$ be the polynomial part of $r_b$ at infinity (of degree at most seven) and
$L_j=\lambda_j/n$. For $0\le j<n$,
\[
  L_j=P_{b,j}+\sum_{\text{poles }v}\frac{R_v\,v^{n-j-1}}{1-v^n},
\]
where $R_v$ is the residue at $v$. For $j\ge8$ the polynomial term vanishes, and for
$b\ne0$ and $0\le j\le6$ the right side equals $\sum_vR_vv^{-j-1}/(1-v^n)$, which is
the same pole formula with $j$ replaced by $n+j$. These identities extend to $b=0$ by
continuity.
\end{lemma}

\begin{proof}
On $|u|=1$ expand $R/(u-v)=\sum_{k\ge1}Rv^{k-1}u^{-k}$ and average over $G$; this
includes $u=1$, where $r_b(1)=-1$ is the value of the rational function. For
$b\ne0$, $r_b$ vanishes to order seven at $0$, so its Taylor coefficients of degree
$j\le6$ vanish, which gives $P_{b,j}=\sum_vR_vv^{-j-1}$.
\end{proof}

\begin{lemma}\label{lem:smallpoles}
For an integer $k\ge-2$ let $\mathcal H_k(b)=\sum_{v\text{ small}}R_vv^{k-1}$. Then
\[
  \mathcal H_k(b)=\gamma_kb^\ell+O(b^{\ell+1}),\qquad \ell=\Bigl\lceil\frac{k+3}4
  \Bigr\rceil,\qquad \gamma_k=(-1)^{\ell+1}a^{-\ell-1}c_\varrho(a,\ell),
\]
where $\varrho=4\ell-k-3\in\{0,1,2,3\}$ and
\[
  c_0=1,\quad c_1=a-1,\quad c_2=\frac{(\ell+1)(a-1)}a,\quad
  c_3=\frac{\ell a^2-2\ell a+2\ell-2a+2}a .
\]
All four are nonzero at $a=-1/6$ for every $\ell\ge1$. Moreover the cyclic small-pole
sum $\sum_{v\text{ small}}R_vv^{k-1}/(1-v^n)=\sum_{m\ge0}\mathcal H_{k+mn}$ has the
same leading term for $n\ge29$.
\end{lemma}

\begin{proof}
The sum is the integral of $-u^{k+6}g_b(u)/q_b(u)$ over a small circle enclosing the
four small poles and no other pole; $u=0$ is not a pole since $k+6\ge4$. On this
circle
$1/q_b=\sum_{t\ge0}(-b)^tq_1^t/(u^{4t+4}h^{t+1})$, and the residue of the $t$-th term
is $-(-b)^t[u^{4t-k-3}](g_bq_1^t/h^{t+1})$, which vanishes for $t<\ell$. At $t=\ell$
only $g_0$ contributes, since the $b$-part of $g_b$ raises the order in $b$, and the
coefficient of $u^\varrho$ in $g_0q_1^\ell h^{-\ell-1}$ is $a^{-\ell-1}c_\varrho$ by
direct expansion through $u^3$. At $a=-1/6$, $c_1=-7/6$, $c_2=7(\ell+1)$ and the
numerator of $c_3$ is $85\ell/36+7/3>0$. The terms with $m\ge1$ have larger order in
$b$.
\end{proof}

\begin{proof}[Proof of Theorem~\ref{thm:special}]
Write $k=n-\chi$ and take rows $j\in\{1,3,7\}$.

\emph{The range $7\le\chi\le n-12$.} Take the columns $e\in\{0,1,7\}$. If
$\chi\le n-15$, every index $\chi+j+e$ lies in $[8,n-1]$; if
$n-14\le\chi\le n-12$ only the entry $(7,7)$ wraps, to $0,1,2$, and
Lemma~\ref{lem:poleformula} gives the same pole formula. The contribution of the
two old poles is an analytic matrix of rank at most two, equal at $b=0$ to
$D_1v_1^{-j-e}+D_2v_2^{-j-e}$ with $D_i=R_iv_i^{k-1}/(1-v_i^n)\ne0$. The cofactor
of the entry $(7,7)$ in this matrix is
\[
  \det\begin{pmatrix}v_1^{-1}&v_2^{-1}\\ v_1^{-3}&v_2^{-3}\end{pmatrix}D_1D_2
  \det\begin{pmatrix}1&v_1^{-1}\\1&v_2^{-1}\end{pmatrix}=(-30)(-5)D_1D_2=150D_1D_2 .
\]
By Lemma~\ref{lem:smallpoles} the small-pole correction at $(7,7)$, where the shift
$j+e=14$ is maximal, has order $\ell=\lceil(k-11)/4\rceil$, while all other
corrections (shifts at most $10$) have order at least $\ell+1$. Since the old-pole
matrix has determinant identically zero, the determinant of the minor is
\[
  150D_1D_2\gamma_{k-14}b^\ell+O(b^{\ell+1})\ne0 .
\]
For $n-14\le\chi\le n-12$ we have $k-14\in\{0,-1,-2\}$ and $\ell=1$.

\emph{The range $0\le\chi\le6$.} At $b=0$ the polynomial part is
$P_0=u^3/6+7u^2/36-203u/216-161/1296$, and
\[
  \partial_bP_b|_{b=0}=-u^7-\tfrac76u^6-\tfrac{49}{36}u^5-\tfrac{307}{216}u^4
  -\tfrac{1897}{1296}u^3-\tfrac{12811}{7776}u^2-\tfrac{41041}{46656}u
  +\tfrac{44093}{279936}.
\]
For the columns listed in Table~\ref{tab:nonneg} the polynomial part has a single
nonzero entry in the minor when $\chi\le2$, and the determinant at $b=0$ is
$\Delta_\chi D_1D_2$. For $3\le\chi\le6$ the polynomial part of $P_0$ does not enter, the
old-pole part has rank two, and all small-pole corrections have order at least three
since $n-(\chi+14)\ge9$; the coefficient of $b$ in the determinant comes from
$\partial_bP_b$ alone and equals $\Delta_\chi D_1D_2$. The values $\Delta_\chi$, obtained from
$3\times3$ determinants of these explicit entries and the old-pole cofactors, are
nonzero and independent of $n$.

\begin{table}[ht]
\caption{Minors for $0\le\chi\le6$ (rows $j=1,3,7$).}\label{tab:nonneg}
\begin{tabular}{cccr}
\toprule
$\chi$ & columns $e$ & order in $b$ & $\Delta_\chi$\\
\midrule
0&1,3,7&0&38329200\\
1&0,3,7&0&38329200\\
2&0,3,7&0&32853600\\
3&0,1,7&1&$-26044200$\\
4&0,1,3&1&119700\\
5&0,1,7&1&$-32853600$\\
6&0,1,3&1&$-421200$\\
\bottomrule
\end{tabular}
\end{table}

\emph{The range $\chi=n-h$, $1\le h\le11$.} Put $x=2^{-n}$ and $y=(-3)^{-n}$, and let
$t_j=[u^j]r_0(u)$, with $t_j=0$ for $j<0$. At $b=0$ every entry of the minors below is
\[
  L_{n-h+j+e}=\sum_{i=1,2}\frac{R_iv_i^{h-j-e-1}}{1-v_i^n}+t_{j+e-h},
\]
by Lemma~\ref{lem:poleformula} and the identity $P_{0,j}-\sum_iR_iv_i^{-j-1}=t_j$ at
wrapped indices. For the columns of Table~\ref{tab:neg}, the determinant is the
rational function $N_h/(K_h(x-1)(y-1))$ listed there. For $n\ge29$ we have
$0<x<10^{-3}$, $|y|<10^{-3}$ and $|y/x|=(2/3)^n<10^{-3}$, and each $N_h$ is nonzero:
for $h=1$ the perturbation of $-112x$ is less than $(2268+15120)x/1000<112x$; for
$h=2,3$ the perturbations are less than $250047x/1000<21952x$ and
$35721x/1000<3136x$; for $h=5$ the perturbation of $1015$ is less than
$(28672+413343)/1000<1015$; the remaining numerators are constants. This covers both
parities of $n$.

\begin{table}[ht]
\caption{Minors for $\chi=n-h$, $1\le h\le11$ (rows $j=1,3,7$), with
$x=2^{-n}$, $y=(-3)^{-n}$.}\label{tab:neg}
\begin{tabular}{cclr}
\toprule
$h$&columns $e$&$N_h$&$K_h$\\
\midrule
1&0,1,3&$15120xy-112x-2268y$&5\\
2&0,1,3&$21952x-250047y$&1620\\
3&0,1,3&$3136x-35721y$&1620\\
4&0,1,3&$-49$&11664\\
5&0,1,7&$1015-28672x-413343y$&233280\\
6&0,1,3&$343$&419904\\
7&0,1,3&$49$&419904\\
8&0,1,7&$343$&1259712\\
9&0,1,7&$343$&15116544\\
10&0,1,7&$343$&544195584\\
11&0,1,7&$49$&544195584\\
\bottomrule
\end{tabular}
\end{table}

The three ranges cover $\ZZ/n$. For fixed $n$ the finitely many nonzero minors define
nonempty open subsets of $\Sigma$, which meet each other and $\Omega_n$. The kernel
dimensions follow from the ranks four and five of the origin systems and from
Corollary~\ref{cor:six}: $6-4=2$, $6-5=1$ and $7-4=3$.
\end{proof}

% =====================================================================
\section{The first opening and the lower boundary}\label{sec:opening}

\subsection{A perturbation lemma}
\begin{lemma}\label{lem:perturb}
Let $M(\varepsilon)=M_0+\varepsilon M_1+O(\varepsilon^2)$ be a matrix of analytic
functions near $\varepsilon=0$. Let the columns of $N$ form a basis of $\ker M_0$ and
let $Q'$ be any matrix whose rows lie in the left kernel of $M_0$. Then
$\rank M(\varepsilon)\ge\rank M_0+\rank(Q'M_1N)$ for all sufficiently small
$\varepsilon\ne0$.
\end{lemma}

\begin{proof}
Let the rows of $Q$ form a basis of the left kernel; then $Q'=TQ$ and
$\rank Q'M_1N\le\rank QM_1N$. Choose invertible $U,V$ such that the last rows of $U$ are $Q$, the last columns of
$V$ are $N$, and
\[
  UM_0V=\begin{pmatrix}I&0\\0&0\end{pmatrix},\qquad
  UM(\varepsilon)V=\begin{pmatrix}I+\varepsilon X&\varepsilon Y\\
  \varepsilon Z&\varepsilon QM_1N+O(\varepsilon^2)\end{pmatrix}.
\]
The Schur complement of the upper left block is $\varepsilon(QM_1N+O(\varepsilon))$.
\end{proof}

We apply Lemma~\ref{lem:perturb} to $M(\varepsilon)=\Gamma^{\mathrm f}(y(\varepsilon))$,
character by character. By \eqref{eq:residual}, $N$ is the special kernel, lifted by
\eqref{eq:pivot}. A left kernel vector $\varpi$ of the residual system $\mathcal R$
lifts to one of $M_0$ by adding to it, on each pivot row $(k,ab)$, the multiple
$-\sum_j\varpi_{j,ab}V_{abk}(b_j)/V_{abk}(p_k)$, the sum running over the auxiliary
indices $j\in\ZZ/n\cup\{o\}$. Hence, if $S^1=M_1N$ denotes the first-order rows on special
kernel vectors, the relevant first-order map is $\varpi\circ\mathcal R^1$ with the
\emph{corrected} rows
\begin{equation}\label{eq:R1}
  \mathcal R^1_{j,ab}=S^1_{j,ab}-\sum_{k\notin\{a,b\}}
  \frac{V_{abk}(b_j)}{V_{abk}(p_k)}S^1_{k,ab}.
\end{equation}
For an assigned row this is divided by $E_{ab}(p_a)$; the derivative of this
normalization does not matter on $\ker\mathcal R$. For the origin row $V_{abk}(o)=1$, so
$\mathcal R^1_o=S^1_o-\sum_kS^1_k/V_{abk}(p_k)$.

\subsection{The support of the origin rows}
Put $\mathbb D=F-F$ (modulo $n$) and
\begin{equation}\label{eq:U}
  U_\chi=\mathbb D\cup(\chi+\mathbb D)\cup\bigcup_{u,v\in F,\ u+v+\chi\equiv0}(\mathbb D-v).
\end{equation}
We regard an origin row as a function of the separation $x\in G\smallsetminus\{1\}$,
the values at $x$ and $x^{-1}$ being related by the reversal of the edge, and we
extend it by zero at $x=1$.

\begin{lemma}\label{lem:support}
For every unknown vector of the character $\chi$ (not necessarily in the special
kernel), the function $f(x)f(x^{-1})\mathcal R_o(x)$ has Fourier support in $U_\chi$.
Consequently, for every $k\notin U_\chi$, the functional
$\mathcal R_o\mapsto\widehat{(f_+f_-\mathcal R_o)}_k$, where $f_\pm=f(x^{\pm1})$, lies
in the left kernel of $\mathcal R$.
\end{lemma}

\begin{proof}
The computation in the proof of Proposition~\ref{prop:origin} did not use the assigned
equations: the terms omitted from the $\beta$ sums contain $f(1)=0$, and
$h(1)=f(1)\beta(1)=0$ holds for every unknown vector. It shows that
$f_+f_-\mathcal R_o$ is a sum of $f(x)q(x^{-1})+x^\chi f(x^{-1})q(x)$ with $q$ supported on
$F$, of $\eta\,(f(x)+x^\chi f(x^{-1}))$, and of $f_+f_-$ times
$\alpha(1+x^\chi)$, constants, $\kappa_e(1+x^{-e})$ (with $x^{-e}=x^\chi$) and
$S_\chi(x)$. These have supports in $\mathbb D$, $\chi+\mathbb D$ and
$\bigcup(\mathbb D-v)$.
\end{proof}

\subsection{The opening on the special kernel}
On the special kernel let $\hat\beta_j=\sum_t\beta(t)t^{-j}$ and
$\mathsf Q_{uv}=\hat\beta_{\chi+u+v}$ for $u,v\in F$. Then
$(\mathsf Qa)_u=\sum_t h(t)t^{-\chi-u}=H_{\chi+u}$, so the classification in
Proposition~\ref{prop:origin} reads
\begin{equation}\label{eq:Qa}
  \mathsf Qa=\mathsf C\mathbf 1+\eta e_0,\qquad \mathsf C=\mathsf L-\alpha,
\end{equation}
where $\mathsf L=0$ in the ordinary and quadratic characters,
$\mathsf L=\kappa_e-nca_e$ in the character $-e$ and $\mathsf L=\frac12(n\alpha+\rho-nc)$
in the character $0$. Put
$\tilde{\mathsf Q}=\mathsf Q-\mathsf C(\mathbf1e_0^{\mathsf t}+e_0\mathbf1^{\mathsf t})
-\eta e_0e_0^{\mathsf t}$. Then $\tilde{\mathsf Q}a=0$, its block
$\mathsf S=(\mathsf Q_{ef})_{e,f\in\{1,3,7\}}$ coincides with that of $\mathsf Q$, and
\begin{equation}\label{eq:Qtilde}
  \sum_{u,v\in F}a_ua_v\tilde{\mathsf Q}_{uv}t^uz^v
   =\sum_{e,f>0}a_ea_f\mathsf S_{ef}(t^e-1)(z^f-1).
\end{equation}
Define regular functions on $G$ by $\varphi_e(t)=a_e(t^e-1)/f(t)$ for $t\ne1$ and
$\varphi_e(1)=a_ee/f'(1)$, $e\in\{1,3,7\}$; their sum is identically one. Put
\begin{equation}\label{eq:TS}
  T_{\mathsf S}(x)=\sum_{t\in G}t^\chi\varphi(t)^{\mathsf t}\mathsf S\varphi(t/x).
\end{equation}

\begin{proposition}\label{prop:openingidentity}
Let $n\ge29$ and let the coefficients lie in $\Omega_n$. On every special kernel
vector, the Fourier coefficients outside $U_\chi$ of $f_+f_-\mathcal R^1_o$ coincide
with those of $f_+f_-T_{\mathsf S}$.
\end{proposition}

\begin{proof}
We take the edge $\{0,b\}$ with separation $x=t_b$ and write $k$ for the internal
labels $t_k\ne1,x$. Call a function on $G$ \emph{$U_\chi$-supported} if after
multiplication by $f_+f_-$ its Fourier support lies in $U_\chi$; we must show that
$\mathcal R^1_o-T_{\mathsf S}$ is $U_\chi$-supported. The radial opening has velocity $-p_k$ at the
$k$-th point and $\nu_k\cdot(-p_k)=1$. Let $u_k=t_{0bk,k}$ be the pivot unknown of
this edge and, by \eqref{eq:pivot},
\begin{equation}\label{eq:Dk}
  D_k=\nu_k\cdot v_k-u_k=k^\chi f(k)f(k/x)\Bigl[\frac{\beta(1/k)}{f_-}
   +\frac{\beta(x/k)}{f_+}+c\Bigr].
\end{equation}
Put $\mathsf H=\sum_{t\ne1}1/f(t)$ and $V_k=V_{0bk}(p_k)$. Since
$\frac{d}{d\varepsilon}\ell_j((1-\varepsilon)p_k)=1-C_{kj}$, differentiation of the
products gives
\begin{equation}\label{eq:S1k}
  \frac{S^1_k}{V_k}=D_k\Bigl[\mathsf H-\frac1{f(k)}-\frac1{f(k/x)}-(n-3)\Bigr]
   +\sum_{l\ne k}\frac{\nu_l\cdot v_k-u_l}{f(k/l)}
   -\sum_j\beta_{kj}\frac{E^1_j}{V_k},
\end{equation}
where $E^1_j=\frac{d}{d\varepsilon}E_{0b}((1-\varepsilon)p_j)$, and
$S^1_o=-\sum_j\beta_{oj}E^1_j$ since the origin does not move. At the endpoints
$E^1_0/E_0=\mathsf H-1/f_--(n-2)$ and $E^1_x/E_x=\mathsf H-1/f_+-(n-2)$.

\emph{Interior terms.} The contribution of the interior $j$ to the $\beta$ sums in
$-\sum_kS^1_k/V_k$ and in $S^1_o$ is
\begin{equation}\label{eq:TQex}
  T^{\mathrm{ex}}_{\mathsf Q}(x)=\sum_{t\ne1,x}t^\chi
  \frac{\sum_{u,v\in F}a_ua_v\mathsf Q_{uv}t^{u+v}x^{-v}-\eta}{f(t)f(t/x)} .
\end{equation}

\emph{Assigned cancellations.} Write $I=G\smallsetminus\{1,x\}$ for the set of internal
labels and decompose $\nu_l\cdot v_k-u_l=D_l+\nu_l\cdot(v_k-v_l)$. In
$\sum_{k\in I}S^1_k/V_k$ the second sum of \eqref{eq:S1k} then splits into
$\sum_{l\in I}D_l\sum_{k\in I,k\ne l}1/f(k/l)$ and the double velocity sum
$\sum_{k\in I}\sum_{l\in I,l\ne k}\nu_l\cdot(v_k-v_l)/f(k/l)$. Hence, before the
endpoint terms, the coefficient of $D_k$ in $\sum_kS^1_k/V_k$ is
$2\mathsf H-(n-3)-1/f(k)-1/f(k/x)-1/f(1/k)-1/f(x/k)$. The two special assigned
equations at the endpoints $0$ and $x$ give exactly
\begin{align*}
  \sum_{k\in I}\frac{D_k}{f(1/k)}&=\beta(1)+\beta(x)\frac{f_-}{f_+}+cf_--\mathcal V_0,&
  \mathcal V_0&=\sum_{k\in I}\frac{\nu_k\cdot(v_0-v_k)}{f(1/k)},\\
  \sum_{k\in I}\frac{D_k}{f(x/k)}&=x^\chi\beta(1)+x^\chi\beta(1/x)\frac{f_+}{f_-}
     +cx^\chi f_+-\mathcal V_x,&
  \mathcal V_x&=\sum_{k\in I}\frac{\nu_k\cdot(v_x-v_k)}{f(x/k)} .
\end{align*}
Substituting these for the terms $-\sum_kD_k/f(1/k)$ and $-\sum_kD_k/f(x/k)$ adds
$\mathcal V_0+\mathcal V_x$ to $\sum_kS^1_k/V_k$, and these are exactly the terms $k=1$ and
$k=x$ missing from the double velocity sum. The velocity part of $\sum_kS^1_k/V_k$ is
therefore
\[
  \sum_{l\in I}\ \sum_{k\in G,\,k\ne l}\frac{\nu_l\cdot(v_k-v_l)}{f(k/l)} .
\]
Substituting $k=ly$ and using
$\nu_l\cdot(v_{ly}-v_l)=l^\chi\bigl(y^\chi\sum_{e>0}W_ey^e-\alpha\bigr)$, the inner sum is
$l^\chi J_\chi$ with $J_\chi=\sum_{y\ne1}(y^\chi\sum_{e>0}W_ey^e-\alpha)/f(y)$, and
$\sum_{l\in I}l^\chi=n[\chi=0]-1-x^\chi$. Thus the velocity part of $\sum_kS^1_k/V_k$ equals
$J_\chi(n[\chi=0]-1-x^\chi)$, in every character, including the global ones; it enters
$\mathcal R^1_o=S^1_o-\sum_kS^1_k/V_k$ with a minus sign. Since $f_+f_-$ is supported in
$\mathbb D$, this term is $U_\chi$-supported, whatever its sign.

\emph{Origin cancellation.} The special origin equation gives
\[
  \sum_kD_k=\alpha(n[\chi=0]-1-x^\chi)+\nu_0\cdot v_o+\nu_x\cdot v_o
   +\eta\Bigl(\frac1{f_-}+\frac{x^\chi}{f_+}\Bigr)+\rho,
\]
which is $U_\chi$-supported; also $\sum_kk^\chi f(k)f(k/x)=nS_\chi(x)$. This accounts for all
constant multiples of $\sum_kD_k$ and for the $c$ part of the endpoint $\beta$ sums. After
these steps the terms that are not yet known to be $U_\chi$-supported, besides
\eqref{eq:TQex}, are
\begin{multline}\label{eq:remaining}
  \sum_kD_k\Bigl[\frac1{f(k)}+\frac1{f(k/x)}\Bigr]
   -\sum_kk^\chi f(k)f(k/x)\Bigl[\frac{\beta(1/k)}{f_-^2}+\frac{\beta(x/k)}{f_+^2}\Bigr]\\
   +\beta(1)(1+x^\chi)+\beta(x)\frac{f_-}{f_+}+x^\chi\beta(1/x)\frac{f_+}{f_-}
   +c(f_-+x^\chi f_+)+\eta\Bigl(\frac1{f_-^2}+\frac{x^\chi}{f_+^2}\Bigr).
\end{multline}
In the first sum, cancelling $f(k)$ or $f(k/x)$ against \eqref{eq:Dk} allows the sum to
be extended over all $k$; the omitted endpoint terms are
$\beta(1)+\beta(x)f_-/f_++cf_-$ and $x^\chi\beta(1)+x^\chi\beta(1/x)f_+/f_-+cx^\chi f_+$,
which enter with a minus sign and cancel the three $\beta$ terms and the term
$c(f_-+x^\chi f_+)$ of \eqref{eq:remaining}; this $c$ term is cancelled, not discarded. The completed $c$ sums equal
$c\sum_kk^\chi(f(k/x)+f(k))=nc\sum_{e\in F,\,e+\chi\equiv0}a_e(x^{-e}+1)$, which vanishes
outside the global characters and is $U_\chi$-supported in them. The completed $\beta$ sums are
$U_\chi$-supported: the substitution $y=x/k$ turns them into $x^\chi$ times a fixed moment or a
polynomial supported on $F$, and the $\beta(1/k)$ sums are supported on $-F$; the
exterior factors $f_\pm$ then give $\mathbb D$ or $\chi+\mathbb D$. For the second
sum, \eqref{eq:Qa} gives
\[
  \sum_kk^\chi\beta(1/k)f(k)f(k/x)=\mathsf Cf_-+\eta,\qquad
  \sum_kk^\chi\beta(x/k)f(k)f(k/x)=x^\chi(\mathsf Cf_++\eta),
\]
where the omitted endpoints have a vanishing factor $f(1)$. The $\eta$ terms cancel
the last term of \eqref{eq:remaining}, and the $\mathsf C$ terms are $U_\chi$-supported. Hence
$\mathcal R^1_o-T^{\mathrm{ex}}_{\mathsf Q}$ is $U_\chi$-supported.

\emph{The affine part of $\mathsf Q$.} By definition of $\tilde{\mathsf Q}$, the
numerator in \eqref{eq:TQex} equals its $\tilde{\mathsf Q}$ version plus
$\mathsf C(f(t)+f(t/x))$, which contributes
$\mathsf C[(1+x^\chi)K_\chi-x^\chi/f_+-1/f_-]$ with $K_\chi=\sum_{t\ne1}t^\chi/f(t)$, a
$U_\chi$-supported function. By \eqref{eq:Qtilde} the remaining summand is
$t^\chi\varphi(t)^{\mathsf t}\mathsf S\varphi(t/x)$, which is regular also at $t=1$ and
$t=x$; adding these two terms changes the sum by
$\varphi(1)^{\mathsf t}\mathsf S\varphi(x^{-1})+x^\chi\varphi(x)^{\mathsf t}\mathsf S
\varphi(1)$, which after multiplication by $f_+f_-$ is supported in $\mathbb D\cup
(\chi+\mathbb D)$. Thus $\mathcal R^1_o-T_{\mathsf S}$ is $U_\chi$-supported.

Finally, every $U_\chi$-supported remainder obtained above contains a factor $f_+$ or $f_-$ that
is not divided out, so it vanishes at $x=1$, as does $f_+f_-T_{\mathsf S}$; hence the
Fourier transformation of functions extended by zero at $x=1$ introduces no extra
term.
\end{proof}

\subsection{Three moments and the detector}
Since $\sum_e\varphi_e=1$, a matrix $\mathsf S$ of the form
$\mathbf1w^{\mathsf t}+w\mathbf1^{\mathsf t}$ contributes to $\widehat{T_{\mathsf S}}$
only at the indices $0$ and $\chi$, hence only to $U_\chi$ after multiplication by
$f_+f_-$. Eliminating $\varphi_7=1-\varphi_1-\varphi_3$, only the symmetric
$2\times2$ matrix $\mathsf m$ with entries
\begin{equation}\label{eq:moments}
\begin{aligned}
  \mathsf m_{11}&=\hat\beta_{\chi+2}-2\hat\beta_{\chi+8}+\hat\beta_{\chi+14},\\
  \mathsf m_{13}&=\hat\beta_{\chi+4}-\hat\beta_{\chi+8}-\hat\beta_{\chi+10}
     +\hat\beta_{\chi+14},\\
  \mathsf m_{33}&=\hat\beta_{\chi+6}-2\hat\beta_{\chi+10}+\hat\beta_{\chi+14}
\end{aligned}
\end{equation}
is visible outside $U_\chi$. Let $\hat\varphi(k)=(\hat\varphi_1(k),\hat\varphi_3(k))$ with
$\hat\varphi_e(k)=\sum_t\varphi_e(t)t^{-k}$. Then $\widehat{T_{\mathsf S}}_k=\hat\varphi(k-\chi)^{\mathsf t}
\mathsf m\hat\varphi(-k)$ modulo the two invisible indices, and the outside-support
coefficients of $f_+f_-T_{\mathsf S}$ are given by the \emph{detector}
\begin{equation}\label{eq:detector}
  (\mathcal O_\chi\mathsf m)_k=\sum_{u,v\in F}a_ua_v\,\hat\varphi(k-u+v-\chi)^{\mathsf t}\,
   \mathsf m\,\hat\varphi(-k+u-v),\qquad k\notin U_\chi .
\end{equation}
Let $\mathcal W_\chi$ denote the linear map from the six- or seven-dimensional unknown
space of Corollary~\ref{cor:six} to $(\mathsf m_{11},\mathsf m_{13},\mathsf m_{33})$, and
$\mathcal P_\chi$ the origin system of Proposition~\ref{prop:origin}. Combining
Lemma~\ref{lem:perturb}, Lemma~\ref{lem:support} and
Proposition~\ref{prop:openingidentity}, the rank of the first opening on the special
kernel of $\chi$ is at least, writing $[\mathcal P;\mathcal X]$ for the matrix obtained by
placing the rows of $\mathcal X$ below those of $\mathcal P$,
\begin{equation}\label{eq:openingbound}
  \rank\bigl(\mathcal O_\chi\mathcal W_\chi|_{\ker\mathcal P_\chi}\bigr)
  =\rank\begin{pmatrix}\mathcal P_\chi\\ \mathcal O_\chi\mathcal W_\chi\end{pmatrix}
   -\rank\mathcal P_\chi .
\end{equation}

\begin{lemma}\label{lem:momentformula}
Let $P_{11}=(u-u^7)^2$, $P_{13}=(u-u^7)(u^3-u^7)$, $P_{33}=(u^3-u^7)^2$ and
$\Psi_A(j)=\sum_tt^{-j}P_A(t)/f(t)$, with value $0$ at $t=1$. In the coordinates of
Corollary~\ref{cor:six},
\begin{equation}\label{eq:momentformula}
  \mathsf m_A=\frac1n\sum_{j\in\ZZ/n}\frac{\Psi_A(j)}{\lambda_j}
  \Bigl[\sum_{e\in F}z'_e\lambda_{\chi+j+e}-nc\,\mathcal B_{\chi+j}\Bigr].
\end{equation}
\end{lemma}

\begin{proof}
For instance $\mathsf m_{11}=\sum_t\beta(t)t^{-\chi}(t^{-1}-t^{-7})^2$; the term $t=1$
vanishes, and $\beta(t)=h(t)/f(t^{-1})$ otherwise. Inserting
$h(t)=\frac1n\sum_kH_kt^k$ gives $\mathsf m_A=\frac1n\sum_jH_{\chi+j}\Psi_A(j)$. By
\eqref{eq:Htilde} the constant $-\alpha$ contributes $\sum_j\Psi_A(j)=0$, and the rest
is \eqref{eq:momentformula}.
\end{proof}

The required opening ranks are $2$ in the ordinary and
global characters and $1$ in the quadratic characters; equivalently
$\rank[\mathcal P_\chi;\mathcal O_\chi\mathcal W_\chi]=6$ in every character. In a global character the remaining
one-dimensional kernel is spanned by a translation or by the dilation: in the
character $-e$ the vector with $z'=u^e$ (that is, $z'_e=1$) and $\kappa_e=1$, and in
the character $0$ the vector with $z'=1$ and $\rho'=2$, have zero moments.

\subsection{The detector is injective for \texorpdfstring{$n\ge89$}{n >= 89}}
\begin{proposition}\label{prop:detector}
Let $n\ge89$. On a nonempty open subset of $\Omega_n$, $\mathcal O_\chi$ is injective on
symmetric $2\times2$ matrices for every $\chi$.
\end{proposition}

\begin{proof}
Write $g=f/(1-u)$. We add to $\Omega_n$ the open conditions that the six roots
$\alpha_1,\dots,\alpha_6$ of $g$ are simple and nonzero, that the thirty ratios
$\alpha_i/\alpha_j$ ($i\ne j$) are distinct, and that none of them is a root of
$f(u)f(u^{-1})$. These conditions are nonempty: for $g=1+3u+3u^2+6u^3+6u^4+6u^5+6u^6$
the polynomial $\operatorname{Res}_z(g(z),g(uz))/(u-1)^6$ has degree $30$, is
squarefree and is coprime to $f$ and to its reciprocal, as one verifies modulo $101$;
a nonzero reduction of these resultants certifies the corresponding characteristic
zero conditions.

The functions $\varphi_1=-a_1/g$ and $\varphi_3=-a_3(1+u+u^2)/g$ are proper rational
functions, and for $0\le k<n$
\[
  \hat\varphi(k)=\sum_iU_i\alpha_i^{-k},\qquad
  U_i=-\frac{n\alpha_i^{n-1}}{(1-\alpha_i^n)g'(\alpha_i)}
   \bigl(a_1,\ a_3(1+\alpha_i+\alpha_i^2)\bigr)^{\mathsf t},
\]
by the same averaging as in Lemma~\ref{lem:poleformula}. Choose one of the intervals
$8\le k\le\chi-8$ or $\chi+8\le k\le n-8$. On it every shifted index $k-u+v$ in
\eqref{eq:detector} lies strictly between $0$ and $\chi$, or strictly between $\chi$ and
$n$, so \eqref{eq:detector} is a sum of the thirty exponentials
$(\alpha_j/\alpha_i)^k$: the diagonal terms are killed by $f(1)^2=0$, and the
coefficient of $(\alpha_j/\alpha_i)^k$ is a nonzero multiple of
$U_i^{\mathsf t}\mathsf mU_j\,f(\alpha_i/\alpha_j)f(\alpha_j/\alpha_i)$, the wrap
factor being $\alpha_i^\chi\alpha_j^{-n}$, or $\alpha_i^{\chi-n}\alpha_j^{-n}$ on the
lower interval. These intervals avoid $\mathbb D\cup(\chi+\mathbb D)$; in a quadratic
character $\chi=n-(u+v)$ the extra sets $\mathbb D-v$ avoid $8\le k\le\chi-8$, and in a
global character there is no extra set. The longer interval has
$\max(\chi,n-\chi)-15\ge30$ elements for $n\ge89$, so if $\mathcal O_\chi\mathsf m=0$ the
Vandermonde matrix of thirty consecutive indices forces
$U_i^{\mathsf t}\mathsf mU_j=0$ for all $i\ne j$. At least three of the projective
points $[a_1:a_3(1+\alpha_i+\alpha_i^2)]$ are distinct, since a quadratic polynomial
takes each value at most twice. A symmetric bilinear form on $\CC^2$ vanishing on
all pairs of three pairwise independent vectors is zero; hence $\mathsf m=0$.
\end{proof}

\subsection{The fixed matrices}\label{sec:fixed}
For each character $\chi$ we regard $\mathcal P_\chi$ (the origin system of
Proposition~\ref{prop:origin}, in the columns $z'_0,z'_1,z'_3,z'_7,c,\eta$ and, in a global
character, $\kappa_e$ or $\rho'$) and $\mathcal W_\chi$ (the three rows
\eqref{eq:momentformula}) as matrices whose entries are functions of the coefficient
vector $a\in\Sigma$. Their entries are polynomials in $n$, $a_1,a_3,a_7$, the $\lambda_j$,
the $\lambda_j^{-1}$ and the values $\Psi_A(j)=\sum_{t\ne1}t^{-j}P_A(t)/f(t)$; no entry
involves a division by $a_1$, $a_3$ or $a_7$. Hence both matrices are regular on the open
set $\Sigma'\subset\Sigma$ where $f(x)\ne0$ for $x\in G\smallsetminus\{1\}$ and
$\lambda_j\ne0$ for all $j$, and every minor is a rational function on $\Sigma$, regular on
$\Sigma'$. The divisions by $a_e$ occur only in the passage from the six or seven fixed
coordinates back to the original unknowns ($V_e=W_e/a_e$, $O_e=\kappa_e/a_e$), that is, in
Corollary~\ref{cor:six}, which is applied on $\Omega_n=\Sigma'\cap\{a_1a_3a_7\ne0,\
f'(1)\ne0,\ \gcd(f,u^7f(u^{-1}))=u-1\}$. Consequently, if a minor of
$[\mathcal P_\chi;\mathcal W_\chi]$ (or of $[\mathcal P_\chi;\mathcal O_\chi\mathcal W_\chi]$)
is nonzero at one point of $\Sigma'$, or has a nonzero coefficient in its expansion along
an analytic curve whose general point lies in $\Sigma'$, then it is a nonzero rational
function on $\Sigma$, hence nonzero on a dense open subset of $\Sigma$, which meets
$\Omega_n$; at the points of this intersection the rank statement is a statement about the
special kernel. In the next two subsections we use two such curves: the \emph{cubic line}
$a=(-\frac76+b,\frac16,-b)$, which lies in $\Omega_n$ for small $b\ne0$ and whose point
$b=0$ lies in $\Sigma'$ but not in $\Omega_n$, and the \emph{six-small-pole line}
$a=(-1-\epsilon,0,\epsilon)$, which lies entirely in the plane $\{a_3=0\}$ and hence
outside $\Omega_n$, but whose general point lies in $\Sigma'$.

\subsection{The bulk characters}
\begin{proposition}\label{prop:bulk}
Let $n\ge29$ and $7\le\chi\le n-16$. On a nonempty open subset of $\Sigma$,
$\rank[\mathcal P_\chi;\mathcal W_\chi]=6$.
\end{proposition}

All characters in this range are ordinary. We use the same curve as before,
\[
  g_b(t)=1-\tfrac t6-\tfrac{t^2}6+b(t+\dots+t^6),\quad f_b=(1-t)g_b,\quad
  a_1=-\tfrac76+b,\ a_3=\tfrac16,\ a_7=-b,
\]
and study the fixed matrices near $b=0$; the point $b=0$ lies in $\Sigma'$, as shown below,
so by Section~\ref{sec:fixed} a minor with a nonzero coefficient in its expansion at $b=0$ is
a nonzero rational function on $\Sigma$. Put
$K=n-\chi\ge16$, $\ell=\lceil(K-15)/4\rceil\ge1$ and $\varrho=4\ell-K+15\in\{0,1,2,3\}$.

\emph{Fourier denominators.} At $b=0$, $g_0=(1-t/2)(1+t/3)$; the old poles of $r_0$ are
$v_1=1/2$, $v_2=-1/3$ with residues $R_1=-21/160$, $R_2=-56/1215$, and the polynomial
part is $P_0=t^3/6+7t^2/36-203t/216-161/1296$. For $4\le j<n$,
$\lambda_j=n\sum_iR_iv_i^{n-j-1}/(1-v_i^n)$, which is nonzero because the ratio of the
second summand to the first is at most $(1792/5103)/(1-3^{-29})<1/2$; for $j\le3$ the
nonzero coefficient $P_{0,j}$ dominates the two tails, each of modulus less than
$10^{-4}$. The values $f_0(x)$, $x\ne1$, are nonzero since $2$ and $-3$ are not roots of
unity. Put
\[
  D_1=\frac{nR_1v_1^{\chi-1}}{1-v_1^n},\quad
  T=\frac{R_2}{R_1}\Bigl(-\frac23\Bigr)^{\chi-1}\frac{1-2^{-n}}{1-(-3)^{-n}},\quad
  d_k=2^k+T(-3)^k .
\]
Then $\lambda_{K+k}=D_1d_k$ whenever $4\le K+k\le n-1$, and in the bulk range
\begin{equation}\label{eq:Tbound}
  0<|T|\le\frac{1792}{5103}\Bigl(\frac23\Bigr)^6\frac1{1-3^{-29}}<\frac1{30},
\end{equation}
uniformly in $n$ and $\chi$; in particular $d_k>0$ for $k<0$.

\emph{The limiting kernel.} Factor $g_b=G_bH_b$ analytically, where $G_b$ is the
quadratic factor with roots near $2,-3$ and $H_b$ the quartic factor whose roots tend to
infinity, both with constant term one. Let $Z_b=t^7-\operatorname{rem}_{G_b}(t^7)$.
A polynomial supported on $F$ and divisible by $G_b$ is uniquely of the form
$z=\mathsf A f_b+\mathsf B Z_b$. At $b=0$,
\begin{gather*}
  Z_0=t^7-463t+798,\\
  l_Z(t)=-tg_0(t)\,\frac{Z_0}{g_0}(t^{-1}),\\
  l_Z(t)=6t^{-4}-7t^{-3}+42t^{-2}-84t^{-1}+336-840t+78t^2+133t^3 .
\end{gather*}
The three origin rows $j=1,3,7$ have rank two at $b=0$ and require exactly that $z$
be divisible by $g_0$; their $c$ coefficients vanish since all $\chi+j$ lie strictly
between $7$ and $n-7$; and the row $j=0$ determines $\eta$. On the resulting
three-dimensional space with coordinates $(\mathsf A,\mathsf B,c)$, the numerators in
\eqref{eq:momentformula} are the Fourier transforms of $r_0(t)z(t^{-1})=\mathsf Af_0+
\mathsf Bl_Z$, supported in $[-4,3]$, and of the $c$ term, supported in $[-3,3]$; so only
the indices $j=K+k$, $-4\le k\le3$, occur, all between $12$ and $n-4$, and the
polynomial parts of $P_A/f_0$ (of degree at most $11$) do not contribute. The two old
poles give two moment conditions
\begin{equation}\label{eq:twooldmoments}
  \mathcal A_i\mathsf A+\mathcal B_i\mathsf B-\mathcal C_i(cD_1)=0\qquad(i=1,2),
\end{equation}
with, for $\alpha_1=2$, $\alpha_2=-3$,
\[
  \mathcal A_i=\sum_k\frac{(f_0)_k\alpha_i^{-k}}{d_k},\quad
  \mathcal B_i=\sum_k\frac{(l_Z)_k\alpha_i^{-k}}{d_k},\quad
  \mathcal C_i=\sum_{u,v\in\{0,1,3\}}a_ua_v\frac{d_{-v}\alpha_i^{-u+v}}{d_{u-v}} .
\]
The two residue vectors are independent, since $t^2+t^4$ takes the values $20$ and
$90$ at $2$ and $-3$. A direct rational simplification gives
\begin{equation}\label{eq:DAC}
  \det(\mathcal A,\mathcal C)=\frac{-9475375\,T(8736T^2+1277T-1134)}
  {10077696(T+1)(3T-2)(4T+9)(9T+4)(27T-8)},
\end{equation}
nonzero under \eqref{eq:Tbound}. So \eqref{eq:twooldmoments} has rank two, every
nonzero limiting kernel vector has $\mathsf B\ne0$, and
\begin{equation}\label{eq:limkernel}
  \frac{\mathsf A}{\mathsf B}=-9870,\qquad \frac{cD_1}{\mathsf B}=\mathsf c(T),
\end{equation}
where
\[
  \mathsf c(T)=\frac{-7776\,\mathsf p(T)}{175T(2T-3)(8T-27)(16T+81)(8736T^2+1277T-1134)}
\]
and
\begin{multline*}
  \mathsf p(T)=609140736T^6-2807074304T^5-4400288360T^4+13769852070T^3\\
  -1887377085T^2-2495601009T+589545216 .
\end{multline*}
These identities are obtained by forming the finite sums in \eqref{eq:twooldmoments}.
Thus the extended matrix has rank exactly five at $b=0$, with this one-dimensional
kernel.

\emph{A moving third moment.} Let $\alpha_1(b),\alpha_2(b)$ be the roots of $G_b$,
$\mathsf b_i=\alpha_i^2+\alpha_i^4$, and
\[
  P_\perp(t,b)=t^2(1-t^2)^2\bigl(t^2+t^4-\mathsf b_1\bigr)\bigl(t^2+t^4-\mathsf b_2\bigr),
\]
which lies in the span of $P_{11},P_{13},P_{33}$ (at $b=0$ with coefficients
$(1800,-3710,1911)$) and vanishes at both old roots. Then
$P_\perp/f_b=N_b/H_b$ with $N_b$ an analytic polynomial of degree at most $11$, and
\[
  N_0=6t^{11}+42t^9-36t^8-378t^7-504t^6-1770t^5-1260t^4+2100t^3+1800t^2 .
\]
The coefficient of $b^{\ell'}$ in $P_\perp/f_b$ has degree at most $11+4\ell'$; for
$z=\mathsf Af_b+\mathsf BZ_b$,
$r_b(t)z(t^{-1})=-tg_b(t)[z/G_b](t^{-1})/H_b(t^{-1})$ has $b^{m}$ coefficient with
exponents between $-4-4m$ and $7$; and the $c$ numerator has lowest exponent $-3$ at
order zero and $-7$ from order one on. In \eqref{eq:momentformula} a term of orders
$\ell'$ and $m$ can pair only if $K-4-4m\le j\le11+4\ell'$, which forces
$\ell'+m\ge\ell$. Hence the moving moment vanishes through order $\ell-1$ on the whole
family, and the positive origin equations vanish on it through order $\ell$, their
first possible alias having order $\lceil(K-11)/4\rceil=\ell+1$.

\emph{Leading coefficients.} The four highest nonconstant coefficients of
$\partial_bH_b|_0$ are $-6t^4+0t^3-42t^2+36t$, and one obtains for $j=0,1,2,3$
\begin{align*}
  [b^{\ell'}t^{11+4\ell'-j}]\,\frac{P_\perp}{f_b}&=6^{\ell'+1}\bigl(1,0,7(\ell'+1),
  -6(\ell'+1)\bigr)_j,\\
  [b^mt^{-4-4m+j}]\,r_bz(t^{-1})&=\mathsf B\,6^m\bigl(6,-7,42(m+1),-84-85m\bigr)_j ;
\end{align*}
the $\mathsf A$ part begins four exponents later. The $c$ terms that can occur at this
order are $a_3\lambda_{K-3}$, $a_1a_3\lambda_{K-3}$, $a_1\lambda_{K-1}$ at the exponents
$-3,-2,-1$ (order zero) and $-\lambda_{K-7}$, $-a_1\lambda_{K-7}$, $0$ at $-7,-6,-5$
(order one). For $\varrho=1$ the first pair cancels because $6a_3=1$; for $\varrho=3$
the surviving difference cancels because $6a_1=-7$; for $\varrho=0$ there is no $c$
term; for $\varrho=2$ the retained term is
$-nca_16^\ell[\lambda_{K-3}/\lambda_{K-2}-\lambda_{K-7}/\lambda_{K-6}]$. Normalizing
$\mathsf B=1$ and substituting \eqref{eq:limkernel}, the coefficient of $b^\ell$ of the
moving moment is $n6^\ell\mathsf S_\varrho(T,\ell)/D_1$, where, with sums over
$m=0,\dots,\ell$,
\begin{align*}
  \mathsf S_0&=36\sum\frac1{d_{-4-4m}},\qquad
  \mathsf S_1=-42\sum\frac1{d_{-3-4m}},\\
  \mathsf S_2&=252\sum\Bigl[\frac{\ell-m+1}{d_{-4-4m}}+\frac{m+1}{d_{-2-4m}}\Bigr]
   +\frac76\mathsf c(T)\Delta(T),\\
  \mathsf S_3&=-6\sum\Bigl[\frac{36(\ell-m+1)}{d_{-4-4m}}
   +\frac{49(\ell-m+1)}{d_{-3-4m}}+\frac{84+85m}{d_{-1-4m}}\Bigr],
\end{align*}
where
\[
  \Delta=\frac{d_{-3}}{d_{-2}}-\frac{d_{-7}}{d_{-6}}=\frac{-1950T}{(4T+9)(64T+729)} .
\]
All $d_k$ occurring here are positive, so $\mathsf S_0>0$ and $\mathsf S_1,\mathsf S_3<0$.
For $\mathsf S_2$, $1/d_k\ge\frac{30}{31}2^{-k}$, and the terms $m=0,1$ alone contribute
more than $252\cdot420\cdot30/31>102000$, while the numerator of
$\frac76\mathsf c\Delta$ is $101088\,\mathsf p(T)$ with $|\mathsf p(T)|<7\cdot10^8$ for
$|T|<1/30$, and its denominator
$(2T-3)(4T+9)(8T-27)(16T+81)(64T+729)(8736T^2+1277T-1134)$ has factors of absolute
values greater than $2,8,26,80,726,1080$; hence $|\frac76\mathsf c\Delta|<2712$ and
$\mathsf S_2>0$.

\begin{proof}[Proof of Proposition~\ref{prop:bulk}]
At $b=0$ choose five independent rows and columns: two positive origin rows, the
$\eta$ row and the two old moment rows; their pivot minor stays invertible for small
$b$. In the family $z=\mathsf A(b)f_b+Z_b$ solve the two chosen moment equations for
$\mathsf A(b),c(b)$, using \eqref{eq:DAC}, and solve the $\eta$ row. The resulting
vector satisfies the chosen origin equations to order $\ell$, so its correction to the
exact five-pivot kernel changes it only at order $\ell+1$. The moving moment vanishes
on the family through order $\ell-1$, and its coefficient of $b^\ell$ depends only on
$\mathsf A(0)=-9870$, $\mathsf B=1$ and $c(0)D_1=\mathsf c(T)$; it is the nonzero number
above. Hence the Schur complement of a sixth row with respect to the five pivots is
nonzero at order $\ell$, and some $6\times6$ minor has a nonzero coefficient of
$b^\ell$.
\end{proof}

\subsection{The boundary characters}
\begin{proposition}\label{prop:boundary}
Let $n\ge89$ and
$\chi\in\{0,1,\dots,6\}\cup\{n-15,\dots,n-1\}$. On a nonempty open subset of
$\Sigma$, $\rank[\mathcal P_\chi;\mathcal W_\chi]=6$.
\end{proposition}

These are $22$ characters. We write $\chi=-h$ for the negative ones and use the two
lines of Section~\ref{sec:fixed}: the six-small-pole line for
$h\in\{1,2,\dots,10,14\}$ and $\chi=0$, and the cubic line for
$\chi\in\{1,\dots,6\}$ and $h\in\{11,12,13,15\}$. Throughout we use the normalized
quantities $l_j=\lambda_j/n$, $\psi_A=\Psi_A/n$ and $\bar c=nc$; in the character $-h$
the moment rows are
\[
  \mathsf W_A=\sum_{j\in\ZZ/n}\psi_A(j)\,\frac{\sum_ez'_el_{j-h+e}-\bar c\,\bar{\mathcal B}_j}
  {l_j},\qquad \bar{\mathcal B}_j=\sum_{u-v=j-h}a_ua_vl_{h-v},
\]
and the origin rows are the numerators for $j=1,3,7$, with $(\bar ca_h-\kappa)l_j$ added
for $h\in\{1,3,7\}$, $(\bar c-\rho')l_j/2$ added for $h=0$, and the extra equation
$\bar c=0$ for the quadratic $h$. The row $j=0$ has coefficient $-l_0$ on $\eta$ and
determines $\eta$.

\emph{The six-small-pole line.} Put $a_1=-1-\epsilon$, $a_3=0$, $a_7=\epsilon$, so that
$f_\epsilon=(1-u)(1-\epsilon A(u))$ with $A=u+\dots+u^6$; with
$Q(u)=u^{-1}+\dots+u^{-6}$ we have $r_\epsilon(u)=-u(1-\epsilon A(u))/(1-\epsilon Q(u))$.
Let $\mathsf L_m(\epsilon)$ be the coefficient of $u^m$ in
\[
  -u+\sum_{q\ge1}\epsilon^q\bigl[-uQ(u)^q+uA(u)Q(u)^{q-1}\bigr],
\]
a Laurent polynomial at each order, with exponents at most $7$. Then
$l_j=\sum_{k\in\ZZ}\mathsf L_{j+kn}$; in particular $l_1=-1+O(\epsilon)$,
$l_0=-\epsilon+O(\epsilon^2)$ and $l_j=\epsilon+O(\epsilon^2)$ for $2\le j\le7$. For
$8\le j<n$ the valuation of $l_j$ is $q_j=\lceil(n-j+1)/6\rceil$ and its leading
coefficient is $-\binom{q_j+\varrho_j-1}{\varrho_j}$ with
$\varrho_j=6q_j-n+j-1\in\{0,\dots,5\}$. So every $l_j$ is a unit in
$\QQ((\epsilon))$; moreover $f_\epsilon$ does not vanish on $G\smallsetminus\{1\}$ near
$\epsilon=0$ and $f'_\epsilon(1)=-(1-6\epsilon)$. Write $n=6\mathsf N+\mathsf s$,
$0\le\mathsf s\le5$, $\mathsf N\ge14$. For fixed $j\ge8$ and
$\delta_j=\lceil(\mathsf s-j+1)/6\rceil$, before a second alias enters, the coefficient
of $\epsilon^{\mathsf N+\delta}$ in $l_j$ is the polynomial in $\mathsf N$
\[
  d_{j,\delta}(\mathsf N)=-C(\mathsf N+\delta,6\delta-\mathsf s+j-1)
  +\sum_{a=1}^6C(\mathsf N+\delta-1,6\delta-\mathsf s+j-a-7),
\]
where $C(q,\mathsf d)=0$ for $\mathsf d<0$ and otherwise
$C(q,\mathsf d)=\sum_{i\le\mathsf d/6}(-1)^i\binom qi\binom{q+\mathsf d-6i-1}{\mathsf d-6i}$;
this follows by reversing the powers in $Q^q$. Finally
$P_A/f_\epsilon=p_A(u)\sum_q\epsilon^qA(u)^q$ with $p_A=P_A/(1-u)$ of degree at most
$13$.

Replacing every $l_m$ in the positive origin rows by $\mathsf L_m$ gives a formal
$3$-row matrix over $\QQ((\epsilon))$. For $h=2,\dots,10,14$ choose the first
three-column minor of least valuation (Table~\ref{tab:cofactor}), and form the
cofactor kernel vectors: the free entry is this determinant, the pivot entries are the
negatives of the determinants with the corresponding pivot column replaced, the other
free entries are zero, and the common power of $\epsilon$ is divided out.

\begin{table}[H]
\caption{Least-valuation pivot minors in the positive origin rows (columns
numbered from $0$ in the order $z'_0,z'_1,z'_3,z'_7,\bar c$).}\label{tab:cofactor}
\begin{tabular}{lcccccccccc}
\toprule
$h$&2&3&4&5&6&7&8&9&10&14\\
\midrule
columns&012&012&023&023&013&013&023&023&023&023\\
valuation&2&1&3&2&3&1&3&2&3&5\\
\bottomrule
\end{tabular}
\end{table}

Every entry of the three origin rows has index at most $12$, so a missing cyclic alias
begins at order at least $\mathsf N-1$; for the $\bar{\mathcal B}$ terms, $u-v=j-h$
implies $h-v=j-u\le7$. Replacing the formal entries by the actual cyclic entries changes
the cofactor vectors only by $O(\epsilon^{\mathsf N-6})$, which is $O(\epsilon^8)$, and
the chosen minors remain nonzero, so the corrected vectors lie in the actual origin
kernel. After multiplication by $\epsilon^{\mathsf N}$ the moment rows are regular, and
through order six their values on these vectors are given by the finite sums
\begin{equation}\label{eq:finitesum}
  \sum_{j=8}^{h+7}\psi_A(j)\,
  \frac{\sum_ez'_e\mathsf L_{j-h+e}-\bar c\sum_{u-v=j-h}a_ua_v\mathsf L_{h-v}}
  {D_j(\epsilon;\mathsf N,\mathsf s)},
\end{equation}
where $D_j$ is the series $\epsilon^{-\mathsf N}l_j$ computed from $d_{j,\delta}$:
indices $j\le7$ contribute only from order $\mathsf N-1$ on, indices $j>h+7$ have
numerators without ordinary Laurent coefficient, cyclic corrections to the
numerators start at order $\mathsf N-2$, a second alias in a denominator is at least
$\mathsf N-1$ orders later than the first, and the degree bound on $p_A$ excludes any
$\psi$ alias. All these bounds exceed six for $n\ge89$.

For the quadratic $h\in\{2,4,6,8,10,14\}$ take the unique cofactor vector and its
$P_{11}$ moment, multiplied by $\epsilon^{\mathsf N}$; its first nonzero coefficient has
order at most three. For $h\in\{3,5,7,9\}$ take the first two free cofactor vectors and
the $2\times2$ determinant of the rows $\epsilon^{\mathsf N}\mathsf W_A$ on them, with
$A\in\{11,13\}$ for $h=3,7,9$ and $A\in\{11,33\}$ for $h=5$; its first nonzero order is at
most six. All $36+24$ orders and leading coefficients, for every residue $\mathsf s$, are
listed in Tables~\ref{tab:quadall} and~\ref{tab:twoall} of Appendix~\ref{app:boundary};
for instance, for $\mathsf s=0$ and $h=2$ the coefficient is
$2/(\mathsf N(\mathsf N-1))$ at order three. In each case all lower coefficients vanish.
Each listed coefficient is a rational function of $\mathsf N$ which does not vanish for
real $\mathsf N\ge14$, by the sign test of Appendix~\ref{app:boundary}. In the translation
characters $h=3,7$ the third free vector is the global direction and has zero moments.
Since the positive origin rows have rank three over $\QQ((\epsilon))$ and $\eta$ is
determined, this gives rank six.

For $h=1$ the pivot columns are $012$ with valuation two; the first two cofactor vectors
$K_0,K_1$ correspond to the free $z'_7$ and $\bar c$, the third is global. Only
$j=8$ produces a moment pole of size $\epsilon^{-\mathsf N}$, with numerator
$C_8(z,\bar c)=z'_0\mathsf L_7-\bar c\epsilon\mathsf L_1$. With $C_i=C_8(K_i)$ the vector
$K_*=(C_1K_0-C_0K_1)/\epsilon^2$ is regular, lies in the formal origin kernel, kills
$C_8$, and begins with $(z'_0,z'_1,z'_3,z'_7,\bar c,\kappa)=(-1,2,0,0,1,0)$. The
normalized moment $\epsilon^{\mathsf N}\mathsf W_{11}(K_0)$ has a nonzero leading
coefficient of order at most three in every residue class (Table~\ref{tab:minusone}). For the second row use the exact cyclic projection
$\mathsf W_*=\mathsf W_{13}-[\psi_{13}(8)/\psi_{11}(8)]\mathsf W_{11}$, where
$\psi_{11}(8)$ is a unit with constant term $-1$; it removes the whole summand $j=8$ on
the full source. On $K_*$ the first nonzero coefficient of $\mathsf W_*$ has order zero,
except for $\mathsf s=3$, where it has order one. Both leading coefficients are listed in
Table~\ref{tab:minusone} for every residue $\mathsf s$, and they are nonzero by the same
sign test. The cyclic origin cofactors differ from the formal
ones only at order at least $\mathsf N-4$; solving the first moment pivot on $K_0$, of
order at most three, changes the vector by $O(\epsilon^{\mathsf N-7})$, and since the
projected row has valuation at least $-2$, its displayed coefficient is unchanged
because $\mathsf N-9\ge5$. Thus the two moments are independent modulo the global
direction.

For $\chi=0$ subtract the global direction to make $z'_0=0$, keep the columns
$(z'_1,z'_3,z'_7,\bar c,\rho')$ and only the two positive origin rows $j=1,3$, whose
least pivot is in columns $(0,3)$ with valuation one. The cofactor construction gives
$K_0,K_1,K_2$ (free columns $z'_3,z'_7,\rho'$); put $K'_2=K_2-\epsilon^2K_1$ and use the
moment rows $\mathsf W_{11}$, $\mathsf W_{13}-\mathsf W_{11}$,
$\mathsf W_{33}-2\mathsf W_{13}+\mathsf W_{11}$, whose numerators after division by
$1-u$ have degrees at most $13,9,5$. Multiply $(K_0,K_1,K'_2)$ by
$(\epsilon,\epsilon^2,\epsilon^{-1})$ and the three rows by $\epsilon^{-r_i}$ with
$(r_0,r_1,r_2)=(0,1,1)$ for $\mathsf s=0,1$ and $(0,0,1)$ for $\mathsf s\ge2$. All entries
become regular. The resulting $3\times3$ matrices $\mathsf M$ of constant terms and their
determinants are listed in Appendix~\ref{app:boundary}; for example, for $\mathsf s=2$,
$\det\mathsf M=15/(\mathsf N(\mathsf N+1)(\mathsf N+2)(\mathsf N+3)(\mathsf N+4))$. All six
determinants are nonzero for $\mathsf N\ge14$. Thus two origin rows and three
moments have rank five after the global gauge, and the $\eta$ row gives rank six; the
third positive origin row is not needed.

\emph{The cubic line.} For the remaining ten characters use $f_b$ of the bulk proof at
$b=0$, with $K=n-\chi$,
$D=R_1(1/2)^{\chi-1}/(1-2^{-n})$, $T$ and $d_k$ as there, so that $l_{K+k}=Dd_k$ for
indices in $[4,n-1]$, and $z=\mathsf Af_0+\mathsf BZ_0$ for $z$ divisible by $g_0$, with
$r_0z(u^{-1})=\mathsf Af_0+\mathsf Bl_Z$. The two old-pole moment equations and their
kernel are \eqref{eq:DAC}--\eqref{eq:limkernel}, with $\bar cD$ in place of $cD_1$.
Two facts about wrapped indices are needed: the Taylor coefficients of $r_0$ at $0$
vanish in degrees $0,1,2$, so the old-pole formula for $l_{K+k}$ continues across
$n,n+1,n+2$; and those of $P_A/f_0$ vanish in degrees $0,1$ and equal $(1,0,0)$ in degree
two, so its old-pole formula continues across $n,n+1$ and acquires this extra vector at
$n+2$. The two old-pole residue vectors together with either $(1,0,0)$ or $(1,1,1)$ are
independent. For $1\le\chi\le6$ the true value of $T$ lies in the interval
$T_\chi(1\pm10^{-8})$, $T_\chi=(1792/5103)(-2/3)^{\chi-1}$, because
$(2^{-29}+3^{-29})/(1-3^{-29})<10^{-8}$.

For $\chi=1,2$, the origin rows $j=3,7$ force divisibility by $g_0$, and all moment
contributions lie in the two old-pole directions, except that for $\chi=1$ the wrapped
index $n+2$ contributes the numerator of the remaining positive origin row
$(\mathsf A+798\mathsf B-\bar cl_K)/6$ times the independent vector $(1,0,0)$. In both cases
zero moments force that row and both old-pole equations, which is impossible unless
$9072+\mathsf c(T)(1+T)=0$; this expression and \eqref{eq:DAC} are nonzero on the
corresponding $T$ intervals, by the interval bounds of Appendix~\ref{app:boundary}.

For $\chi=3,4,5,6$ the three positive origin rows at $b=0$ involve only the old-pole
directions and impose divisibility by $g_0$, and the extended matrix has rank five with
kernel \eqref{eq:limkernel}. The left annihilator of the two old positive-origin columns
on the rows $j=1,3,7$ is $(468,-133,1)$. Moving $z$ along the family divisible by the
analytic factor $G_b$, the first $b$-derivatives of $r_bz(u^{-1})$ in the degrees
$4,5,6,7$ are $-504,-546,-468,9072$, and the $c$ contributions to these degrees are
$\bar ca_3l_{K-3}$, $0$, $\bar ca_1l_{K-1}$, $\bar cl_K$. The projected first-order
obstructions are
\begin{align*}
  \chi=3&:\ -173628+\mathsf c(T)\bigl[78d_{-3}+\tfrac{931}6d_{-1}\bigr], &
  \chi=4&:\ -1462104-133\,\mathsf c(T)d_0,\\
  \chi=5&:\ 468\bigl[-468-\tfrac76\mathsf c(T)d_{-1}\bigr], &
  \chi=6&:\ 468\bigl[9072+\mathsf c(T)d_0\bigr],
\end{align*}
each nonzero on its $T$ interval (Appendix~\ref{app:boundary}). Since the projected row
annihilates every column at
$b=0$, these are the first Schur obstructions after the five pivots, independently of
the derivative of the lifted kernel vector.

For $h\in\{11,12,13,15\}$ ($\chi=n-h\ge74$, so $|T|<10^{-4}$), the first two positive
origin rows force divisibility by $g_0$. The polynomial parts of the three moment
functions agree in the degrees $8,\dots,11$ (coefficients $-36,42,0,6$), so all
polynomial contributions lie in the direction $(1,1,1)$, independent of the old-pole
directions, and a zero-moment vector satisfies the two equations of
\eqref{eq:limkernel} and a remaining scalar equation. For $h=11$ the third positive
origin row is $6\mathsf B=0$; for $h=15$ the scalar is $36\mathsf B/(Dd_{-4})$; for
$h=12$, after multiplication by $D/\mathsf B$, it is
$-216/d_{-4}-294/d_{-3}-504/d_{-1}<0$ (the $c$ terms $\mp7\bar c$ cancel); for $h=13$ it is
$F_{13}(T)=252/d_{-4}+252/d_{-2}+\frac76\mathsf c(T)d_{-3}/d_{-2}$, and $TF_{13}(T)$ has
numerator and denominator bounded away from zero on $|T|\le10^{-4}$
(Appendix~\ref{app:boundary}). In each case
$\mathsf B=0$, and then the invertible old system forces $\mathsf A=\bar c=0$. This
completes the proof of Proposition~\ref{prop:boundary}.

\subsection{The lower boundary}
\begin{theorem}\label{thm:lowerboundary}
Let $d\ge27$ and $s=d+3$. Then a general configuration $P$ of $r=\binom{d+2}3+d+3$
points satisfies $\rank J_P=3r-4$, and $T^1_j(C(Z))=0$ for all $j<0$.
\end{theorem}

\begin{proof}
Let $n=d+2\ge29$. Consider first $n\ge89$. Fix coefficients in the intersection of
the finitely many nonempty open subsets of $\Sigma$ given by Lemma~\ref{lem:denom},
Theorem~\ref{thm:special} and Propositions~\ref{prop:detector}, \ref{prop:bulk} and
\ref{prop:boundary}. Since $\mathcal O_\chi$ is injective, \eqref{eq:openingbound} gives
opening ranks $2$, $1$ and $2$ on the special kernels of the ordinary, quadratic and
global characters, so the total opening rank is
\[
  2(n-10)+6+2\cdot4=2n-6 .
\]
For $29\le n\le88$ the same ranks are certified directly by
$\rank[\mathcal P_\chi;\mathcal O_\chi\mathcal W_\chi]=6$ for every character (Appendix~\ref{app:finite}); there the
injectivity of $\mathcal O_\chi$ is not used. By \eqref{eq:residual},
Theorem~\ref{thm:special} and Lemma~\ref{lem:perturb}, for small $\varepsilon\ne0$,
\[
  \rank\Gamma^{\mathrm f}(y(\varepsilon))\ge 3L_d+(n^2+5n+4)-(2n-2)+(2n-6)=N_{n+1}-4 .
\]
Let $Z\subset Y_{n+1}$ be the affine subspace in which the last auxiliary point is the
origin. It is irreducible, contains $y(\varepsilon)$, and meets $Y^\circ_{n+1}$, since
the evaluation at the origin is the nonzero vector $(1,\dots,1)$ and can be completed by
$n$ general points. Lemma~\ref{lem:transfer}(a) gives $\rank J_P=3r-4$ for a general
configuration; explicitly, after eliminating the injective $(n+1)^2$ coefficients
$\beta$,
\[
  \rank J_P=3\tbinom n3+n^2+5n-(n+1)^2=3\tbinom n3+3n-1=3r-4 .
\]
Then $T^1_{-1}=0$ by Lemma~\ref{lem:Wker}. Since $m=M_d-d-3\ge4$,
Propositions~\ref{prop:Psi} and~\ref{prop:lowerneg} give $T^1_j=0$ for all $j<0$ on
the intersection of the two nonempty open sets.
\end{proof}

% =====================================================================
\section{Coupled Horace transitions}\label{sec:horace}

\subsection{Setting}\label{sec:horacesetting}
Let $e\ge2$, $d=e+1$ and $D=\binom{e+2}2=M_e$. Choose integers $0\le s'<D$ and
$0\le c\le e+2$, and put
\[
  a=\tbinom{e+2}3+s',\quad b=D+c,\quad k=D-s',\quad t=e+2-c,\quad s=s'+c,\quad m=k+t .
\]
Then $a+b=L_d+s$ and $m=M_d-s$. Let $H=\{z=0\}$, let $A$ be $a$ distinct points of
$\A^3\smallsetminus H$ and $B$ be $b$ distinct points of $H$, and assume:
\begin{enumerate}[label=(H\arabic*)]
\item evaluation of $F_{e-1}$ on $A$ is injective and evaluation of $F_e$ on $A$ is
surjective;
\item evaluation of $F_e(H)$ (plane polynomials of degree at most $e$) on $B$ is
injective and evaluation of $F_{e+1}(H)$ on $B$ is surjective.
\end{enumerate}
These are nonempty open conditions. Put $E_A=K(A)\subset F_e$, of dimension $k$, and
$V_A=V_{e-1}(A)$. Restriction of $E_A$ to $H$ is injective, since an element divisible
by $z$ would give an equation of degree $e-1$ on $A$. Let
$L_A=E_A|_H$, $L_B=\ev_B(L_A)$, $W=\ev_B F_e(H)$, and
\[
  Q_A=\CC^A/V_A\ (\dim s'),\qquad Q_B=\CC^B/L_B\ (\dim s),\qquad C_B=\CC^B/W\ (\dim c).
\]
The union $\Gamma=A\cup B$ has no equation of degree $e$ (restrict to $B$, divide by
$z$, and use (H1)), evaluation of $F_{e+1}$ on $\Gamma$ is surjective (interpolate on
$B$, then correct on $A$ by $z$ times a polynomial of degree $e$), and restriction to
$H$ gives an exact sequence
\begin{equation}\label{eq:horaceseq}
  0\to zE_A\to K(\Gamma)\to I_{B,H}(\le e+1)\to0,
\end{equation}
where $I_{B,H}(\le e+1)$ is the $t$-dimensional space of plane polynomials of degree at
most $e+1$ vanishing on $B$. Thus $\Gamma\in U^{\mathrm{ev}}_{d,s}$.

\subsection{The coupled identity}
A velocity of $\Gamma$ is a triple $(v_A,u,w)$ with $v_A\in(\CC^3)^A$, tangential
components $u\in(\CC^2)^B$ and normal components $w\in\CC^B$. Define the \emph{normal
constraint}
\[
  R_1(w)(h)=[w\,h(B)]\in C_B\quad(h\in E_A),\qquad W_1=\ker R_1,\qquad\rho_1=\rank R_1 ,
\]
where $w\,h(B)$ is the coordinatewise product. For $w\in W_1$ let
$g_{w,h}\in F_e(H)$ be the unique plane polynomial with values $w\,h(B)$, regarded as a
polynomial independent of $z$, and put
\[
  \mathcal K(w)(h)=\bigl[g_{w,h}(A)/z(A)\bigr]\in Q_A,\qquad
  J_A(v_A)(h)=[D_{v_A}h(A)]\in Q_A .
\]
Thus $J_A$ is the moving differential of $A$ in degree $e$; let $j_A$ be its rank.

\begin{lemma}\label{lem:horacenormal}
The linearized conditions of all equations $zh$, $h\in E_A$, hold if and only if
$w\in W_1$ and $J_A(v_A)=\mathcal K(w)$. Constant vectors lie in $W_1$, and
$\mathcal K(\lambda\mathbf1)=0$.
\end{lemma}

\begin{proof}
The derivative data of $zh$ are $z(A)D_{v_A}h(A)$ on $A$ and $w\,h(B)$ on $B$. A
correcting polynomial $G\in F_e$ must restrict on $H$ to $g_{w,h}$, which requires
$w\in W_1$, so $G=g_{w,h}+zq$ with $q\in F_{e-1}$; comparing values on $A$ gives the
condition. For $w=\lambda\mathbf1$ write $h=h_0+zh_1$ with $h_0=h|_H$; then
$g=\lambda h_0$ and $h_0(A)/z(A)=-h_1(A)\in V_A$.
\end{proof}

Let $\pi$ be the quotient by $\operatorname{im}J_A$, $\tau=\rank(\pi\mathcal K)$ and
$\mathcal Z=\{(v_A,w):w\in W_1,\ J_A(v_A)=\mathcal K(w)\}$. Projection to $w$ gives an
exact sequence $0\to\ker J_A\to\mathcal Z\to\ker(\pi\mathcal K)\to0$, so
$\dim\mathcal Z=(3a-j_A)+(b-\rho_1-\tau)$.

For $f\in I_{B,H}(\le e+1)$ choose a lift $F_f=f+zH_f\in K(\Gamma)$, and a right inverse
$T_A$ of evaluation of $F_e$ on $A$. Define on $\mathcal Z\oplus(\CC^2)^B$
\[
  \mathcal T(v_A,w,u)(f)=\bigl[D_uf(B)+w\,H_f(B)-\ev_BT_A(D_{v_A}F_f(A))\bigr]\in Q_B .
\]
Any correcting polynomial of degree $e$ for $F_f$ must interpolate the derivative data
on $A$, and two such differ by $E_A$, whose values on $B$ form $L_B$. Hence the
remaining linearized conditions are exactly $\mathcal T=0$. Changing $T_A$ or the lift
changes $\mathcal T$ by terms that vanish on $\mathcal Z$, so $\rho_2=\rank\mathcal T$ is
well defined. By \eqref{eq:horaceseq},
\begin{equation}\label{eq:coupled}
  \rank J_\Gamma=j_A+\rho_1+\tau+\rho_2 .
\end{equation}
This is an identity at one configuration; it involves no independence assumption on
its summands.

\begin{proposition}\label{prop:separated}
Suppose that $\ker J_A$ consists of the four motions, that $W_1=\CC\mathbf1$, and that
the operator
\begin{equation}\label{eq:R2op}
  (u,\eta)\longmapsto\bigl(f\mapsto[D_uf(B)+\eta\,N_A(f)(B)]\in Q_B\bigr),
\end{equation}
where $N_A(f)=H_f|_H\bmod L_A$,
on $(\CC^2)^B\oplus\CC$ is injective. Then $\rank J_\Gamma=3(a+b)-4$.
\end{proposition}

\begin{proof}
By Lemma~\ref{lem:horacenormal}, $\mathcal K=0$ on $W_1$, so $\tau=0$, $\rho_1=b-1$, and
$\mathcal Z=\ker J_A\oplus\CC\mathbf1$ has dimension five. Together with $u$ this source
has dimension $2b+5$, and contains the four motions of $\Gamma$. Subtracting the normal
translation we may take $w=0$; subtracting the two tangential translations and the
dilation, $v_A=\lambda e_z$. On this $(2b+1)$-dimensional representative,
$D_{v_A}F_f=\lambda\partial_zF_f$, which has degree at most $e$ and interpolates its own
values on $A$, and $\partial_zF_f|_H=H_f|_H$. Hence $\mathcal T$ becomes
\eqref{eq:R2op} with $\eta=-\lambda$. By \eqref{eq:coupled},
$\rank J_\Gamma=(3a-4)+(b-1)+(2b+1)$.
\end{proof}

\begin{remark}\label{rem:capacity}
Since $\tau\le b-\rho_1-1$, $\tau\le s'k-j_A$, $\rho_1\le ck$ and $\rho_2\le st$, rank
$3(a+b)-4$ requires $st\ge2b+1+(3a-4-j_A)$; in particular a positive residual defect
$3a-4-j_A$ cannot be ignored. For $c=1$ and $s'>0$ we have $\dim W_1\ge s'+1>1$, so the
separated form of Proposition~\ref{prop:separated} is impossible; we use only $c\ge2$.
\end{remark}

The conditions of Proposition~\ref{prop:separated} other than the one on $\ker J_A$ will
be called (R1) ($W_1=\CC\mathbf1$) and (R2) (injectivity of \eqref{eq:R2op}). On the
evaluation pivot charts both are rank conditions on matrices of rational functions of
$(A,B)$, hence open in the irreducible space of pairs $(A,B)$.

\begin{corollary}\label{cor:transition}
Fix $(e,s',c)$. Suppose that a general configuration in degree $e$ and position $s'$
has $T^1_{-1}=0$, and that (R1) and (R2) hold for one pair $(A,B)$ satisfying (H1),
(H2), with $A$ arbitrary. Then a general configuration in degree $e+1$ and position
$s'+c$ has $T^1_{-1}=0$.
\end{corollary}

\begin{proof}
The three open conditions on pairs $(A,B)$ (residual rank, (R1), (R2)) are nonempty
and meet. At a common point Proposition~\ref{prop:separated} gives
$\rank J_\Gamma=3(a+b)-4$, and this is an open condition on all configurations of
$a+b$ points.
\end{proof}

\subsection{Cartesian residuals and the normal operator}
For a lower set $\mathsf T\subset\{(i,j):i+j\le e\}$ with $|\mathsf T|=k$ let
$\Lambda_{\mathsf T}$ consist of all triples of degree at most $e-1$ and the triples
$(i,j,e-i-j)$ with $(i,j)\notin\mathsf T$. Choose pairwise distinct \emph{plane nodes}
$\xi_0,\dots,\xi_{e+1}$ and pairwise distinct nonzero \emph{normal nodes}
$\zeta_0,\dots,\zeta_e$, and let
$A_{\xi,\zeta}=\{(\xi_i,\xi_j,\zeta_l):(i,j,l)\in\Lambda_{\mathsf T}\}$. Put
$N_i(x)=\prod_{v<i}(x-\xi_v)$ and $Z_l(z)=\prod_{v<l}(z-\zeta_v)$. Unless stated otherwise
$\xi_v=v$; we write $A_\zeta$ in that case.

\begin{lemma}\label{lem:cartesian}
For all nodes as above, $A=A_{\xi,\zeta}$ satisfies (H1), $|A|=a$, and
$L_A=\operatorname{span}\{x^iy^j:(i,j)\in\mathsf T\}$. For a plane polynomial
$f=\sum f_{ij}N_i(x)N_j(y)$ of degree at most $e+1$, written in the Newton basis of the
plane nodes, with lift $f+zH_f$ vanishing on $A$,
\begin{equation}\label{eq:Nzeta}
  N_A(f)=-\sum_{i+j\le e,\ (i,j)\notin\mathsf T}w_{e-i-j+1}f_{ij}N_i(x)N_j(y)\bmod L_A,
\end{equation}
where $w_v=\sum_{l<v}\zeta_l^{-1}$.
\end{lemma}

\begin{proof}
The polynomials $N_iN_jZ_l$, $(i,j,l)\in\Lambda_{\mathsf T}$, form a triangular
interpolation basis on $A$, as in Lemma~\ref{lem:newton}, since the nodes on each axis are
distinct. For $(i,j)\in\mathsf T$ the polynomial $N_iN_jZ_{e-i-j}$ vanishes on $A$; these
$k$ polynomials span $E_A$, and their restrictions to $H$ are nonzero multiples of
$N_iN_j$, which span the monomial space of the lower set $\mathsf T$ because the change
from Newton to monomial basis is triangular. Interpolate $-f(x,y)/z$ on $A$ by the Newton
basis. Its coefficients are products of divided differences; the divided difference of
$f$ over the nodes $\xi_0,\dots,\xi_i$ and $\xi_0,\dots,\xi_j$ is $f_{ij}$, and
$[\zeta_0,\dots,\zeta_l](1/z)=(-1)^l/(\zeta_0\cdots\zeta_l)$ by induction from the
recurrence. Since $Z_l(0)=(-1)^l\zeta_0\cdots\zeta_{l-1}$, the restriction to $H$
receives $-f_{ij}/\zeta_l$ from each $(i,j,l)\in\Lambda_{\mathsf T}$. For $(i,j)\notin
\mathsf T$ all $l\le e-i-j$ occur; for $(i,j)\in\mathsf T$ the contribution lies in
$L_A$; and $i+j=e+1$ gives nothing.
\end{proof}

Formula \eqref{eq:Nzeta} is linear in $(w_1,\dots,w_{e+1})$, and the map from
$(\zeta_0^{-1},\dots,\zeta_e^{-1})$ to $w$ is triangular and invertible. Hence, if a
maximal minor of the (R2) matrix with the normal column \eqref{eq:Nzeta} is nonzero at
some \emph{formal} weight vector $w$, it is a nonzero linear polynomial in $w$ and is
nonzero for some pairwise distinct nonzero nodes. This is how formal normal operators
are used below.

Below, $\mathsf T$ is often a \emph{first-row} lower set: the first $k$ pairs
$(i,j)$ in the lexicographic order, that is, ordered by increasing $i$ and then by
increasing $j$.

\subsection{Plane points and canonical modules}\label{sec:R1}
In this and the following subsections $S=\CC[X_0,X,Y]$ denotes the homogeneous
coordinate ring of the projective closure of the plane $H$, with affine coordinates
$x=X/X_0$, $y=Y/X_0$ (the closure of $H$ in $\PP^3$ is the plane $X_3=0$, with homogeneous coordinates $X_0$, $X=X_1$, $Y=X_2$). Let $B$ be
$b=D+c$ general points of $H=\{X_0=1\}$, with homogeneous coordinate ring $R_B$; its $h$-vector is $(1,2,\dots,e+1,c)$.
By Sauer \cite[Proposition~1.7]{Sauer}, for points in a dense open set the ideal $I_B$
has $\max(t,c+1)$ minimal generators, where $t=e+2-c=\dim(I_B)_{e+1}$. Hence:
\begin{itemize}
\item if $c\le(e+1)/2$ (\emph{low} $c$), with $l=e+1-2c\ge0$, the minimal resolution is
$0\to S(-e-2)^l\oplus S(-e-3)^c\to S(-e-1)^t\to I_B\to0$;
\item if $c\ge(e+1)/2$ (\emph{high} $c$), with $q=2c-e-1\ge0$, it is
$0\to S(-e-3)^c\to S(-e-1)^t\oplus S(-e-2)^q\to I_B\to0$.
\end{itemize}
In both cases the graded ranks are forced by the number of generators and the Hilbert
series. Let $\omega$ be the canonical module of $R_B$, graded so that $dt_i$ has degree
one on the normalization. By the residue description \cite[Section~1.3.2]{Stevens},
\[
  \omega_j=(\ev_BS_{-j})^\perp\subset\CC^b ,
\]
multiplication by a form being coordinatewise multiplication by its values; thus
$\omega_j=0$ for $j<-e$, $\omega_{-e}=W^\perp$ has dimension $c$, and
$\omega_0=\mathbf1^\perp$. For a subspace $\mathcal U$ of polynomials put
$\mathcal U_B=\ev_B\mathcal U$. Since $\langle w\cdot l,\theta\rangle=\langle w,l\cdot\theta
\rangle$, we have
\begin{equation}\label{eq:R1dual}
  \{w:w\,L_B\subset W\}=(W^\perp\cdot L_B)^\perp ,
\end{equation}
and $W^\perp\cdot L_B\subset\mathbf1^\perp$ because $L_B\subset W$. So (R1) holds as soon
as $L_A\cdot\omega_{-e}=\omega_0$.

\begin{lemma}\label{lem:presentation}
Let $c\le(e+1)/2$, and let $\mathcal P$ be the affine space of $t\times(t-1)$ matrices
$H=[\mathsf A\mid\mathsf B_q]$ over $S$ whose first $l$ columns consist of linear forms
and whose last $c$ columns consist of quadratic forms. Let $\mathcal P_{\mathrm{red}}$ be
the set of $H$ whose maximal minors generate the ideal of $b$ distinct reduced points of
the affine plane $X_0\ne0$.
\begin{enumerate}[label=\textup{(\alph*)}]
\item $\mathcal P_{\mathrm{red}}$ is a nonempty open subset of $\mathcal P$, and the map
from $\mathcal P_{\mathrm{red}}$ to the space of $b$-point subsets of the affine plane is
dominant.
\item Let $\mathcal U\subset\mathcal P$ be a nonempty open subset defined by the
nonvanishing of minors in the entries of $H$, and suppose that at every point of
$\mathcal U\cap\mathcal P_{\mathrm{red}}$ a given rank condition on point evaluations of $B$
holds. Then this rank condition holds for a general set $B$ of $b$ points.
\end{enumerate}
\end{lemma}

\begin{proof}
(a) The space $\mathcal P$ is affine, hence irreducible. All entries have positive degree,
so at every point of $\PP^2$ the values of the entries can be prescribed independently, and
the universal evaluation map is a submersion onto the space of $t\times(t-1)$ matrices. In
that space the matrices of rank $t-2$ form a smooth stratum of codimension two and those of
rank at most $t-3$ have codimension six. Hence, for general $H$, the degeneracy locus
$\{p:\rank H(p)\le t-2\}$ is a finite reduced subscheme of $\PP^2$ which does not meet the
line $X_0=0$. Its ideal of maximal minors then has height two, the Hilbert--Burch theorem gives
the resolution of Section~\ref{sec:R1} with fixed graded ranks, and therefore the locus
consists of $b=D+c$ reduced points. These are open conditions, so $\mathcal P_{\mathrm{red}}$
is a nonempty open subset. By Sauer's count, the ideal of $b$ general points has exactly $t$
minimal generators, all of degree $e+1$, and the low-$c$ resolution; so it is generated by
the maximal minors of some $H\in\mathcal P$, which lies in $\mathcal P_{\mathrm{red}}$. This
proves dominance; for ordered configurations one passes to the finite \'etale cover given by
the orderings.

(b) Since $\mathcal P$ is irreducible, $\mathcal U\cap\mathcal P_{\mathrm{red}}$ is a
nonempty open subset of $\mathcal P_{\mathrm{red}}$. Its image under the dominant map of (a)
is constructible and dense, hence contains a nonempty open set of configurations, and at
every point of the image the rank condition holds.
\end{proof}

\begin{lemma}[high $c$]\label{lem:R1high}
Let $c\ge(e+1)/2$ and let $L_A$ contain all monomials $y^j$ ($j\le e$) and $xy^j$
($j\le e-1$). Then (R1) holds for general $B$.
\end{lemma}

\begin{proof}
Every element of $\omega_{-e}$ is a minimal generator, and by Sauer's count the type of
$R_B$ is $c=\dim\omega_{-e}$, so $\omega$ is generated in degree $-e$. Dualizing the
resolution presents $\omega$ by $c$ generators of degree $-e$, $t$ quadratic and $q$
linear relations. Choose coordinates with $P_\infty=(0:1:0)\notin B$. The maximal minors
of the Hilbert--Burch matrix generate $I_B$ and do not all vanish at $P_\infty$, so the
presentation matrix has rank $c$ there. Multiply the linear relations by $X$ (which is
$1$ at $P_\infty$) and choose $c$ relations independent at $P_\infty$; constant
combinations give, for a basis $\gamma_1,\dots,\gamma_c$ of $\omega_{-e}$,
\[
  X^2\gamma_i+\sum_j\bigl(Xa_{ij}(X_0,Y)+b_{ij}(X_0,Y)\bigr)\gamma_j=0 .
\]
Repeated substitution reduces every $p\gamma_i$, with $p$ of degree $e$, to a
combination of the $\gamma_j$ with coefficients in
$\CC[X_0,Y]_e\oplus X\CC[X_0,Y]_{e-1}$, which is contained in the homogenization of $L_A$.
Hence $L_A\cdot\omega_{-e}=\omega_0$.
\end{proof}

\begin{lemma}[low $c$]\label{lem:R1low}
Let $e\ge3$, $1\le c\le(e+1)/2$, $l=e+1-2c$ and $q_0=\lceil l/c\rceil+1$. For general $B$
and every integer $\nu>l$, the multiplication map $\mathcal U_{\nu,q_0}\otimes\omega_{-e}
\to\omega_{\nu-e}$ is surjective, where $\mathcal U_{\nu,q_0}$ is spanned by the
monomials $x^iy^j$ with $i\le q_0$ and $i+j\le\nu$. In particular
\[
  \mathcal U_{e,q_0}\cdot W^\perp=\mathbf1^\perp,\qquad
  \mathcal U_{e-2,q_0}\cdot W^\perp=(\ev_BF_2)^\perp\quad(c\ge2),
\]
and (R1) holds whenever $L_A\supset\mathcal U_{e,q_0}$; for a first-row lower set this
means $k\ge(q_0+1)(e+1)-q_0(q_0+1)/2$.
\end{lemma}

\begin{proof}
Put $\mathcal N_\nu=\omega_{\nu-e}$. Dualizing the low-$c$ resolution, $\mathcal N$ is
generated by $c$ elements $g_i$ of degree $0$ and $l$ elements $h_j$ of degree $1$, with
$t$ relations of degree $2$; the presentation matrix has $c$ quadratic and $l$ linear
rows. We first exhibit a presentation with the required property, not necessarily of
reduced points. Partition $l$ into $c$ integers $a_i$ differing by at most one. For
$a_i>0$ take lines $\ell_0,\dots,\ell_{a_i}$ in general position avoiding $P_\infty$, and
let $\mathsf L_i$ be the $a_i\times(a_i+1)$ matrix whose row $j$ has $\ell_{j-1}$ in
column $j-1$ and $-\ell_j$ in column $j$. Its maximal minors are, up to sign, the
products of all but one of the lines, so $\mathsf L_i$ drops rank exactly at the pairwise
intersections of these lines, and only by one; choose the arrangements for different $i$
with disjoint intersection sets. Add a quadratic first row to obtain a square block
$\mathsf P_i$ with $F_i=\det\mathsf P_i$ nonzero at $P_\infty$, the $F_i$ pairwise without
common components (possible since the linear minors have no common curve). Form
$\mathcal B=[\operatorname{diag}(\mathsf P_1,\dots,\mathsf P_c)\mid v]$ with a column $v$,
quadratic in the $g$ rows and linear in the $h$ rows, chosen so that at each of the
finitely many rank-drop points of the $\mathsf L_i$ the linear part of $v$ is outside the
image. Then the $l\times(l+c+1)$ linear submatrix $\mathcal L$ has full row rank at every
point of $\PP^2$, and for generic $v$ the maximal minors of $\mathcal B$ have no common
factor (one of them is $\prod F_i$, and at a general point of a component of $F_i$ some
other minor is nonzero), so they generate an ideal of height two and Hilbert--Burch
applies; the cokernel of $\mathcal B$ is the corresponding shifted canonical module.

Quotient $\mathcal N$ by the $g_i$. The quotient is presented by the linear matrix
$\mathcal L$. For $l>0$ choose $l+2$ general constant combinations of its columns; since
$\mathcal L$ is surjective at every point and the bad choices have codimension three at
each point of the two-dimensional $\PP^2$, the resulting $l\times(l+2)$ matrix is
surjective everywhere, its maximal minors have grade three, and by the Buchsbaum--Rim
criterion \cite[Introduction and Theorem~8.4]{Kustin} its cokernel $Q'$ has the exact
resolution
\[
  0\to S(-l-3)^l\to S(-l-2)^{l+2}\to S(-2)^{l+2}\to S(-1)^l\to Q'\to0 .
\]
Thus $\operatorname{reg}Q'\le l$ and $Q'_\nu=0$ for $\nu>l$; the same holds for the
quotient of $\mathcal N$, which is a quotient of $Q'$. For $l=0$ the quotient is zero.
Hence $\mathcal N_\nu$ is generated by the $g_i$ for $\nu>l$. The adjugate of
$\mathsf P_i$ gives $F_ig_i=0$, and since $F_i(P_\infty)\ne0$ its leading $X$ term is a
nonzero multiple of $X^{a_i+2}$; repeated substitution bounds the $X$ exponent of the
coefficients by $a_i+1\le q_0$, preserving degrees. This proves the surjectivity for the
witness.

The presentation matrix of $\mathcal N$ is the transpose of a Hilbert--Burch matrix,
so the witness is a point of the space $\mathcal P$ of Lemma~\ref{lem:presentation}. On
the height-two locus the graded ranks of the resolutions are constant, so surjectivity in a
fixed degree is the nonvanishing of a minor in the entries of $H$, a nonempty open condition
on $\mathcal P$. At points of $\mathcal P_{\mathrm{red}}$ the residue description turns it
into a rank condition on point evaluations of $B$. By Lemma~\ref{lem:presentation}(b),
surjectivity in the finitely many degrees $l+1,\dots,e+1$ holds for general $B$; for $\nu\ge e+1$, $\omega_{\nu-e}$ has dimension $b$ and multiplication by $X_0$ is an
isomorphism, so the higher degrees follow without raising the $X$ exponent.
\end{proof}

\subsection{The second condition by syzygies}
\begin{lemma}[column count]\label{lem:columns}
Let $\mathsf T$ be a first-row lower set of size $k$ with $2e+1\le k<D$, let $m_0$ be its
last row index and $h_0$ the last $j$ in that row, and let $\Delta_p$ be the set of pairs of
total degree at most $p\ge1$. Then
\[
  \bigl|(\mathsf T+\Delta_p)\cup\Delta_{m_0+1+p}\bigr|=k+p(e+1)+\tfrac{p(p+1)}2+1+[h_0=0]
  \le k+p(e+1)+\tfrac{p(p+1)}2+2 .
\]
\end{lemma}

\begin{proof}
Here $1\le m_0$ and $m_0+1\le e$. The set $\mathsf T$ has full columns $i<m_0$ of height
$e-i$ and a last column $m_0$ of height $h_0$. In $\mathsf T+\Delta_p$ the columns
$i\le m_0+p-1$ have height $e+p-i$ (use a full column and the maximal vertical shift)
and the last column $m_0+p$ has height $h_0$. Adjoining $\Delta_{m_0+1+p}$ raises the
height of that column to $\max(h_0,1)$, adds the column $m_0+p+1$ of height zero, and
changes nothing else since $m_0+1\le e$. Counting gives the formula.
\end{proof}

\begin{proposition}\label{prop:R2syz}
Let $\mathsf T$ be a first-row lower set with $2e+1\le k<D$, $L_A$ its monomial span, and
suppose that the plane trace generators $f_1,\dots,f_t$ of $I_B(\le e+1)$ admit
polynomial syzygies of degree at most $p$ spanning their syzygy space over
$\CC(x,y)$, for general $B$, and that $t\ge2$. If
\begin{equation}\label{eq:R2dim}
  k+p(e+1)+\tfrac{p(p+1)}2+2\le b,
\end{equation}
then (R2) holds for the Cartesian residual $A_\zeta$ of $\mathsf T$, for suitable nodes
$\zeta$ and general $B$.
\end{proposition}

\begin{proof}
Let $\ell'=m_0+1\le e$, take the plane nodes $\xi_v=v$, and use the formal weight vector
with $w_{e-\ell'+1}=-1$ and all other $w_v=0$; then \eqref{eq:Nzeta} is the projection
$P_{\ell'}(f)=\sum_{i+j=\ell'}f_{ij}N_iN_j\bmod L_A$ onto Newton total degree $\ell'$. Suppose
$(D_{u_p}f_i(p)+\eta P_{\ell'}(f_i)(p))_{p\in B}\in\ev_BL_A$ for all $i$. Choose
$g_i\in L_A$ with these values and put $G_i=g_i-\eta P_{\ell'}(f_i)\in L_A+F_{\ell'}$, so
that $G_i(p)=D_{u_p}f_i(p)$. Differentiating a syzygy $\sum a_if_i=0$ at $p\in B$ gives
$\sum a_i(p)G_i(p)=0$, and $\sum a_iG_i$ lies in $(L_A+F_{\ell'})F_p$, whose dimension is
at most $b$ by Lemma~\ref{lem:columns} and \eqref{eq:R2dim}. For general $B$ evaluation is
injective on this fixed space, so $\sum a_iG_i=0$ for all syzygies, and $G=\varrho f$ for a
rational function $\varrho$. The $f_i$ have no common factor: a common factor of degree
$j$, $1\le j\le e+1$, would contain at most $\binom{j+2}2-1$ of the general points, and
dividing it out would place the $t$-dimensional space $I_B(\le e+1)$ in a space of
dimension at most $\max(0,t-j(e+1-j))<t$, or of dimension at most one if $j=e+1$. So
$\varrho$ is a polynomial, and since $\deg G_i\le e<e+1=\deg f_i$, $G=0$.

Let $\lambda$ be the coefficient functional of $N_{\ell'}(x)N_0(y)$. Since
$(\ell',0)\notin\mathsf T$, $P_{\ell'}(f)\in L_A$ implies $\lambda(f)=0$. On the other
hand $\lambda$ does not vanish on $I_B(\le e+1)$ for general $B$: choosing
$\dim F_{e+1}$ points whose evaluations form a basis of the dual, some $b$ of them have
evaluation span not containing $\lambda$. If $\eta\ne0$ then $P_{\ell'}(f_i)=g_i/\eta\in L_A$
for all $i$, contradicting $\lambda\ne0$ on $I_B(\le e+1)$. So $\eta=0$, $g_i=0$ and
$D_{u_p}f_i(p)=0$. Finally, at a general point $p$ the derivatives of
$I_B(\le e+1)$ span the cotangent space: the polynomials of degree at most $e+1$
vanishing to order two at $p$ have codimension three, and $b-1$ further general points
impose independent conditions on them because $t\ge2$. Hence $u=0$. Passing from the
formal weights to actual nodes as explained after Lemma~\ref{lem:cartesian}, (R2) holds
for some $A_\zeta$, and then on a nonempty open set.
\end{proof}

\begin{lemma}\label{lem:syzdeg}
Let $c\ge2$.
\begin{enumerate}[label=\textup{(\alph*)}]
\item If $l=e+1-2c\ge0$, the syzygies of the $f_i$ are spanned by the columns of the
Hilbert--Burch matrix, of degrees one and two; so $p=2$.
\item If $l<0$ and $t=e+2-c\ge3$, put $q=2c-e-1$, $h_0'=\lfloor(t-1)/2\rfloor$ and
$\mathsf h=\lceil q/h_0'\rceil$. Then the syzygies among $f_1,\dots,f_t$ of degree at most
$p=\mathsf h+2$ span their syzygy space over the fraction field.
\end{enumerate}
\end{lemma}

\begin{proof}
(a) is the low-$c$ resolution. In (b) the Hilbert--Burch matrix is a $t\times c$ quadratic
block $U$ over a $q\times c$ linear block $V$, with $t+q=c+1$. The block $V$ has rank $q$:
otherwise its kernel would have dimension at least $t$, $U$ would be injective on it,
and its image would be annihilated by the nonzero row $(f_1,\dots,f_t)$. Syzygies among
the $f_i$ alone are $U\kappa$ with $\kappa\in\ker V$. For general $V$ the kernel of
$V$ is spanned over the fraction field by polynomial vectors of degree $\mathsf h$: for
a witness, partition $q$ into $h_0'$ balanced parts $a_i$ and take a block diagonal matrix
with $a_i\times(a_i+2)$ blocks, everywhere of full rank on $\PP^2$, adding a zero column if
$t-1$ is odd; by the Buchsbaum--Rim complex each block has kernel generated in coefficient
degree $a_i\le\mathsf h$, and empty blocks contribute free vectors. Surjectivity of the
coefficient-degree-$\mathsf h$ kernel map is open, a fraction-field spanning minor
persists on an open set, and this open set meets the plane point presentations.
\end{proof}

\begin{proposition}[low-$c$ and high-$c$ transitions]\label{prop:gatesBC}
Let $e\ge3$ and $c\ge2$.
\begin{enumerate}[label=\textup{(\alph*)}]
\item If $l=e+1-2c\ge0$, put $q_0=\lceil l/c\rceil+1$, $p=2$ and
$k_{\mathrm{req}}=(q_0+1)(e+1)-q_0(q_0+1)/2$.
\item If $l<0$ and $t\ge3$, put $p$ as in Lemma~\ref{lem:syzdeg}(b) and $k_{\mathrm{req}}=2e+1$.
\end{enumerate}
Let $s'$ be a position such that a general configuration in degree $e$ and position
$s'$ has $T^1_{-1}=0$ (so that $s'\ge e+3$), and suppose
\begin{equation}\label{eq:gateBC}
  p(e+1)+\tfrac{p(p+1)}2+2-c\le s'\le D-k_{\mathrm{req}} .
\end{equation}
Then a general configuration in degree $e+1$ and position $s=s'+c$ has $T^1_{-1}=0$.
\end{proposition}

\begin{proof}
Let $\mathsf T$ be the first-row lower set of size $k=D-s'\ge k_{\mathrm{req}}\ge2e+1$. By
Lemmas~\ref{lem:R1low} and~\ref{lem:R1high}, (R1) holds for general $B$. The lower bound
in \eqref{eq:gateBC} is \eqref{eq:R2dim}, so Proposition~\ref{prop:R2syz} and
Lemma~\ref{lem:syzdeg} give (R2); here $t\ge2$ in case (a) because $t=l+c+1$, and
$k<D$ because $s'>0$. Corollary~\ref{cor:transition} applies.
\end{proof}

\subsection{The sparse transition}
\begin{proposition}[sparse transition]\label{prop:gateA}
Let $c\ge2$ and $e\ge2c^2+2c-3$, and put $l=e+1-2c$, $a_0=\lfloor l/(c+1)\rfloor$,
$q_0=\lceil l/c\rceil+1$ and
\[
  K_{\min}=(q_0+1)(e-1)-\tfrac{q_0(q_0+1)}2+2e-2-2q_0 .
\]
If a general configuration in degree $e$ and position $s'$ has $T^1_{-1}=0$, and
$e+c+1\le s'\le D-K_{\min}$, then a general configuration in degree $e+1$ and position
$s'+c$ has $T^1_{-1}=0$.
\end{proposition}

\emph{The residual.} Let $\mathbb D_0=\{a_0,2a_0+1,\dots,(c+1)a_0+c\}$, a subset of
$\{2,\dots,e-1\}$. Let $\mathsf T_{\max}$ consist of all pairs of degree at most $e-1$
except the $c+1$ pairs $(e-1-\delta,\delta)$ with $\delta\in\mathbb D_0$, together with
$(e,0)$; it is a lower set of size $D-e-c-1$. Let $\mathsf T_{\min}$ be the union of the
pairs of degree at most $e-2$ with $i\le q_0$, of the axis $(i,0)$, $i\le e$, and of the
pairs $(i,1)$, $i\le e-2$; it is a lower set of size $K_{\min}$ contained in
$\mathsf T_{\max}$. Every size between $K_{\min}$ and $D-e-c-1$ is realized by a lower set
$\mathsf T$ with $\mathsf T_{\min}\subset\mathsf T\subset\mathsf T_{\max}$; this is
exactly the range $e+c+1\le s'\le D-K_{\min}$. Use the Cartesian residual $A_\zeta$ of
$\mathsf T$.

\emph{The binary data.} Write the low-$c$ Hilbert--Burch matrix as $[\mathsf A\mid
\mathsf B_q]$, with $l$ linear and $c$ quadratic columns, and let a subscript $0$ denote
restriction to the line $X_0=0$. Let $\Lambda_{e+1}\subset\CC[X,Y]_{e+1}$ be the
$t$-dimensional space of leading forms of $I_B(\le e+1)$, with $c$-dimensional quotient
$\mathsf Q$. We use Lemma~\ref{lem:presentation} and the following open conditions on the binary
restriction, each shown to be nonempty by a witness.
\begin{enumerate}[label=(B\arabic*)]
\item \emph{Admissibility}: $\Lambda_{e+1}$ is base-point-free, that is, the forms in
$\Lambda_{e+1}$ have no common zero on $\PP^1$, and
$\CC[X,Y]_1\Lambda_{e+1}=\CC[X,Y]_{e+2}$. A witness is the monomial space whose quotient
positions (exponents of $Y$) are $1,3,\dots,2c-1$.
\item \emph{Detection}: the maps $p_2\mapsto([X^{e-1}p_2],[X^{e-2}Yp_2])\in\mathsf Q^2$
on binary quadrics and $p_1\mapsto[X^ep_1]\in\mathsf Q$ on binary linear forms are
injective. A witness is any $\Lambda_{e+1}$ on which the two coefficient functionals
$c_0+c_2$ and $c_1$ vanish ($c_j$ the coefficient of $X^{e+1-j}Y^j$): for
$p_2=\alpha X^2+\beta XY+\gamma Y^2$ they read $(\alpha+\gamma,\beta)$ on $X^{e-1}p_2$ and
$(\beta,\alpha)$ on $X^{e-2}Yp_2$, and for $p_1=\delta X+\epsilon Y$ they read
$(\delta,\epsilon)$ on $X^ep_1$.
\item \emph{The rank conditions in the proof of (R2) below}: $\mathsf A_0^{\mathsf t}$ has
no constant kernel vector, $\pi_{\mathbb D_0}$ is injective on the kernel of
$\mathsf A_0^{\mathsf t}$ in degree $e-1$, and the matrix $[\pi_{\mathbb D_0}T_0\mid
\pi_{\mathbb D_0}f_2]$ has rank $ct+1$. A witness is the binary chain matrix
$\mathsf A_0$ with $c+1$ blocks of balanced lengths $a_i\ge a_0\ge1$ summing to $l$,
together with the bidiagonal matrix $\mathsf B_0$ and the lift described there.
\end{enumerate}
Conditions (B1) and (B2) are open in the Grassmannian $\Gr(t,\CC[X,Y]_{e+1})$, which is
irreducible, so they hold simultaneously on a nonempty open subset. On the open set (B1)
the Hilbert numerator of the binary ideal generated by $\Lambda_{e+1}$ yields a minimal
binary Hilbert--Burch presentation with $l$ linear and $c$ quadratic columns; placing it in
$\mathcal P$ with entries independent of $X_0$ gives a point at which the minors expressing
(B1) and (B2) do not vanish. For (B3) the chain witness is realized inside $\mathcal P$: the
quadratic entries that pair with the kernel vector $v_i$ of a chain of length $a_i$ span all
binary forms of degree $a_i+2$, so the bidiagonal matrix below is realized by actual
quadratic columns, and the linear entries span the binary forms of degree $a_i+1$, so the
lift $X_0\mathsf B_1$ used below is realized as well. Each of (B1)--(B3) is therefore a
nonempty open condition on $\mathcal P$, and by Lemma~\ref{lem:presentation}, together with
the finitely many open conditions on point evaluations used below, they hold for general
$B$.

\begin{proof}[Proof of Proposition~\ref{prop:gateA}: condition (R1)]
By Lemma~\ref{lem:R1low} and $L_A\supset\mathcal U_{e-2,q_0}$, a vector $w$ with
$wL_B\subset W$ lies in $(\mathcal U_{e-2,q_0}\cdot W^\perp)^\perp=\ev_BF_2$, so $w$ is
the evaluation of a polynomial $p$ of degree at most two. If $g\in L_A$ and $pg$ has
degree $e+1$, then $w\cdot g(B)\in W$ means that $pg$ agrees on $B$ with a polynomial of
degree at most $e$, so the homogeneous part of degree $e+1$ of $pg$ lies in
$\Lambda_{e+1}$. Testing with $g=x^{e-1}$ and
$g=x^{e-2}y$, which lie in $\mathsf T_{\min}$, condition (B2) kills the quadratic part of
$p$, and testing with $g=x^e$ kills its linear part. So $w$ is constant.
\end{proof}

\begin{proof}[Proof of Proposition~\ref{prop:gateA}: condition (R2)]
\emph{Two degenerations.} Take the plane nodes $\xi_v=\epsilon v$ and the formal weights
$w_2=-1$, $w_v=0$ for $v\ne2$. By Lemma~\ref{lem:cartesian}, which holds for every
$\epsilon\ne0$, the normal column \eqref{eq:Nzeta} is then the projection
$P^{(\epsilon)}_{e-1}(f)=\sum_{i+j=e-1}f_{ij}N_iN_j\bmod L_A$ onto Newton total degree
$e-1$ for the scaled Newton basis. As $\epsilon\to0$ the scaled Newton basis and the Newton
coefficients of $f$ tend polynomially to the monomial basis and the monomial
coefficients, while $L_A$ is the same monomial space for all $\epsilon$; hence
$P^{(\epsilon)}_{e-1}$ tends to the projection $P^{\mathrm{hom}}_{e-1}$ onto the homogeneous part
of degree $e-1$, and the (R2) matrix with this normal column is polynomial in $\epsilon$.
It therefore suffices to show that the (R2) operator with normal column $P^{\mathrm{hom}}_{e-1}$ is
injective; then a $(2b+1)$-minor is nonzero for some $\epsilon\ne0$, it is a nonzero
linear polynomial in $w$ for these plane nodes, and it is nonzero for suitable pairwise
distinct nonzero normal nodes $\zeta$, by the remark after Lemma~\ref{lem:cartesian}.

\emph{The homogeneous operator.} Suppose
$D_uf_i(B)+\eta P^{\mathrm{hom}}_{e-1}(f_i)(B)\in\ev_BL_A$, choose $g_i\in L_A$ and put
$G_i=g_i-\eta P^{\mathrm{hom}}_{e-1}(f_i)\in F_{e-1}+\CC x^e$, so $G_i(B)=D_uf_i(B)$. The space
$(F_{e-1}+\CC x^e)F_1$ has dimension $D+2\le b$, so differentiating the linear syzygies
gives $\mathsf A^{\mathsf t}G=0$ identically. The degree $e$ part of $G$ is $\alpha x^e$
with $\alpha\in\CC^t$; since $\mathsf A_0^{\mathsf t}$ has no constant kernel vector
by (B3) (at the chain witness because every chain has positive length), $\alpha=0$ and the homogenization of $G$ (in degree $e$) is
divisible by $X_0$.

The quadratic syzygies give $\mathsf B_q^{\mathsf t}G\in I_B$ in degree $e+2$, so there is
a $t\times c$ matrix $\mathsf C$ of linear forms with
$\mathsf B_q^{\mathsf t}G+\mathsf C^{\mathsf t}f=0$. Let $\delta f(\mathsf C)$ be the
derivative of the maximal minors of $[\mathsf A\mid\mathsf B_q+\epsilon\mathsf C]$ at
$\epsilon=0$; it has degree $e$, and $\mathsf A^{\mathsf t}\delta f=0$,
$\mathsf B_q^{\mathsf t}\delta f+\mathsf C^{\mathsf t}f=0$. The polynomial kernel of
$[\mathsf A\mid\mathsf B_q]^{\mathsf t}$ is generated by the primitive vector $f$ of
degree $e+1$, so $G=\delta f(\mathsf C)$. At $X_0=0$ these derivative minors vanish, so
each column of $\mathsf C_0$ is a linear syzygy of the binary $f_0$, hence equals
$\mathsf A_0$ times a constant vector by the binary Hilbert--Burch theorem; subtracting
this column-operation gauge we get $\mathsf C=X_0\mathsf C^0$ with $\mathsf C^0$ constant.
Thus
\begin{equation}\label{eq:GZ}
  G=X_0\,\delta f(\mathsf C^0),\qquad (G/X_0)|_{X_0=0}=T_0(\mathsf C^0):=\delta f_0(\mathsf C^0),
\end{equation}
and $T_0$ is injective on the $ct$-dimensional space of constant $\mathsf C^0$, since a
kernel vector would give constant syzygies of the minimal binary ideal. Let $f_2$ be the
coefficient of $X_0^2$ in $f$ and $\pi_{\mathbb D_0}$ the vector of coefficients of the
deleted monomials $x^{e-1-\delta}y^\delta$, $\delta\in\mathbb D_0$, in each component.
Comparing the parts of degree $e-1$ of $G=g-\eta P^{\mathrm{hom}}_{e-1}(f)$, and using that $g\in L_A$
has no deleted coefficients, we get
\begin{equation}\label{eq:detect}
  \pi_{\mathbb D_0}\bigl(T_0(\mathsf C^0)+\eta f_2\bigr)=0 .
\end{equation}

\emph{The detector is injective on the binary kernel.} At the chain witness of (B3),
the kernel of $\mathsf A_0^{\mathsf t}$ on $\CC[X,Y]_{e-1}^t$ is $\bigoplus_iv_i\CC[X,Y]_{e-1-a_i}$ with
$v_i=(X^{a_i},X^{a_i-1}Y,\dots,Y^{a_i})$ on the $i$-th block. Deleting the coefficient
of $Y$-degree $\mathsf h$ in all components of one block detects the coefficients of the
cofactor in the degrees $[\mathsf h-a_i,\mathsf h]$. For the minimal length $a_0$ the
$c+1$ intervals attached to $\mathbb D_0$ are consecutive from $0$ to $(c+1)a_0+c$.
Writing $e=(c+1)a_0+2c-1+\varrho$ with $0\le\varrho\le c$, the required last degree is
$e-1-a_0=ca_0+2c-2+\varrho$, which is at most $(c+1)a_0+c$ because $a_0\ge2c-2\ge
c-2+\varrho$; this uses $e\ge2c^2+2c-3$. Longer chains only widen the intervals.

\emph{A normal witness.} Let $w_i=a_i+2$ and take, in the quotient by the chain columns,
the $(c+1)\times c$ bidiagonal quadratic matrix $\mathsf R$ with $\mathsf R_{i,i}=X^{w_i}$
($0\le i<c$) and $\mathsf R_{i,i-1}=Y^{w_i}$ ($1\le i\le c$). Its omitted-row minors are,
up to sign, $F_i=\prod_{h<i}X^{w_h}\prod_{h>i}Y^{w_h}$, and the component of $f_0$ in the
block $i$ is $v_iF_i$; since $F_0$ and $F_c$ are pure powers of $Y$ and $X$, this is a
genuine binary Hilbert--Burch matrix, realized by quadratic columns $\mathsf B_0$. A
constant variation changes an entry in row $u$ and column $v$; deleting row $1$, a nonzero
cofactor has $u\le v<1$ or the reverse, and its ratio to $F_1$ does not decrease the $Y$
exponent (for $u<1$) or the $X$ exponent (for $u>1$). Hence the monomial $F_1/(XY)$ does
not occur in the constant tangent image in the component $1$. Now take the lift
$\mathsf B_q=\mathsf B_0+X_0\mathsf B_1$ whose only nonzero entries of $\mathsf B_1$ replace
$\mathsf R_{0,0}$ by $X^{w_0-1}$ and $\mathsf R_{2,1}$ by $Y^{w_2-1}$, realized by linear
entries of the actual blocks; the coefficient of $X_0^2$ in the component $1$ is exactly
$F_1/(XY)$. Therefore $f_2\notin\operatorname{im}T_0$, and since
$\pi_{\mathbb D_0}$ is injective on the binary kernel, the matrix
$[\pi_{\mathbb D_0}T_0\mid\pi_{\mathbb D_0}f_2]$ has rank $ct+1$. This lift need not define
reduced points; it only certifies a nonzero polynomial minor, which persists on an open
subset of the irreducible presentation space, where it meets the open sets of reduced point
sets, of the low-$c$ resolution and of the injectivity of evaluation on
$(F_{e-1}+\CC x^e)F_1$.

At such a presentation \eqref{eq:detect} forces $\eta=0$ and $\mathsf C^0=0$, then $G=0$
by \eqref{eq:GZ}, $g=0$ and $D_uf_i(B)=0$, and $u=0$ as in Proposition~\ref{prop:R2syz}.

By the first step this gives (R2) for an actual Cartesian residual $A_{\xi,\zeta}$, and
Corollary~\ref{cor:transition} gives the proposition.
\end{proof}

% =====================================================================
\section{A rational seed and the propagation to large degrees}\label{sec:seed}

\subsection{The seed \texorpdfstring{$(d,s)=(11,14)$}{(d,s)=(11,14)}}
\begin{proposition}\label{prop:seed}
For $(d,s)=(11,14)$, $r=300$, a general configuration satisfies
$\rank J_P=896=3r-4=sm$, so $T^1_{-1}=0$.
\end{proposition}

\begin{proof}
Let $n=13$. Appendix~\ref{app:seed} describes thirteen planes
$\ell_i=1+\nu_i\cdot x$ with rational normals, points $p_i$ with $\ell_i(p_i)=0$ and
directions $w_i$ with $\nu_i\cdot w_i=1$; the third coordinates are fixed by
$\nu_{i,3}=(-1-\nu_{i,1}p_{i,1}-\nu_{i,2}p_{i,2})/p_{i,3}$ and
$w_{i,3}=(1-\nu_{i,1}w_{i,1}-\nu_{i,2}w_{i,2})/\nu_{i,3}$, so these identities hold
over $\QQ$. Every three normals are independent and no four planes meet. Let $A$ be the
$286$ triple points and, for $\lambda\in\{1,2\}$, let $B_\lambda$ consist of the
points $p_i+\lambda w_i$ and the origin. One checks over $\QQ$ that each triple point
lies on exactly its three planes, that no point of $B_\lambda$ lies on a plane, and
that the $300$ points are distinct and span $\A^3$.

By Proposition~\ref{prop:starjac}, after choosing a basis $\mathcal K$ of
$\ker\mathcal E$ ($\mathcal E$ the $14\times78$ edge-product evaluation matrix), $J_P$
is the $896\times900$ matrix with rows $(j,\alpha)$ and entries
\begin{gather*}
  J_{(j,\alpha),(T,k)}=-V_T(b_j)\mathcal K_{T\smallsetminus k,\alpha},\\
  J_{(j,\alpha),(j,\beta)}=\sum_{ab}\mathcal K_{ab,\alpha}\partial_\beta E_{ab}(b_j),
  \qquad J_{(j,\alpha),(j',\beta)}=0\ \ (j'\ne j).
\end{gather*}
All coordinates and the pivots used to compute $\mathcal K$ have denominators prime to
$P=32003$. Modulo $P$ one finds $\rank\mathcal E=14$ and $\rank J=896$, with a selected
$896\times896$ minor whose determinant is $19954$ for $\lambda=1$ and $31377$ for
$\lambda=2$ modulo $P$; either configuration alone suffices for the argument, and both
are recorded in the supplementary material (Appendix~\ref{app:data}). The same pivots define a rational basis
$\mathcal K$ reducing to the
modular one, so the corresponding rational minor has a denominator prime to $P$ and a
nonzero reduction; it is therefore nonzero. Hence $\rank J_P=896$ over $\QQ$ and over
$\CC$, and the kernel consists of the four motions. Lemma~\ref{lem:open} gives the
statement for general configurations.
\end{proof}

\subsection{Finite propagation up to degree 84}
For each degree $d$ we write $[\sigma^-_d,\sigma^+_d]$ for an interval of positions $s$ on
which a general configuration has already been shown to satisfy $\rank J_P=3r-4$.
Starting from Proposition~\ref{prop:seed}, we apply the transitions of
Propositions~\ref{prop:gateA} and~\ref{prop:gatesBC} degree by degree. At degree $e$ the
input is an interval $[\sigma^-_e,\sigma^+_e]$ of positions already proved; for each admissible $c$ and
each transition, its residual interval is intersected with $[\sigma^-_e,\sigma^+_e]$ and shifted by $c$.
This is a finite integer computation, and each of its steps is an instance of the cited
propositions, whose hypotheses are inequalities between explicit integers. The resulting
$188$ transitions are listed in Appendix~\ref{app:propagation}. Their lower endpoints are
always $2d-8$, and the last row is
\begin{equation}\label{eq:row84}
  d=84:\qquad 160\le s\le2640,\qquad M_{84}-2640=1015=7\cdot84+427 .
\end{equation}

\subsection{The analytic induction}
\begin{theorem}\label{thm:broad}
For every $d\ge84$ and $2d-8\le s\le M_d-7d-427$, a general configuration satisfies
$\rank J_P=3r-4$.
\end{theorem}

\begin{proof}
The case $d=84$ is \eqref{eq:row84}. Let $e\ge84$ and assume the statement for
$[\sigma^-_e,\sigma^+_e]=[2e-8,D-7e-427]$, $D=M_e$. We cover $[2(e+1)-8,M_{e+1}-7(e+1)-427]=
[2e-6,\sigma^+_e+e-5]$ by three transitions.

\emph{Sparse, $c=2$.} Here $e\ge9$ and $q_0=\lfloor e/2\rfloor$. Writing $K_2$ for
$K_{\min}$, for $e=2h$ we have $K_2=(3h^2+5h-6)/2$ and $D-K_2-(2e-2)=(h-3)(h-4)/2$, and
for $e=2h+1$, $K_2=(3h^2+7h)/2$ and $D-K_2-(2e-2)=(h-2)(h-3)/2$. In either case
$K_2\ge3e^2/8\ge7e+427$, so $D-K_2\le \sigma^+_e$; also $\sigma^-_e\ge e+3$ and $D-K_2\ge \sigma^-_e$.
Proposition~\ref{prop:gateA} gives $I_2=[2e-6,D-K_2+2]$.

\emph{Low $c$, $c_H=\lfloor(e+1)/2\rfloor$.} Then $l\in\{0,1\}$, $p=2$ and
$k_{\mathrm{req}}\le3e$, which is at most $k=D-s'$ for all $s'\le \sigma^+_e$. The lower bound of
\eqref{eq:gateBC} is $2e+7-c_H\le2e-8$ since $c_H\ge15$. Proposition~\ref{prop:gatesBC}
gives $I_H=[2e-8+c_H,\sigma^+_e+c_H]$. There is no gap after $I_2$: with
$\Delta=D-K_2-(2e-2)$ we need $\Delta\ge c_H-9$, and $\Delta-c_H+9$ equals
$(h^2-9h+30)/2$ for $e=2h$ and $(h^2-7h+22)/2$ for $e=2h+1$, both positive.

\emph{High $c$, $c_7=e-5$.} Then $t=7$, $q=e-11$, $h_0'=3$,
$p=\lceil(e-11)/3\rceil+2\le e/3-1$, $k_{\mathrm{req}}=2e+1$, and the image is
$I_7=[T_e,\sigma^+_e+e-5]$ with $T_e=p(e+1)+p(p+1)/2+2\ge3e-13$. Since
$T_e\le7e^2/18-5e/6+1$ and $\sigma^+_e+c_H\ge e^2/2-5e-426$,
\[
  (\sigma^+_e+c_H)-(T_e-1)\ge e^2/9-25e/6-426 ,
\]
which equals $8$ at $e=84$ and increases. So $I_7$ meets $I_H$, and
$I_2\cup I_H\cup I_7\supset[2e-6,\sigma^+_e+e-5]$.
\end{proof}

\section{Proof of the main theorems}\label{sec:proof}

\begin{proposition}\label{prop:bands}
A general configuration satisfies $\rank J_P=3r-4$ in each of the following cases:
\begin{enumerate}[label=\textup{(\alph*)}]
\item $d\ge27$ and $d+3\le s\le2d-24$;
\item $d\ge28$ and $d+3\le s\le3d-52$.
\end{enumerate}
\end{proposition}

\begin{proof}
(a) Put $\mathsf t=s-d-3\ge0$ and $d_0=d-\mathsf t=2d+3-s\ge27$. Start from
$(d_0,d_0+3)$, given by Theorem~\ref{thm:lowerboundary}, and apply the sparse transition
with $c=2$ exactly $\mathsf t$ times; for $\mathsf t=0$ there is nothing to do. Before the
$j$-th step the residual has degree $e=d_0+j-1\ge27$ and position $s'=d_0+2j+1$, so
$s'-(e+3)=j-1\ge0$ and $(2e-2)-s'=d_0-5\ge22$. By the formulas for $D-K_2-(2e-2)$ in the
proof of Theorem~\ref{thm:broad}, $D-K_2\ge2e-2\ge s'$, and $e\ge9$; so the hypotheses
of Proposition~\ref{prop:gateA} hold, and the residual is the position reached at the
previous step. The final position is $s$.

(b) At $d=28$ the interval $[31,32]$ is covered by (a). Suppose (b) holds in degree
$e\ge28$, and apply the sparse transition with $c=3$ to $s'\in[e+4,3e-52]$; here
$e\ge21$ and the lower bound is $e+c+1$. With $q_0=\lceil(e-5)/3\rceil+1\le(e+1)/3$ and
$K_{\min}$ increasing in $q_0$ in this range,
\[
  D-K_{\min}\ge\frac{2e^2-5e+47}9\ge3e-52,
\]
the last difference being $(2X^2+80X+1187)/9$ with $e=X+28$. This gives
$[d+6,3d-52]$ in degree $d=e+1$, and (a) gives $[d+3,2d-24]\supset[d+3,d+5]$ for $d\ge29$.
\end{proof}

\begin{proof}[Proof of Theorems~\ref{thm:main}, \ref{thm:lower} and~\ref{thm:rank}]
The rank assertions are Theorem~\ref{thm:lowerboundary}, Proposition~\ref{prop:bands}
and, for $d\ge84$, the union of Proposition~\ref{prop:bands}(b) with
Theorem~\ref{thm:broad}: the two intervals overlap because $(3d-52)-(2d-9)=d-43>0$, so
together they cover $d+3\le s\le M_d-7d-427$. In each case the open set of
Lemma~\ref{lem:open} is nonempty, and Proposition~\ref{prop:leading} translates the
rank statement into the dimension of the image of $\Phi$. By Lemma~\ref{lem:Wker},
$T^1_{-1}=0$. In all cases $m\ge4$: at $s=d+3$, $m=(d^2+d-4)/2$; in the band
(a), $m\ge(d^2-d+50)/2$; in the band (b), $m\ge(d^2-3d+106)/2$; in the broad range,
$m\ge7d+427$. Propositions~\ref{prop:Psi} and~\ref{prop:lowerneg} give $T^1_j=0$ for all
$j<0$ on the intersection with the open set where $\Psi_P$ is injective, which is
nonempty by irreducibility. The cones are reduced, have an isolated singularity at the
vertex and no equations of degree less than $d\ge27$, so
Proposition~\ref{prop:nonsmooth} excludes algebraic and analytic smoothings. The
statement at $s=d+2$ is Theorem~\ref{thm:outer}.

For the density, every integer $r\ge1$ can be written as $L_d+s$ with $0\le s<M_d$, since
$L_{d+1}-L_d=M_d$. In the layer $d\ge84$ the omitted positions number
$M_d-(M_d-8d-429)=8d+429$. Up to the layer $D$ the omissions are $O(D^2)$ among
$L_{D+1}=\Theta(D^3)$ integers, and a partial last layer below a bound $R$ contributes
$O(D^2)=O(R^{2/3})$ further omissions. Hence the set of branch numbers covered has
natural density one.
\end{proof}

\begin{remark}\label{rem:numbers}
At $d=84$ the covered interval is $87\le s\le2640$, that is $102427\le r\le104980$; the
broad range of Theorem~\ref{thm:broad} alone begins at $s=160$. The first case of
Theorem~\ref{thm:lower}(a) is $d=27$, $r=3684$, with $\rank J_P=11048$. The uniform part of
the lower-boundary proof starts at $d=87$, $r=113654$; the degrees $27\le d\le86$ are
covered by the finite certificates of Appendix~\ref{app:finite}.
\end{remark}

\begin{remark}\label{rem:finiterange}
The same arguments give further finite ranges: Proposition~\ref{prop:seed} at
$(d,s)=(11,14)$ and Table~\ref{tab:prop} for $12\le d\le83$ provide positions
$\sigma^-_d\le s\le\sigma^+_d$ (with $\sigma^-_d=2d-8$) at which a general configuration has
$\rank J_P=3r-4$. In all of them $m\ge4$ (in fact $m\ge64$), so
Propositions~\ref{prop:Psi}, \ref{prop:lowerneg} and~\ref{prop:nonsmooth} give
$T^1_j(C(Z))=0$ for all $j<0$ and the absence of smoothings for general $Z$ at these
positions. We do not claim the lower endpoint $2d-8$ in degrees outside this table other
than through Theorem~\ref{thm:broad}.
\end{remark}

\begin{remark}\label{rem:scope}
The results leave open the positions $M_d-7d-426\le s\le M_d-d-6$ near the upper end of
the window \eqref{eq:window} ($6d+421$ positions in each degree $d\ge84$), the degrees
below those covered, the question whether $J_P$ has the expected rank $\min(3r-4,sm)$ at a
general configuration for every position $(d,s)$, and the smoothability of the cones when
$T^1_{-1}\ne0$. In particular Theorem~\ref{thm:outer} makes no assertion about
smoothability at $s=d+2$.
\end{remark}

% =====================================================================
% =====================================================================
\appendix
\section{Finite characteristic-zero certificates}\label{app:finite}

This appendix describes the computation used in the proof of
Theorem~\ref{thm:lowerboundary} for $29\le n\le88$, that is, $27\le d\le86$; it covers all
$3510$ characters of these $60$ degrees.

\emph{Data.} The coefficient vector is $(a_1,a_3,a_7)=(2,3,-6)$ throughout, so
$f=1+2u+3u^3-6u^7$, which satisfies the last condition of \eqref{eq:denomopen} and has no
root on the unit circle other than $1$ (Lemma~\ref{lem:denom}). For each $n$ a prime
$p\equiv1\pmod n$ and an element $\omega\in\mathbb F_p$ of exact order $n$ are chosen; the
order is checked at every prime divisor of $n$. For $n=87$ the first suitable prime has a
vanishing Fourier denominator, and $p=32713$ is used instead.

\emph{Checks of the denominator conditions.} The computation verifies, modulo $p$, that
$f(\omega^j)\ne0$ for $j\ne0$ (hence $\Pi\ne0$), that every $\lambda_j$ is nonzero, that
$f'(1)=-31\ne0$, that $2$, $n$ and all incidence pivots are nonzero, and that the Euclidean
algorithm gives $\gcd(g,u^6g(u^{-1}))=1$ for $g=f/(1-u)$. It also records a nonzero minor of
the auxiliary evaluation matrix on the $n+1$ edges $\{0,j\}$ ($1\le j<n$), $\{1,2\}$,
$\{1,3\}$; after division of the assigned rows by $\Pi$ its endpoint entries are
$1/f(\omega^{a-b})$ and its origin row consists of ones.

\emph{Rank statements.} For each character $\chi$ the computation forms the origin system
$\mathcal P_\chi$ of Proposition~\ref{prop:origin} on the six unknowns of
Corollary~\ref{cor:six} (seven in a global character, with $c$, $\eta$ and $\rho'$ or
$\kappa_e$ retained), the three moment rows $\mathcal W_\chi$ from
\eqref{eq:momentformula}, and the detector $\mathcal O_\chi$ from \eqref{eq:detector} on all
indices outside $U_\chi$, including the collision sets. It certifies
\[
  \rank\mathcal P_\chi=\begin{cases}4&\text{ordinary and global }\chi,\\
  5&\text{quadratic }\chi,\end{cases}\qquad
  \rank\begin{pmatrix}\mathcal P_\chi\\\mathcal O_\chi\mathcal W_\chi\end{pmatrix}=6,
\]
each rank being witnessed by a selected square minor, with its row and column indices and
its nonzero determinant modulo $p$. Since
$\rank[\mathcal P;\mathcal X]-\rank\mathcal P$
is the rank of $\mathcal X$ on $\ker\mathcal P$, this gives opening ranks $2$, $1$ and $2$ in
the three types of characters without using the injectivity of $\mathcal O_\chi$, which may
fail for small $n$ (for instance for $n=29$, $\chi=-14$ its rank outside $U_\chi$ is one,
which suffices for that one-dimensional special kernel).

\emph{Characteristic zero.} Let $\ZZ[\zeta_n]$ be localized at $n$ and at all coefficients and
denominators occurring in the computation. Since $p\nmid n$ and $\omega$ has order $n$,
$\zeta_n\mapsto\omega$ defines a ring homomorphism to $\mathbb F_p$, under which these
denominators map to the nonzero recorded residues. Every selected minor is the image of the
corresponding minor over $\QQ(\zeta_n)$, which is therefore nonzero. The modular residues are
not used as rational coordinates: the characteristic-zero witness is the cyclotomic point,
and the transfer to genuine configurations is that of Lemma~\ref{lem:transfer}.

\emph{Result.} Table~\ref{tab:finite} lists the data. In every degree the special kernel has
dimension $2n-2$, the opening has rank $2n-6$, the remaining kernel has dimension four, and
$\rank J_P=3\binom n3+3n-1=3r-4$.

\begin{table}[tp]\small\setlength{\tabcolsep}{4pt}
\caption{Finite certificates for $29\le n\le88$: $n=d+2$, $r=\binom n3+n+1$, the prime $p$, the element $\omega\in\mathbb F_p$ of order $n$, and $\rank J_P=3r-4$. In every case the special kernel has dimension $2n-2$ and the opening rank is $2n-6$.}\label{tab:finite}
\begin{tabular}{rrrrrr@{\qquad}rrrrrr}\toprule
$n$&$d$&$r$&$p$&$\omega$&$3r-4$&$n$&$d$&$r$&$p$&$\omega$&$3r-4$\\\midrule
29&27&3684&32191&18740&11048&59&57&32569&32569&6207&97703\\
30&28&4091&32191&205&12269&60&58&34281&32341&16632&102839\\
31&29&4527&32117&5220&13577&61&59&36052&32941&30384&108152\\
32&30&4993&32257&2292&14975&62&60&37883&32117&19658&113645\\
33&31&5490&32077&8638&16466&63&61&39775&32257&227&119321\\
34&32&6019&32029&25858&18053&64&62&41729&32257&624&125183\\
35&33&6581&32341&11264&19739&65&63&43746&33151&24488&131234\\
36&34&7177&32077&4077&21527&66&64&45827&32077&11439&137477\\
37&35&7808&32117&17815&23420&67&65&47973&32027&7157&143915\\
38&36&8475&32377&7278&25421&68&66&50185&32029&6271&150551\\
39&37&9179&32059&13433&27533&69&67&52464&32569&18953&157388\\
40&38&9921&32321&27732&29759&70&68&54811&32341&1612&164429\\
41&39&10702&32063&3482&32102&71&69&57227&32377&15197&171677\\
42&40&11523&32089&14261&34565&72&70&59713&32257&23701&179135\\
43&41&12385&32251&19979&37151&73&71&62270&32413&25819&186806\\
44&42&13289&32077&12182&39863&74&72&64899&32117&27518&194693\\
45&43&14236&32401&12258&42704&75&73&67601&32251&14038&202799\\
46&44&15227&32063&2282&45677&76&74&70377&32377&1183&211127\\
47&45&16263&32713&4574&48785&77&75&73228&32341&8220&219680\\
48&46&17345&32257&12683&52031&78&76&76155&32059&2712&228461\\
49&47&18474&32341&29183&55418&79&77&79159&32233&30818&237473\\
50&48&19651&32051&28637&58949&80&78&82241&32321&27827&246719\\
51&49&20877&32029&23450&62627&81&79&85402&32077&31795&256202\\
52&50&22153&32189&24068&66455&82&80&88643&32063&29394&265925\\
53&51&23480&32119&4239&70436&83&81&91965&32371&18905&275891\\
54&52&24859&32077&3568&74573&84&82&95369&32089&2318&286103\\
55&53&26291&32341&24631&78869&85&83&98856&33151&9733&296564\\
56&54&27777&32089&12505&83327&86&84&102427&32251&25757&307277\\
57&55&29318&32377&24283&87950&87&85&106083&32713&9871&318245\\
58&56&30915&32191&22087&92741&88&86&109825&32297&10780&329471\\
\bottomrule\end{tabular}\end{table}

\section{Leading coefficients for the boundary characters}\label{app:boundary}

This appendix lists the data used in the proof of Proposition~\ref{prop:boundary}. Throughout,
$n=6\mathsf N+\mathsf s$ with $0\le\mathsf s\le5$ and $\mathsf N\ge14$, and the \emph{order}
of a Laurent series in $\epsilon$ is the exponent of its first nonzero coefficient; in each
case all coefficients of lower order vanish.

\subsection*{Normalizations on the six-small-pole line}
For $h\in\{2,\dots,10,14\}$ the cofactor kernel vectors are those of Table~\ref{tab:cofactor}:
for each free column the free entry is the pivot minor, the pivot entries are the negatives
of the minors with the corresponding pivot column replaced, and the common power of $\epsilon$
is divided out. Table~\ref{tab:quadall} gives, for the quadratic characters, the leading term
of $\epsilon^{\mathsf N}\mathsf W_{11}$ on the unique cofactor vector. Table~\ref{tab:twoall}
gives, for $h\in\{3,5,7,9\}$, the leading term of the determinant of the $2\times2$ matrix
$(\epsilon^{\mathsf N}\mathsf W_A(K_i))$, where $K_0,K_1$ are the first two free cofactor
vectors and $A$ runs over the indicated pair of rows. Table~\ref{tab:minusone} gives the two
leading terms used for $\chi=-1$, with
\begin{equation}\label{eq:p3}
  \mathsf p_3(\mathsf N)=24\mathsf N^7+125\mathsf N^6+69\mathsf N^5-475\mathsf N^4
  -549\mathsf N^3+350\mathsf N^2+17736\mathsf N-20160 .
\end{equation}
For $\chi=0$ the matrix $\mathsf M$ is the constant term of the $3\times3$ matrix whose
$(i,k)$ entry is $\epsilon^{-r_i}$ times the $i$-th row of
$(\mathsf W_{11},\ \mathsf W_{13}-\mathsf W_{11},\ \mathsf W_{33}-2\mathsf W_{13}+\mathsf W_{11})$
evaluated on the $k$-th vector of $(\epsilon K_0,\epsilon^2K_1,\epsilon^{-1}K'_2)$; the row
weights $(r_0,r_1,r_2)$ are given with each matrix.

\subsection*{The sign test}
Each of the $36+24+12+6=78$ rational functions of $\mathsf N$ listed below (the coefficients in
Tables~\ref{tab:quadall}, \ref{tab:twoall} and~\ref{tab:minusone}, and the six determinants
$\det\mathsf M$) is written in lowest terms. After the substitution $\mathsf N=X+14$ its
numerator has all coefficients of one sign and a nonzero constant term, and its denominator
has positive coefficients. Hence none of these functions vanishes for real $\mathsf N\ge14$.
This is a statement about finitely many explicit polynomials in $X$, not a test at finitely
many values of $\mathsf N$.

{\small\renewcommand{\arraystretch}{1.9}
\begin{longtable}{cccl}
\caption{Quadratic characters $\chi=-h$: order and first nonzero coefficient of $\epsilon^{\mathsf N}\mathsf W_{11}$ on the cofactor kernel vector, for $n=6\mathsf N+\mathsf s$.}\label{tab:quadall}\\
\toprule $h$&$\mathsf s$&order&coefficient\\\midrule\endfirsthead
\toprule $h$&$\mathsf s$&order&coefficient\\\midrule\endhead
\bottomrule\endfoot
2&0&3&$\displaystyle \frac{2}{\mathsf N(\mathsf N - 1)}$\\
2&1&3&$\displaystyle \frac{1}{\mathsf N - 1}$\\
2&2&3&$\displaystyle 1$\\
2&3&2&$\displaystyle \frac{120}{\mathsf N(\mathsf N + 1)(\mathsf N + 2)(\mathsf N + 3)(\mathsf N + 4)}$\\
2&4&2&$\displaystyle \frac{24}{\mathsf N(\mathsf N + 1)(\mathsf N + 2)(\mathsf N + 3)}$\\
2&5&2&$\displaystyle \frac{6}{\mathsf N(\mathsf N + 1)(\mathsf N + 2)}$\\
4&0&2&$\displaystyle \frac{48}{\mathsf N(\mathsf N - 1)(\mathsf N + 1)(\mathsf N + 2)}$\\
4&1&2&$\displaystyle \frac{12}{\mathsf N(\mathsf N - 1)(\mathsf N + 1)}$\\
4&2&2&$\displaystyle \frac{4}{\mathsf N(\mathsf N - 1)}$\\
4&3&2&$\displaystyle \frac{2}{\mathsf N - 1}$\\
4&4&2&$\displaystyle 2$\\
4&5&1&$\displaystyle \frac{240}{\mathsf N(\mathsf N + 1)(\mathsf N + 2)(\mathsf N + 3)(\mathsf N + 4)}$\\
6&0&2&$\displaystyle -\frac{48(\mathsf N + 8)}{\mathsf N(\mathsf N - 1)(\mathsf N + 1)(\mathsf N + 2)(\mathsf N + 3)}$\\
6&1&2&$\displaystyle -\frac{12(\mathsf N + 6)}{\mathsf N(\mathsf N - 1)(\mathsf N + 1)(\mathsf N + 2)}$\\
6&2&2&$\displaystyle -\frac{4(\mathsf N + 4)}{\mathsf N(\mathsf N - 1)(\mathsf N + 1)}$\\
6&3&2&$\displaystyle -\frac{2(\mathsf N + 2)}{\mathsf N(\mathsf N - 1)}$\\
6&4&2&$\displaystyle -\frac{2\mathsf N}{\mathsf N - 1}$\\
6&5&1&$\displaystyle -\frac{240}{\mathsf N(\mathsf N + 1)(\mathsf N + 2)(\mathsf N + 3)(\mathsf N + 4)}$\\
8&0&1&$\displaystyle -\frac{2(\mathsf N - 2)}{\mathsf N(\mathsf N - 1)}$\\
8&1&1&$\displaystyle -\frac{2(\mathsf N - 2)}{\mathsf N - 1}$\\
8&2&0&$\displaystyle -\frac{240}{\mathsf N(\mathsf N + 1)(\mathsf N + 2)(\mathsf N + 3)(\mathsf N + 4)}$\\
8&3&0&$\displaystyle -\frac{48(\mathsf N - 1)}{\mathsf N(\mathsf N + 1)(\mathsf N + 2)(\mathsf N + 3)(\mathsf N + 4)}$\\
8&4&0&$\displaystyle -\frac{12(\mathsf N - 1)}{\mathsf N(\mathsf N + 1)(\mathsf N + 2)(\mathsf N + 3)}$\\
8&5&0&$\displaystyle -\frac{4(\mathsf N - 1)}{\mathsf N(\mathsf N + 1)(\mathsf N + 2)}$\\
10&0&1&$\displaystyle -\frac{6(\mathsf N - 2)}{\mathsf N(\mathsf N - 1)(\mathsf N + 1)(\mathsf N + 2)}$\\
10&1&1&$\displaystyle -\frac{2(\mathsf N - 2)}{\mathsf N(\mathsf N - 1)(\mathsf N + 1)}$\\
10&2&1&$\displaystyle -\frac{\mathsf N - 2}{\mathsf N(\mathsf N - 1)}$\\
10&3&1&$\displaystyle -\frac{\mathsf N - 2}{\mathsf N - 1}$\\
10&4&0&$\displaystyle -\frac{120}{\mathsf N(\mathsf N + 1)(\mathsf N + 2)(\mathsf N + 3)(\mathsf N + 4)}$\\
10&5&0&$\displaystyle -\frac{24(\mathsf N - 1)}{\mathsf N(\mathsf N + 1)(\mathsf N + 2)(\mathsf N + 3)(\mathsf N + 4)}$\\
14&0&3&$\displaystyle -\frac{2(3\mathsf N - 8)}{(\mathsf N - 2)(\mathsf N - 1)}$\\
14&1&3&$\displaystyle -\frac{6\mathsf N - 17}{\mathsf N - 2}$\\
14&2&2&$\displaystyle -\frac{720}{\mathsf N(\mathsf N - 1)(\mathsf N + 1)(\mathsf N + 2)(\mathsf N + 3)}$\\
14&3&2&$\displaystyle -\frac{24(6\mathsf N - 7)}{\mathsf N(\mathsf N - 1)(\mathsf N + 1)(\mathsf N + 2)(\mathsf N + 3)}$\\
14&4&2&$\displaystyle -\frac{12(3\mathsf N - 4)}{\mathsf N(\mathsf N - 1)(\mathsf N + 1)(\mathsf N + 2)}$\\
14&5&2&$\displaystyle -\frac{6(2\mathsf N - 3)}{\mathsf N(\mathsf N - 1)(\mathsf N + 1)}$\\
\end{longtable}}
{\small\renewcommand{\arraystretch}{1.9}
\begin{longtable}{ccccl}
\caption{Characters $\chi=-h$, $h\in\{3,5,7,9\}$: order and first nonzero coefficient of the $2\times2$ determinant of the rows $\epsilon^{\mathsf N}\mathsf W_A$ on the first two cofactor kernel vectors.}\label{tab:twoall}\\
\toprule $h$&rows&$\mathsf s$&order&coefficient\\\midrule\endfirsthead
\toprule $h$&rows&$\mathsf s$&order&coefficient\\\midrule\endhead
\bottomrule\endfoot
3&$(P_{11},P_{13})$&0&6&$\displaystyle -\frac{6(\mathsf N - 2)}{\mathsf N^{2}(\mathsf N - 1)^{2}(\mathsf N + 1)}$\\
3&$(P_{11},P_{13})$&1&6&$\displaystyle -\frac{2(\mathsf N - 2)}{\mathsf N(\mathsf N - 1)^{2}}$\\
3&$(P_{11},P_{13})$&2&5&$\displaystyle -\frac{120}{\mathsf N(\mathsf N - 1)(\mathsf N + 1)(\mathsf N + 2)(\mathsf N + 3)(\mathsf N + 4)}$\\
3&$(P_{11},P_{13})$&3&5&$\displaystyle -\frac{24(\mathsf N - 1)}{\mathsf N(\mathsf N + 1)(\mathsf N + 2)(\mathsf N + 3)(\mathsf N + 4)}$\\
3&$(P_{11},P_{13})$&4&4&$\displaystyle -\frac{720(\mathsf N - 1)}{\mathsf N^{2}(\mathsf N + 1)^{2}(\mathsf N + 2)^{2}(\mathsf N + 3)^{2}(\mathsf N + 4)}$\\
3&$(P_{11},P_{13})$&5&4&$\displaystyle -\frac{48(\mathsf N - 1)}{\mathsf N^{2}(\mathsf N + 1)^{2}(\mathsf N + 2)^{2}(\mathsf N + 3)}$\\
5&$(P_{11},P_{33})$&0&6&$\displaystyle -\frac{2880}{\mathsf N^{2}(\mathsf N - 1)^{2}(\mathsf N + 1)^{2}(\mathsf N + 2)^{2}(\mathsf N + 3)}$\\
5&$(P_{11},P_{33})$&1&6&$\displaystyle -\frac{144}{\mathsf N^{2}(\mathsf N - 1)^{2}(\mathsf N + 1)^{2}(\mathsf N + 2)}$\\
5&$(P_{11},P_{33})$&2&6&$\displaystyle -\frac{12}{\mathsf N^{2}(\mathsf N - 1)^{2}(\mathsf N + 1)}$\\
5&$(P_{11},P_{33})$&3&6&$\displaystyle -\frac{2}{\mathsf N(\mathsf N - 1)^{2}}$\\
5&$(P_{11},P_{33})$&4&6&$\displaystyle -\frac{1}{\mathsf N - 1}$\\
5&$(P_{11},P_{33})$&5&5&$\displaystyle -\frac{120}{\mathsf N(\mathsf N + 1)(\mathsf N + 2)(\mathsf N + 3)(\mathsf N + 4)}$\\
7&$(P_{11},P_{13})$&0&4&$\displaystyle -\frac{120}{\mathsf N(\mathsf N - 1)^{2}(\mathsf N + 1)(\mathsf N + 2)(\mathsf N + 3)}$\\
7&$(P_{11},P_{13})$&1&4&$\displaystyle -\frac{24}{\mathsf N(\mathsf N - 1)(\mathsf N + 1)(\mathsf N + 2)}$\\
7&$(P_{11},P_{13})$&2&3&$\displaystyle -\frac{720}{\mathsf N^{2}(\mathsf N - 1)(\mathsf N + 1)^{2}(\mathsf N + 2)(\mathsf N + 3)(\mathsf N + 4)}$\\
7&$(P_{11},P_{13})$&3&3&$\displaystyle -\frac{48}{\mathsf N^{2}(\mathsf N - 1)(\mathsf N + 1)(\mathsf N + 2)(\mathsf N + 3)}$\\
7&$(P_{11},P_{13})$&4&3&$\displaystyle -\frac{6(\mathsf N^{5} + 10\mathsf N^{4} + 35\mathsf N^{3} + 50\mathsf N^{2} - 96\mathsf N + 120)}{\mathsf N^{2}(\mathsf N - 1)(\mathsf N + 1)^{2}(\mathsf N + 2)^{2}(\mathsf N + 3)(\mathsf N + 4)}$\\
7&$(P_{11},P_{13})$&5&3&$\displaystyle -\frac{2(\mathsf N^{5} + 10\mathsf N^{4} + 35\mathsf N^{3} + 50\mathsf N^{2} - 216)}{\mathsf N^{2}(\mathsf N + 1)^{2}(\mathsf N + 2)(\mathsf N + 3)(\mathsf N + 4)}$\\
9&$(P_{11},P_{13})$&0&3&$\displaystyle -\frac{36}{\mathsf N^{2}(\mathsf N - 1)^{2}(\mathsf N + 1)}$\\
9&$(P_{11},P_{13})$&1&3&$\displaystyle -\frac{6}{\mathsf N(\mathsf N - 1)^{2}}$\\
9&$(P_{11},P_{13})$&2&3&$\displaystyle -\frac{3(\mathsf N + 5)(\mathsf N^{4} + 5\mathsf N^{3} + 10\mathsf N^{2} + 24)}{\mathsf N(\mathsf N - 1)(\mathsf N + 1)(\mathsf N + 2)(\mathsf N + 3)(\mathsf N + 4)}$\\
9&$(P_{11},P_{13})$&3&2&$\displaystyle -\frac{360}{\mathsf N(\mathsf N + 1)(\mathsf N + 2)(\mathsf N + 3)(\mathsf N + 4)}$\\
9&$(P_{11},P_{13})$&4&1&$\displaystyle -\frac{8640}{\mathsf N^{2}(\mathsf N + 1)^{2}(\mathsf N + 2)^{2}(\mathsf N + 3)^{2}(\mathsf N + 4)}$\\
9&$(P_{11},P_{13})$&5&1&$\displaystyle -\frac{432}{\mathsf N^{2}(\mathsf N + 1)^{2}(\mathsf N + 2)^{2}(\mathsf N + 3)}$\\
\end{longtable}}

\begin{table}[H]\small\renewcommand{\arraystretch}{2.1}\setlength{\tabcolsep}{4pt}
\caption{The character $\chi=-1$: the first nonzero coefficient of $\epsilon^{\mathsf N}\mathsf W_{11}$ on $K_0$, and of the projected row $\mathsf W_*$ on $K_*$, with their orders; $\mathsf p_3$ is the polynomial in \eqref{eq:p3}.}\label{tab:minusone}
\begin{tabular}{cclcl}\toprule $\mathsf s$&order&$\epsilon^{\mathsf N}\mathsf W_{11}(K_0)$&order&$\mathsf W_*(K_*)$\\\midrule
0&3&$\displaystyle \frac{1}{\mathsf N - 1}$&0&$\displaystyle \frac{24(6\mathsf N - 7)}{\mathsf N(\mathsf N - 1)(\mathsf N + 1)(\mathsf N + 2)(\mathsf N + 3)}$\\
1&3&$\displaystyle 1$&0&$\displaystyle \frac{24(6\mathsf N - 7)}{\mathsf N(\mathsf N - 1)(\mathsf N + 1)(\mathsf N + 2)(\mathsf N + 3)}$\\
2&2&$\displaystyle \frac{120}{\mathsf N(\mathsf N + 1)(\mathsf N + 2)(\mathsf N + 3)(\mathsf N + 4)}$&0&$\displaystyle \frac{24}{\mathsf N(\mathsf N - 1)(\mathsf N + 1)(\mathsf N + 2)}$\\
3&2&$\displaystyle \frac{24}{\mathsf N(\mathsf N + 1)(\mathsf N + 2)(\mathsf N + 3)}$&1&$\displaystyle \frac{\mathsf p_3(\mathsf N)}{30\mathsf N(\mathsf N - 1)(\mathsf N + 1)(\mathsf N + 2)(\mathsf N + 3)}$\\
4&2&$\displaystyle \frac{6}{\mathsf N(\mathsf N + 1)(\mathsf N + 2)}$&0&$\displaystyle \frac{120(\mathsf N - 1)}{\mathsf N(\mathsf N + 1)(\mathsf N + 2)(\mathsf N + 3)(\mathsf N + 4)}$\\
5&2&$\displaystyle \frac{2}{\mathsf N(\mathsf N + 1)}$&0&$\displaystyle \frac{24(6\mathsf N - 1)}{\mathsf N(\mathsf N + 1)(\mathsf N + 2)(\mathsf N + 3)(\mathsf N + 4)}$\\
\bottomrule\end{tabular}\end{table}

\subsection*{The character zero}
\paragraph{$\mathsf s=0$} Row weights $(0,1,1)$; the constant matrix has entries
{\small\allowdisplaybreaks\begin{align*}
  \mathsf M_{00}&=-\dfrac{9(\mathsf N - 2)(5\mathsf N + 3)}{\mathsf N(\mathsf N - 1)(\mathsf N + 1)(\mathsf N + 2)(\mathsf N + 3)},\\
  \mathsf M_{01}&=-\dfrac{60}{\mathsf N(\mathsf N - 1)(\mathsf N + 1)(\mathsf N + 2)(\mathsf N + 3)},\\
  \mathsf M_{02}&=-\dfrac{\mathsf N^{5} + 5\mathsf N^{4} + 5\mathsf N^{3} - 5\mathsf N^{2} - 6\mathsf N - 840}{2\mathsf N(\mathsf N - 1)(\mathsf N + 1)(\mathsf N + 2)(\mathsf N + 3)},\\
  \mathsf M_{10}&=\dfrac{(\mathsf N - 2)(7\mathsf N^{7} + 76\mathsf N^{6} + 254\mathsf N^{5} + 240\mathsf N^{4} - 177\mathsf N^{3} - 316\mathsf N^{2} + 21516\mathsf N + 12960)}{240\mathsf N(\mathsf N - 1)(\mathsf N + 1)(\mathsf N + 2)(\mathsf N + 3)},\\
  \mathsf M_{11}&=\dfrac{5(\mathsf N^{6} + 3\mathsf N^{5} - 5\mathsf N^{4} - 15\mathsf N^{3} + 4\mathsf N^{2} + 12\mathsf N + 288)}{12\mathsf N(\mathsf N - 1)(\mathsf N + 1)(\mathsf N + 2)(\mathsf N + 3)},\\
  \mathsf M_{12}&=-\dfrac{35\mathsf N^{6} + 93\mathsf N^{5} - 235\mathsf N^{4} - 585\mathsf N^{3} + 200\mathsf N^{2} + 492\mathsf N + 10080}{12\mathsf N(\mathsf N - 1)(\mathsf N + 1)(\mathsf N + 2)(\mathsf N + 3)},\\
  \mathsf M_{20}&=\dfrac{\mathsf N^{4} + 2\mathsf N^{3} - \mathsf N^{2} - 38\mathsf N + 24}{2\mathsf N(\mathsf N - 1)(\mathsf N + 1)(\mathsf N + 2)},\\
  \mathsf M_{21}&=0,\\
  \mathsf M_{22}&=0,
\end{align*}}
and
\[\det\mathsf M=\frac{5(\mathsf N - 2)(\mathsf N^{4} + 2\mathsf N^{3} - \mathsf N^{2} - 38\mathsf N + 24)}{48\mathsf N(\mathsf N - 1)(\mathsf N + 1)(\mathsf N + 2)}.\]
\paragraph{$\mathsf s=1$} Row weights $(0,1,1)$; the constant matrix has entries
{\small\allowdisplaybreaks\begin{align*}
  \mathsf M_{00}&=-\dfrac{9(\mathsf N - 2)(5\mathsf N + 3)}{\mathsf N(\mathsf N - 1)(\mathsf N + 1)(\mathsf N + 2)(\mathsf N + 3)},\\
  \mathsf M_{01}&=-\dfrac{60}{\mathsf N(\mathsf N - 1)(\mathsf N + 1)(\mathsf N + 2)(\mathsf N + 3)},\\
  \mathsf M_{02}&=-\dfrac{\mathsf N^{5} + 5\mathsf N^{4} + 5\mathsf N^{3} - 5\mathsf N^{2} - 6\mathsf N - 840}{2\mathsf N(\mathsf N - 1)(\mathsf N + 1)(\mathsf N + 2)(\mathsf N + 3)},\\
  \mathsf M_{10}&=\dfrac{(\mathsf N - 2)(5\mathsf N^{7} + 38\mathsf N^{6} + 78\mathsf N^{5} - 20\mathsf N^{4} - 155\mathsf N^{3} - 18\mathsf N^{2} + 4392\mathsf N + 2592)}{48\mathsf N(\mathsf N - 1)(\mathsf N + 1)(\mathsf N + 2)(\mathsf N + 3)},\\
  \mathsf M_{11}&=\dfrac{3(\mathsf N^{6} + 3\mathsf N^{5} - 5\mathsf N^{4} - 15\mathsf N^{3} + 4\mathsf N^{2} + 12\mathsf N + 160)}{4\mathsf N(\mathsf N - 1)(\mathsf N + 1)(\mathsf N + 2)(\mathsf N + 3)},\\
  \mathsf M_{12}&=-\dfrac{21\mathsf N^{6} + 59\mathsf N^{5} - 125\mathsf N^{4} - 335\mathsf N^{3} + 104\mathsf N^{2} + 276\mathsf N + 3360}{4\mathsf N(\mathsf N - 1)(\mathsf N + 1)(\mathsf N + 2)(\mathsf N + 3)},\\
  \mathsf M_{20}&=\dfrac{(\mathsf N - 2)(\mathsf N^{2} + 2\mathsf N + 3)}{2\mathsf N(\mathsf N - 1)(\mathsf N + 1)},\\
  \mathsf M_{21}&=0,\\
  \mathsf M_{22}&=0,
\end{align*}}
and
\[\det\mathsf M=\frac{3(\mathsf N - 2)^{2}(\mathsf N^{2} + 2\mathsf N + 3)}{16\mathsf N(\mathsf N - 1)(\mathsf N + 1)}.\]
\paragraph{$\mathsf s=2$} Row weights $(0,0,1)$; the constant matrix has entries
{\small\allowdisplaybreaks\begin{align*}
  \mathsf M_{00}&=-\dfrac{3(15\mathsf N^{3} + 9\mathsf N^{2} + 8\mathsf N - 52)}{\mathsf N(\mathsf N - 1)(\mathsf N + 1)(\mathsf N + 2)(\mathsf N + 3)(\mathsf N + 4)},\\
  \mathsf M_{01}&=-\dfrac{60}{\mathsf N(\mathsf N + 1)(\mathsf N + 2)(\mathsf N + 3)(\mathsf N + 4)},\\
  \mathsf M_{02}&=-\dfrac{\mathsf N^{5} + 10\mathsf N^{4} + 35\mathsf N^{3} + 50\mathsf N^{2} + 24\mathsf N - 840}{2\mathsf N(\mathsf N + 1)(\mathsf N + 2)(\mathsf N + 3)(\mathsf N + 4)},\\
  \mathsf M_{10}&=\dfrac{30(\mathsf N - 2)}{\mathsf N(\mathsf N + 2)(\mathsf N + 3)(\mathsf N + 4)},\\
  \mathsf M_{11}&=\dfrac{60}{\mathsf N(\mathsf N + 1)(\mathsf N + 2)(\mathsf N + 3)(\mathsf N + 4)},\\
  \mathsf M_{12}&=-\dfrac{420}{\mathsf N(\mathsf N + 1)(\mathsf N + 2)(\mathsf N + 3)(\mathsf N + 4)},\\
  \mathsf M_{20}&=\dfrac{1}{2},\\
  \mathsf M_{21}&=0,\\
  \mathsf M_{22}&=0,
\end{align*}}
and
\[\det\mathsf M=\frac{15}{\mathsf N(\mathsf N + 1)(\mathsf N + 2)(\mathsf N + 3)(\mathsf N + 4)}.\]
\paragraph{$\mathsf s=3$} Row weights $(0,0,1)$; the constant matrix has entries
{\small\allowdisplaybreaks\begin{align*}
  \mathsf M_{00}&=-\dfrac{9(5\mathsf N^{3} - \mathsf N^{2} - 4\mathsf N + 20)}{\mathsf N(\mathsf N - 1)(\mathsf N + 1)(\mathsf N + 2)(\mathsf N + 3)(\mathsf N + 4)},\\
  \mathsf M_{01}&=-\dfrac{60}{\mathsf N(\mathsf N + 1)(\mathsf N + 2)(\mathsf N + 3)(\mathsf N + 4)},\\
  \mathsf M_{02}&=-\dfrac{\mathsf N^{5} + 10\mathsf N^{4} + 35\mathsf N^{3} + 50\mathsf N^{2} + 24\mathsf N - 840}{2\mathsf N(\mathsf N + 1)(\mathsf N + 2)(\mathsf N + 3)(\mathsf N + 4)},\\
  \mathsf M_{10}&=\dfrac{6(7\mathsf N^{2} + \mathsf N - 18)}{\mathsf N(\mathsf N + 1)(\mathsf N + 2)(\mathsf N + 3)(\mathsf N + 4)},\\
  \mathsf M_{11}&=\dfrac{60}{\mathsf N(\mathsf N + 1)(\mathsf N + 2)(\mathsf N + 3)(\mathsf N + 4)},\\
  \mathsf M_{12}&=-\dfrac{420}{\mathsf N(\mathsf N + 1)(\mathsf N + 2)(\mathsf N + 3)(\mathsf N + 4)},\\
  \mathsf M_{20}&=\dfrac{\mathsf N^{4} + 9\mathsf N^{3} + 26\mathsf N^{2} - 36\mathsf N + 120}{2\mathsf N(\mathsf N + 2)(\mathsf N + 3)(\mathsf N + 4)},\\
  \mathsf M_{21}&=-\dfrac{60}{\mathsf N(\mathsf N + 1)(\mathsf N + 2)(\mathsf N + 3)(\mathsf N + 4)},\\
  \mathsf M_{22}&=\dfrac{420}{\mathsf N(\mathsf N + 1)(\mathsf N + 2)(\mathsf N + 3)(\mathsf N + 4)},
\end{align*}}
and
\[\det\mathsf M=\frac{15(\mathsf N^{4} + 6\mathsf N^{3} + 11\mathsf N^{2} + 30\mathsf N - 24)}{\mathsf N^{2}(\mathsf N + 1)^{2}(\mathsf N + 2)^{2}(\mathsf N + 3)^{2}(\mathsf N + 4)}.\]
\paragraph{$\mathsf s=4$} Row weights $(0,0,1)$; the constant matrix has entries
{\small\allowdisplaybreaks\begin{align*}
  \mathsf M_{00}&=-\dfrac{9(\mathsf N - 1)(5\mathsf N + 8)}{\mathsf N(\mathsf N + 1)(\mathsf N + 2)(\mathsf N + 3)(\mathsf N + 4)},\\
  \mathsf M_{01}&=-\dfrac{60}{\mathsf N(\mathsf N + 1)(\mathsf N + 2)(\mathsf N + 3)(\mathsf N + 4)},\\
  \mathsf M_{02}&=-\dfrac{\mathsf N^{5} + 10\mathsf N^{4} + 35\mathsf N^{3} + 50\mathsf N^{2} + 24\mathsf N - 840}{2\mathsf N(\mathsf N + 1)(\mathsf N + 2)(\mathsf N + 3)(\mathsf N + 4)},\\
  \mathsf M_{10}&=\dfrac{3(5\mathsf N - 1)}{\mathsf N(\mathsf N + 1)(\mathsf N + 2)(\mathsf N + 3)},\\
  \mathsf M_{11}&=0,\\
  \mathsf M_{12}&=0,\\
  \mathsf M_{20}&=\dfrac{\mathsf N^{5} + 10\mathsf N^{4} + 35\mathsf N^{3} - 94\mathsf N^{2} + 72\mathsf N + 336}{2\mathsf N(\mathsf N + 1)(\mathsf N + 2)(\mathsf N + 3)(\mathsf N + 4)},\\
  \mathsf M_{21}&=-\dfrac{120}{\mathsf N(\mathsf N + 1)(\mathsf N + 2)(\mathsf N + 3)(\mathsf N + 4)},\\
  \mathsf M_{22}&=\dfrac{840}{\mathsf N(\mathsf N + 1)(\mathsf N + 2)(\mathsf N + 3)(\mathsf N + 4)},
\end{align*}}
and
\[\det\mathsf M=\frac{180(5\mathsf N - 1)}{\mathsf N^{2}(\mathsf N + 1)^{2}(\mathsf N + 2)^{2}(\mathsf N + 3)^{2}(\mathsf N + 4)}.\]
\paragraph{$\mathsf s=5$} Row weights $(0,0,1)$; the constant matrix has entries
{\small\allowdisplaybreaks\begin{align*}
  \mathsf M_{00}&=-\dfrac{9(\mathsf N - 1)(5\mathsf N + 8)}{\mathsf N(\mathsf N + 1)(\mathsf N + 2)(\mathsf N + 3)(\mathsf N + 4)},\\
  \mathsf M_{01}&=-\dfrac{60}{\mathsf N(\mathsf N + 1)(\mathsf N + 2)(\mathsf N + 3)(\mathsf N + 4)},\\
  \mathsf M_{02}&=-\dfrac{\mathsf N^{5} + 10\mathsf N^{4} + 35\mathsf N^{3} + 50\mathsf N^{2} + 24\mathsf N - 840}{2\mathsf N(\mathsf N + 1)(\mathsf N + 2)(\mathsf N + 3)(\mathsf N + 4)},\\
  \mathsf M_{10}&=\dfrac{3}{\mathsf N(\mathsf N + 1)(\mathsf N + 2)},\\
  \mathsf M_{11}&=0,\\
  \mathsf M_{12}&=0,\\
  \mathsf M_{20}&=\dfrac{\mathsf N^{5} + 10\mathsf N^{4} + 35\mathsf N^{3} - 64\mathsf N^{2} - 102\mathsf N + 240}{2\mathsf N(\mathsf N + 1)(\mathsf N + 2)(\mathsf N + 3)(\mathsf N + 4)},\\
  \mathsf M_{21}&=-\dfrac{60}{\mathsf N(\mathsf N + 1)(\mathsf N + 2)(\mathsf N + 3)(\mathsf N + 4)},\\
  \mathsf M_{22}&=\dfrac{420}{\mathsf N(\mathsf N + 1)(\mathsf N + 2)(\mathsf N + 3)(\mathsf N + 4)},
\end{align*}}
and
\[\det\mathsf M=\frac{90}{\mathsf N^{2}(\mathsf N + 1)^{2}(\mathsf N + 2)^{2}(\mathsf N + 3)(\mathsf N + 4)}.\]


\subsection*{The cubic line}
Let $\mathsf c(T)$ be as in \eqref{eq:limkernel} and $d_k=2^k+T(-3)^k$. The obstruction
functions of the proof of Proposition~\ref{prop:boundary} are, as identities of rational
functions,
\begin{align*}
  \mathcal F_1&=\mathcal F_2=9072+\mathsf c(T)(1+T)
   =-\frac{1296(27T-8)\,\mathsf q_1(T)}{175\,\Delta_0(T)},\\
  \mathcal F_3&=-173628+\mathsf c(T)\bigl[78d_{-3}+\tfrac{931}6d_{-1}\bigr]
   =\frac{36\,\mathsf q_3(T)}{175\,\Delta_0(T)},\\
  \mathcal F_4&=-1462104-133\,\mathsf c(T)d_0=\frac{216\,\mathsf q_4(T)}{25\,\Delta_0(T)},\\
  \mathcal F_5&=468\bigl[-468-\tfrac76\mathsf c(T)d_{-1}\bigr]\\
   &=-\frac{16848(9T+4)\,\mathsf q_5(T)}{25T(8T-27)(16T+81)(8736T^2+1277T-1134)},\\
  \mathcal F_6&=468\bigl[9072+\mathsf c(T)d_0\bigr]=468\,\mathcal F_1,\\
  \mathcal F_{n-13}&=TF_{13}(T)=\frac{216\,\mathsf q_{13}(T)}{25(2T-3)(16T+81)(8736T^2+1277T-1134)},
\end{align*}
where $\Delta_0(T)=T(2T-3)(8T-27)(16T+81)(8736T^2+1277T-1134)$ and
\begin{align*}
  \mathsf q_1&=135364608T^6-549788672T^5-1798393680T^4+3550155660T^3\\
   &\qquad+1385393220T^2-1202021127T-442158912,\\
  \mathsf q_3&=7185424121856T^7-46490670348288T^6+417002960096T^5\\
   &\qquad+282658729970760T^4-324936129057390T^3-5712883184349T^2\\
   &\qquad+60310985662062T-11121180954624,\\
  \mathsf q_4&=416652263424T^7-1881844058112T^6-5056118526176T^5\\
   &\qquad+13871807147640T^4-477371748510T^3-5379518823921T^2\\
   &\qquad-44678717262T+403248927744,\\
  \mathsf q_5&=406093824T^5-2011489280T^4-1965486960T^3+9368244540T^2\\
   &\qquad-5531611905T+884317824,\\
  \mathsf q_{13}&=152285184T^5-310586240T^4+462770950T^3-369645495T^2\\
   &\qquad-145657845T+65505024 .
\end{align*}
For $1\le\chi\le6$ the value of $T$ lies in the interval $[T^-,T^+]$ of
Table~\ref{tab:cubicint}, which contains $T_\chi(1\pm10^{-8})$; for $h=13$ the interval is
$[-10^{-4},10^{-4}]$.

\emph{Horner interval evaluation.} For a polynomial $a_mT^m+\dots+a_0$ with rational
coefficients and a rational interval $[T^-,T^+]$, start with $V=[0,0]$ and replace $V$
successively by $V\cdot[T^-,T^+]+[a_i,a_i]$ for $i=m,\dots,0$, where the product of two
intervals is the interval spanned by the four products of endpoints. The result is a
rational interval containing all values of the polynomial on $[T^-,T^+]$. Applied to the
numerator and the denominator of each $\mathcal F_\chi$ and of $\det(\mathcal A,\mathcal C)$
in \eqref{eq:DAC}, it gives exact rational bounds, and Table~\ref{tab:cubicbounds} lists
these bounds rounded outward to integers. The exact rational bounds for the numerators and
denominators of the seven functions $\mathcal F_\chi$ are stored in the boundary cubic
certificate (Appendix~\ref{app:data}). For $\det(\mathcal A,\mathcal C)$ the certificate
program checks the nonvanishing but does not store the exact bounds; the integer bounds in
Table~\ref{tab:cubicbounds} are reproduced by applying the procedure above to the
expression \eqref{eq:DAC}. None of the
intervals contains $0$, so all the functions are nonzero on the relevant intervals. For
$h=13$ the factor $T$ is removed from the numerator of $\det(\mathcal A,\mathcal C)$; since
$T\ne0$, this suffices.

\begin{table}[H]\small\centering
\caption{The cubic line: outward integer bounds, obtained by Horner interval evaluation on the intervals of Table~\ref{tab:cubicint}, for the numerator $\mathcal N$ and the denominator $\mathcal D$ of the obstruction function $\mathcal F_\chi$ (first block) and of $\det(\mathcal A,\mathcal C)$ (second block). For $h=13$ the function is $TF_{13}(T)$, and the factor $T$ is removed from the numerator of $\det(\mathcal A,\mathcal C)$.}\label{tab:cubicbounds}
\begin{tabular}{crr}\toprule
$\chi$&$\mathcal N(\mathcal F_\chi)\in$&$\mathcal D(\mathcal F_\chi)\in$\\\midrule
1&$[1093854649755,1093854946711]$&$[115903399,115903422]$\\
2&$[-2512434530478,-2512434417064]$&$[302411944,302411954]$\\
3&$[-104962328052535,-104962320384825]$&$[-113951301,-113951295]$\\
4&$[76006215754635,76006216235690]$&$[21601255,21601257]$\\
5&$[-42459944942637,-42459944579737]$&$[3778287,3778288]$\\
6&$[-2153885545348037,-2153885544034488]$&$[64562671,64562674]$\\
$n-13$&$[14145938176,14152232192]$&$[6887950,6890150]$\\
\bottomrule\end{tabular}\par\medskip
\begin{tabular}{crr}\toprule
$\chi$&$\mathcal N(\det)\in$&$\mathcal D(\det)\in$\\\midrule
1&$[-1303490736,-1303490536]$&$[-1422521450,-1422521141]$\\
2&$[-2116596437,-2116596365]$&$[4559499630,4559499719]$\\
3&$[1067577420,1067577461]$&$[3514487454,3514487528]$\\
4&$[-1155769381,-1155769354]$&$[5934244452,5934244468]$\\
5&$[659494726,659494742]$&$[5076129515,5076129536]$\\
6&$[-514586161,-514586149]$&$[6006176585,6006176592]$\\
$n-13$&$[10743864416,10746286084]$&$[5804067075,5805438717]$\\
\bottomrule\end{tabular}\end{table}

\begin{table}[H]\small\renewcommand{\arraystretch}{1.6}
\caption{The intervals for $T$ on the cubic line.}\label{tab:cubicint}
\begin{tabular}{ccc}\toprule $\chi$&$T^-$&$T^+$\\\midrule
1&$\frac{11111111}{31640625}$&$\frac{100000001}{284765625}$\\
2&$-\frac{200000002}{854296875}$&$-\frac{22222222}{94921875}$\\
3&$\frac{44444444}{284765625}$&$\frac{400000004}{2562890625}$\\
4&$-\frac{800000008}{7688671875}$&$-\frac{88888888}{854296875}$\\
5&$\frac{177777776}{2562890625}$&$\frac{1600000016}{23066015625}$\\
6&$-\frac{3200000032}{69198046875}$&$-\frac{355555552}{7688671875}$\\
$n-13$&$-\frac{1}{10000}$&$\frac{1}{10000}$\\
\bottomrule\end{tabular}\end{table}

\section{The rational seed}\label{app:seed}

Table~\ref{tab:seed} lists the rational data used in Proposition~\ref{prop:seed}. The planes
are $\ell_i=1+\nu_i\cdot x$, the points $p_i$ satisfy $\ell_i(p_i)=0$, and the auxiliary sets
are $B_\lambda=\{p_i+\lambda w_i\}\cup\{0\}$ for $\lambda=1,2$. The $286$ triple points are
computed exactly over $\QQ$. All denominators are prime to $P=32003$; modulo $P$, the
$896\times896$ minor selected in the matrix of Proposition~\ref{prop:seed} has determinant
$19954$ for $\lambda=1$ and $31377$ for $\lambda=2$. The four motions lie in the kernel:
a translation by $b$ has star coordinates $t_{T,k}=\nu_k\cdot b$ and auxiliary velocities $b$,
and the dilation has $t_{T,k}=-1$ and auxiliary velocities $b_j$.

\begin{table}[H]\small
\caption{Rational seed data ($n=13$): first two coordinates of the normals $\nu_i$ and of the directions $w_i$, and the points $p_i$. The third coordinates are determined by $\nu_i\cdot p_i=-1$ and $\nu_i\cdot w_i=1$.}\label{tab:seed}
\begin{tabular}{rrrrrr}\toprule $i$&$\nu_{i,1}$&$\nu_{i,2}$&$p_i$&$w_{i,1}$&$w_{i,2}$\\\midrule
0&17521&5775&$(28534,25025,26602)$&8751&30617\\
1&6234&5692&$(13687,16771,7786)$&14540&22697\\
2&24365&2617&$(28262,15768,2758)$&25222&12214\\
3&14749&32001&$(29918,3997,20899)$&4513&12875\\
4&17881&2230&$(14805,20269,15010)$&31451&23560\\
5&29589&13848&$(16491,25687,1222)$&13330&26386\\
6&7165&21490&$(10509,30595,11898)$&15671&28459\\
7&9362&26419&$(10250,16293,22095)$&18449&25337\\
8&21180&6727&$(931,13174,4843)$&13326&6294\\
9&17687&21176&$(22553,28062,10781)$&6191&13668\\
10&5041&20138&$(24957,23154,17238)$&4835&8131\\
11&13357&23912&$(19500,17261,3025)$&7472&516\\
12&11928&25717&$(8088,29149,29577)$&24519&13733\\
\bottomrule\end{tabular}\end{table}

\section{The finite propagation}\label{app:propagation}

Table~\ref{tab:prop} lists the $188$ transitions from degree $12$ to degree $84$ used in
Section~\ref{sec:seed}. Each residual interval is contained in the interval proved in the
previous degree, each transition satisfies the hypotheses of
Proposition~\ref{prop:gateA} or~\ref{prop:gatesBC}, and the images cover the stated
interval.

{\small\begin{longtable}{rrp{0.68\textwidth}}
\caption{Finite propagation. For each degree $d$ the proved interval $[\sigma^-_d,\sigma^+_d]$ of positions and the transitions used, written $X_c[\ell,u]$: $X=\mathrm S$ (sparse, Proposition~\ref{prop:gateA}), $\mathrm L$ (low $c$) or $\mathrm H$ (high $c$, Proposition~\ref{prop:gatesBC}), with residual positions $[\ell,u]$ in degree $d-1$ and image $[\ell+c,u+c]$.}\label{tab:prop}\\
\toprule $d$&$[\sigma^-_d,\sigma^+_d]$&transitions\\\midrule\endfirsthead
\toprule $d$&$[\sigma^-_d,\sigma^+_d]$&transitions\\\midrule\endhead
\bottomrule\endfoot
11&$[14,14]$&seed (Proposition~\ref{prop:seed})\\
12&$[16,16]$&$\mathrm{S}_{2}[14,14]$\\
13&$[18,18]$&$\mathrm{S}_{2}[16,16]$\\
14&$[20,20]$&$\mathrm{S}_{2}[18,18]$\\
15&$[22,22]$&$\mathrm{S}_{2}[20,20]$\\
16&$[24,24]$&$\mathrm{S}_{2}[22,22]$\\
17&$[26,26]$&$\mathrm{S}_{2}[24,24]$\\
18&$[28,28]$&$\mathrm{S}_{2}[26,26]$\\
19&$[30,30]$&$\mathrm{S}_{2}[28,28]$\\
20&$[32,32]$&$\mathrm{S}_{2}[30,30]$\\
21&$[34,34]$&$\mathrm{S}_{2}[32,32]$\\
22&$[36,37]$&$\mathrm{S}_{2}[34,34]$, $\mathrm{S}_{3}[34,34]$\\
23&$[38,40]$&$\mathrm{S}_{2}[36,37]$, $\mathrm{S}_{3}[36,37]$\\
24&$[40,43]$&$\mathrm{S}_{2}[38,40]$, $\mathrm{S}_{3}[38,40]$\\
25&$[42,46]$&$\mathrm{S}_{2}[40,43]$, $\mathrm{S}_{3}[40,43]$\\
26&$[44,49]$&$\mathrm{S}_{2}[42,46]$, $\mathrm{S}_{3}[42,46]$\\
27&$[46,52]$&$\mathrm{S}_{2}[44,49]$, $\mathrm{S}_{3}[44,49]$\\
28&$[48,55]$&$\mathrm{S}_{2}[46,52]$, $\mathrm{S}_{3}[46,52]$\\
29&$[50,58]$&$\mathrm{S}_{2}[48,55]$, $\mathrm{S}_{3}[48,55]$\\
30&$[52,61]$&$\mathrm{S}_{2}[50,58]$, $\mathrm{S}_{3}[50,58]$\\
31&$[54,64]$&$\mathrm{S}_{2}[52,61]$, $\mathrm{S}_{3}[52,61]$\\
32&$[56,67]$&$\mathrm{S}_{2}[54,64]$, $\mathrm{S}_{3}[54,64]$\\
33&$[58,83]$&$\mathrm{S}_{2}[56,67]$, $\mathrm{S}_{3}[56,67]$, $\mathrm{L}_{15}[56,67]$, $\mathrm{L}_{16}[56,67]$\\
34&$[60,100]$&$\mathrm{S}_{2}[58,83]$, $\mathrm{L}_{17}[58,83]$\\
35&$[62,121]$&$\mathrm{S}_{2}[60,100]$, $\mathrm{L}_{17}[60,100]$, $\mathrm{H}_{21}[92,100]$\\
36&$[64,142]$&$\mathrm{S}_{2}[62,121]$, $\mathrm{H}_{21}[95,121]$\\
37&$[66,166]$&$\mathrm{S}_{2}[64,142]$, $\mathrm{H}_{22}[97,142]$, $\mathrm{H}_{24}[136,142]$\\
38&$[68,191]$&$\mathrm{S}_{2}[66,166]$, $\mathrm{H}_{25}[139,166]$\\
39&$[70,218]$&$\mathrm{S}_{2}[68,191]$, $\mathrm{H}_{25}[143,191]$, $\mathrm{H}_{27}[185,191]$\\
40&$[72,246]$&$\mathrm{S}_{2}[70,212]$, $\mathrm{H}_{26}[146,218]$, $\mathrm{H}_{28}[189,218]$\\
41&$[74,276]$&$\mathrm{S}_{2}[72,214]$, $\mathrm{H}_{27}[149,246]$, $\mathrm{H}_{30}[239,246]$\\
42&$[76,307]$&$\mathrm{S}_{2}[74,233]$, $\mathrm{H}_{30}[197,276]$, $\mathrm{H}_{31}[244,276]$\\
43&$[78,340]$&$\mathrm{S}_{2}[76,235]$, $\mathrm{H}_{30}[202,307]$, $\mathrm{H}_{33}[298,307]$\\
44&$[80,374]$&$\mathrm{S}_{2}[78,255]$, $\mathrm{H}_{31}[206,340]$, $\mathrm{H}_{34}[304,340]$\\
45&$[82,409]$&$\mathrm{S}_{2}[80,257]$, $\mathrm{H}_{31}[211,374]$, $\mathrm{H}_{35}[310,374]$\\
46&$[84,445]$&$\mathrm{S}_{2}[82,278]$, $\mathrm{H}_{32}[215,409]$, $\mathrm{H}_{36}[370,409]$\\
47&$[86,482]$&$\mathrm{S}_{2}[84,280]$, $\mathrm{H}_{33}[219,445]$, $\mathrm{H}_{37}[377,445]$\\
48&$[88,520]$&$\mathrm{S}_{2}[86,302]$, $\mathrm{H}_{34}[223,482]$, $\mathrm{H}_{38}[384,482]$\\
49&$[90,560]$&$\mathrm{S}_{2}[88,304]$, $\mathrm{H}_{35}[227,520]$, $\mathrm{H}_{40}[507,520]$\\
50&$[92,601]$&$\mathrm{S}_{2}[90,327]$, $\mathrm{H}_{37}[286,560]$, $\mathrm{H}_{41}[516,560]$\\
51&$[94,644]$&$\mathrm{S}_{2}[92,329]$, $\mathrm{H}_{37}[292,601]$, $\mathrm{H}_{43}[586,601]$\\
52&$[96,688]$&$\mathrm{S}_{2}[94,353]$, $\mathrm{H}_{38}[297,644]$, $\mathrm{H}_{44}[596,644]$\\
53&$[98,733]$&$\mathrm{S}_{2}[96,355]$, $\mathrm{H}_{39}[302,688]$, $\mathrm{H}_{45}[671,688]$\\
54&$[100,779]$&$\mathrm{S}_{2}[98,380]$, $\mathrm{H}_{40}[307,733]$, $\mathrm{H}_{46}[682,733]$\\
55&$[102,826]$&$\mathrm{S}_{2}[100,382]$, $\mathrm{H}_{41}[312,779]$, $\mathrm{H}_{47}[693,779]$\\
56&$[104,874]$&$\mathrm{S}_{2}[102,408]$, $\mathrm{H}_{42}[317,826]$, $\mathrm{H}_{48}[704,826]$\\
57&$[106,923]$&$\mathrm{S}_{2}[104,410]$, $\mathrm{H}_{42}[323,874]$, $\mathrm{H}_{49}[785,874]$\\
58&$[108,973]$&$\mathrm{S}_{2}[106,437]$, $\mathrm{H}_{44}[392,923]$, $\mathrm{H}_{50}[797,923]$\\
59&$[110,1024]$&$\mathrm{S}_{2}[108,439]$, $\mathrm{H}_{43}[334,973]$, $\mathrm{H}_{51}[809,973]$\\
60&$[112,1076]$&$\mathrm{S}_{2}[110,467]$, $\mathrm{H}_{46}[404,1024]$, $\mathrm{H}_{52}[821,1024]$\\
61&$[114,1129]$&$\mathrm{S}_{2}[112,469]$, $\mathrm{H}_{47}[410,1076]$, $\mathrm{H}_{53}[908,1076]$\\
62&$[116,1183]$&$\mathrm{S}_{2}[114,498]$, $\mathrm{H}_{48}[416,1129]$, $\mathrm{H}_{54}[921,1129]$\\
63&$[118,1238]$&$\mathrm{S}_{2}[116,500]$, $\mathrm{H}_{49}[422,1183]$, $\mathrm{H}_{55}[934,1183]$\\
64&$[120,1294]$&$\mathrm{S}_{2}[118,530]$, $\mathrm{H}_{49}[429,1238]$, $\mathrm{H}_{56}[947,1238]$\\
65&$[122,1351]$&$\mathrm{S}_{2}[120,532]$, $\mathrm{H}_{50}[435,1294]$, $\mathrm{H}_{57}[1040,1294]$\\
66&$[124,1409]$&$\mathrm{S}_{2}[122,563]$, $\mathrm{H}_{52}[514,1351]$, $\mathrm{H}_{58}[1054,1351]$\\
67&$[126,1468]$&$\mathrm{S}_{2}[124,565]$, $\mathrm{H}_{51}[448,1409]$, $\mathrm{H}_{59}[1068,1409]$\\
68&$[128,1528]$&$\mathrm{S}_{2}[126,597]$, $\mathrm{H}_{54}[528,1468]$, $\mathrm{H}_{60}[1082,1468]$\\
69&$[130,1589]$&$\mathrm{S}_{2}[128,599]$, $\mathrm{H}_{55}[535,1528]$, $\mathrm{H}_{61}[1181,1528]$\\
70&$[132,1651]$&$\mathrm{S}_{2}[130,632]$, $\mathrm{H}_{56}[542,1589]$, $\mathrm{H}_{62}[1196,1589]$\\
71&$[134,1714]$&$\mathrm{S}_{2}[132,634]$, $\mathrm{H}_{56}[550,1651]$, $\mathrm{H}_{63}[1211,1651]$\\
72&$[136,1778]$&$\mathrm{S}_{2}[134,668]$, $\mathrm{H}_{57}[557,1714]$, $\mathrm{H}_{64}[1226,1714]$\\
73&$[138,1843]$&$\mathrm{S}_{2}[136,670]$, $\mathrm{H}_{57}[565,1778]$, $\mathrm{H}_{65}[1331,1778]$\\
74&$[140,1910]$&$\mathrm{S}_{2}[138,705]$, $\mathrm{H}_{58}[572,1843]$, $\mathrm{H}_{67}[1816,1843]$\\
75&$[142,1977]$&$\mathrm{S}_{2}[140,707]$, $\mathrm{H}_{59}[579,1910]$, $\mathrm{H}_{67}[1363,1910]$\\
76&$[144,2046]$&$\mathrm{S}_{2}[142,743]$, $\mathrm{H}_{62}[669,1977]$, $\mathrm{H}_{69}[1957,1977]$\\
77&$[146,2116]$&$\mathrm{S}_{2}[144,745]$, $\mathrm{H}_{63}[677,2046]$, $\mathrm{H}_{70}[1979,2046]$\\
78&$[148,2188]$&$\mathrm{S}_{2}[146,782]$, $\mathrm{H}_{63}[686,2116]$, $\mathrm{H}_{72}[2102,2116]$\\
79&$[150,2260]$&$\mathrm{S}_{2}[148,784]$, $\mathrm{H}_{64}[694,2188]$, $\mathrm{H}_{72}[2126,2188]$\\
80&$[152,2334]$&$\mathrm{S}_{2}[150,822]$, $\mathrm{H}_{64}[614,2260]$, $\mathrm{H}_{73}[2149,2260]$, $\mathrm{H}_{74}[2253,2260]$\\
81&$[154,2409]$&$\mathrm{S}_{2}[152,824]$, $\mathrm{H}_{65}[711,2334]$, $\mathrm{H}_{75}[2277,2334]$\\
82&$[156,2485]$&$\mathrm{S}_{2}[154,863]$, $\mathrm{H}_{66}[719,2409]$, $\mathrm{H}_{75}[2302,2409]$, $\mathrm{H}_{76}[2409,2409]$\\
83&$[158,2562]$&$\mathrm{S}_{2}[156,865]$, $\mathrm{H}_{67}[727,2485]$, $\mathrm{H}_{77}[2434,2485]$\\
84&$[160,2640]$&$\mathrm{S}_{2}[158,905]$, $\mathrm{H}_{70}[827,2562]$, $\mathrm{H}_{78}[2459,2562]$\\
\end{longtable}}

\section{Supplementary material}\label{app:data}
The finite computations used in this paper are collected in the supplementary archive
\file{verification.zip}. It contains $26$ files: a \file{README.md}, a file
\file{manifest.json} with the SHA-256 hash of every other file, and $24$ programs and
data files. The programs use Python~3.12, NumPy and SymPy~1.14.0. A program
\file{name.py} at the top level is run from the root of the unpacked archive as
\texttt{python name.py}, and a program in the directory \file{boundary_moment_v1} as
\texttt{python boundary\_moment\_v1/name.py}. Each program writes its JSON result to the
location of the distributed result file, so that a copy of the distributed file should be
kept for comparison.

The computations are of three kinds.

\emph{(i) Fixed symbolic identities and nonvanishing certificates.} These are used as input
to arguments valid for infinitely many degrees.
\begin{itemize}
\item \file{check_cyclic_uniform_origin_rank_v1.py} and
  \file{check_cyclic_bulk_moment_symbolic_v1.py}, with their JSON results: the
  identities of rational functions used in Theorem~\ref{thm:special},
  Proposition~\ref{prop:origin} and Proposition~\ref{prop:bulk}.
\item The directory \file{boundary_moment_v1}: five programs
  \file{certify_boundary_*.py} with their five JSON certificates and two shared modules.
  The certificates contain the data of Appendix~\ref{app:boundary}: the $36$ and $24$
  leading coefficients of Tables~\ref{tab:quadall} and~\ref{tab:twoall}, the coefficients
  for $\chi=-1$ and $\chi=0$, the seven cubic obstruction functions $\mathcal F_\chi$ as
  rational functions of $T$, and the exact rational Horner bounds for their numerators and
  denominators. The program for the cubic line also checks that $\det(\mathcal A,\mathcal C)$
  is nonzero on each interval, but its JSON result does not store the corresponding exact
  bounds; these are reproduced from \eqref{eq:DAC} as described in
  Appendix~\ref{app:boundary}.
\end{itemize}

\emph{(ii) Certificates for finitely many explicit degrees.}
\begin{itemize}
\item \file{certify_cyclic_finite_composition_v1.py} with its JSON result: the
  computation of Appendix~\ref{app:finite} for $29\le n\le88$ and all $3510$ characters,
  with the chosen primes, the elements of order $n$ and the checks of their orders, the
  denominator conditions, the nonzero minors, and the passage to characteristic zero.
  The SHA-256 hash of this certificate file is
  \begin{center}\small\file{8d1b83ef61e872208c3034a5fb07980bf202fe0d86eb4bf91db07bdb2678b501}.\end{center}
\item \file{verify_direct_rational_seed_v1.py}, the data file
  \file{FINITE_SEED_RATIONAL_DATA_v1.json} and the record
  \file{verify_direct_rational_seed_v1_execution.json}: the rational configuration
  of Appendix~\ref{app:seed} for $(d,s)=(11,14)$ and the two nonzero minors of size $896$.
\end{itemize}

\emph{(iii) Bases of inductions and finite transitions.}
\file{check_inductive_ranges.py} with its JSON result, and
\file{data/compile_one_seed_bridge_v2.json}: the $188$ transitions of
Appendix~\ref{app:propagation} up to degree $84$, and a check of the integer inequalities
used in Theorem~\ref{thm:broad}. Theorem~\ref{thm:broad} is proved in the text, and the
check repeats its inequalities. The bands of Proposition~\ref{prop:bands}, with upper
endpoints $2d-24$ and $3d-52$, are proved only by the induction in the text. The
small-band part of \file{check_inductive_ranges.py} concerns the weaker endpoints $2d-26$
and $3d-56$, and it does not verify the bands stated here.

Each certificate records the data needed to check it independently of the search that
produced it. A rank certificate (Appendices~\ref{app:finite} and~\ref{app:seed}) records
the selected nonzero minors, with their rows and columns. A symbolic or sign certificate
records the exact rational functions together with the coefficients or exact interval
bounds that establish their identities or their nonvanishing.

\section*{Use of AI-assisted tools}
AI-assisted tools were used in the research and preparation of this manuscript.

\begin{thebibliography}{9}

\bibitem{Bozlee}
S.~Bozlee, \emph{Close connectedness of the moduli stack of reduced curves}, preprint
(2026), arXiv:2604.03930v2.

\bibitem{deBoorRon}
C.~de~Boor and A.~Ron, \emph{The least solution for the polynomial interpolation problem},
Math. Z. \textbf{210} (1992), no.~3, 347--378.

\bibitem{DJNT}
T.~Douvropoulos, J.~Jelisiejew, B.~I.~U.~N\o dland, and Z.~Teitler, \emph{The Hilbert
scheme of 11 points in $\mathbb A^3$ is irreducible}, in: Combinatorial Algebraic Geometry,
Fields Inst. Commun. \textbf{80}, Springer, New York, 2017, 321--352,
doi:10.1007/978-1-4939-7486-3\_15.

\bibitem{JelElem}
J.~Jelisiejew, \emph{Elementary components of Hilbert schemes of points}, J. Lond. Math.
Soc. (2) \textbf{100} (2019), no.~1, 249--272; theorem numbers as in arXiv:1710.06124v5.

\bibitem{JelOpen}
J.~Jelisiejew, \emph{Open problems in deformations of Artinian algebras, Hilbert schemes
and around}, preprint (2023, revised 2026), arXiv:2307.08777v2.

\bibitem{Kustin}
A.~R.~Kustin, \emph{Canonical complexes associated to a matrix}, preprint (2015),
arXiv:1509.04367v1.

\bibitem{Pinkham}
H.~C.~Pinkham, \emph{Deformations of algebraic varieties with $G_m$ action}, Ast\'erisque
\textbf{20}, Soci\'et\'e Math\'ematique de France, Paris, 1974.

\bibitem{Sauer}
T.~Sauer, \emph{The number of equations defining points in general position}, Pacific J.
Math. \textbf{120} (1985), no.~1, 199--213.

\bibitem{Stevens}
J.~Stevens, \emph{Non-smoothable curve singularities}, J. Singul. \textbf{29} (2026),
100--122, doi:10.5427/jsing.2026.29d; arXiv:2504.00854v2.

\end{thebibliography}

\end{document}
